\documentclass{article}

\usepackage[backend=biber,style=gb7714-2015,sorting=gb7714-2015,gbalign=left,sorting=gb7714-2015]{biblatex}
\usepackage{amsmath, amsthm, amssymb, amsfonts}
\usepackage{thmtools}
\usepackage{setspace}
\usepackage{graphicx}
\usepackage{geometry}
\usepackage{float}
\usepackage{hyperref}
\usepackage[utf8]{inputenc}
\usepackage[english]{babel}
\usepackage{framed}
\usepackage[dvipsnames]{xcolor}
\usepackage{tcolorbox}
\usepackage{tikz-cd}
\usepackage{tikz} 
\usepackage{array}
\usetikzlibrary{decorations.markings, arrows.meta, calc}


\tcbuselibrary{breakable}

\addbibresource{references.bib}

\colorlet{LightGray}{White!90!Periwinkle}
\colorlet{LightOrange}{Orange!15}
\colorlet{LightGreen}{Green!15}

\newcommand{\HRule}[1]{\rule{\linewidth}{#1}}

\declaretheoremstyle[name=Theorem,]{thmsty}
\declaretheorem[style=thmsty,numberwithin=section]{theorem}
\tcolorboxenvironment{theorem}{colback=White}

\declaretheoremstyle[name=Definition,]{prosty}
\declaretheorem[style=prosty,numberlike=theorem]{definition}
\tcolorboxenvironment{definition}{colback=White}

\declaretheoremstyle[name=Lemma,]{prcpsty}
\declaretheorem[style=prcpsty,numberlike=theorem]{lemma}
\tcolorboxenvironment{lemma}{colback=White}

\newtheorem{proposition}[theorem]{Proposition}
\newtheorem{corollary}[theorem]{Corollary}
\newtheorem{remark}[theorem]{Remark}
\newtheorem{example}[theorem]{Example}
\newtheorem{exercise}[theorem]{Exercise}


\def\E{\mathop{\mathcal E}}
\def\D{\mathop{\mathcal D}}
\def\lhdeq{\mathop{\trianglelefteq}}
\def\<{\mathop{\langle}}
\def\>{\mathop{\rangle}}
\def\l{\mathop{\ell}}
\newcommand{\g}{\mathfrak{g}}  
\newcommand{\h}{\mathfrak{h}}  
\newcommand{\n}{\mathfrak{n}}  
\newcommand{\U}{\mathcal{U}}  
\newcommand{\C}{\mathbb{C}} 
\newcommand{\borel}{\mathfrak{b}}
\newcommand{\Ind}{\mathrm{Ind}} 
\newcommand{\Hom}{\mathrm{Hom}}  
\newcommand{\Ext}{\mathrm{Ext}}

\setstretch{1.2}
\geometry{
    textheight=9in,
    textwidth=5.5in,
    top=1in,
    headheight=12pt,
    headsep=25pt,
    footskip=30pt
}

% ------------------------------------------------------------------------------

\begin{document}

% ------------------------------------------------------------------------------
% Cover Page and ToC
% ------------------------------------------------------------------------------

\title{ \normalsize \textsc{}
		\\ [2.0cm]
		\HRule{1.5pt} \\
		\LARGE \textbf{\uppercase{Topics in Algebra}
		\HRule{2.0pt} \\ [0.6cm] \LARGE{By Prof. Iryna Kashuba} \vspace*{10\baselineskip}}
		}
\date{}
\author{\textbf{Leo}}

\maketitle
\newpage
 
\tableofcontents
\newpage

\section{Lecture 1}

In our course for the base field $ k $ we assume $ \operatorname{Char}k \neq 2 $. 

Jordan algebras arose from the search for an “exceptional” setting for quantum mechanics. This question is related to sixth problem of Hilbert, which be presented in 1900, where he suggested axiomatic those branches of physics in which mathmatics is prevalent. For both classical and quantum mechanics physical observables are fundemental objects. In classical mechanics they form associative commutative algebra. In case of quantum mechanics observables are Hermitian or self-adjoint operators on Hilbert space (or Hermitian matrix).

Consider field $\mathbb{R}$, then the set of observables forms a vector space, which is not closed with respect to product (or composition of operators), since for $(AB)^* = B^*A^* = BA$ not necessarily $= AB$.

Jordan observed that $(A^*)^2 = A^2$ and since $(A+B)^2 = A^2 + (AB+BA) + B^2$, we get that $AB+BA = (A+B)^2 - A^2 - B^2$ is Hermitian. Thus, if we consider a new operation $A \circ B = AB+BA$ on a vector space of observables, we obtain an algebra. This algebra is not associative, it satisfies the following relations:
\begin{align*}
    A \circ B &= B \circ A \\
    ((A \circ A) \circ B) \circ A &= (A \circ A) \circ (B \circ A)
\end{align*}

Observe also that $A \circ A = 2A^2$. To fix it
\[
A \cdot B = \frac{1}{2}(AB+BA)
\]
Then we obtain an algebra isomorphic to origin one, with the difference that powers of $A$ in a new algebra coincide with the powers in the original algebra.

In 1933 the physicist Pascual Jordan proposed a program to discover a new algebraic setting for quantum mechanics, which would be freed from dependence on an invisible all-determining meta physical matrix structure.

In their classification (together with Neumann and Wigner) they get only one algebra of new type.

\begin{itemize}
\item commutator: $ [a,b] =a b-b a$
\item associator: $ (a,b,c)=(ab) c-a (b c) $.  
\end{itemize}   

\begin{definition}
An algebra $J$ over $F$ is a Jordan algebra if $[x,y]=0$ and $(x^2,y,x)=0$ for all $x,y \in J$.
\end{definition}

\begin{example}[$ A^{(+)} $ algebras] 
Let $A$ be an algebra. Define on vector space $ A $ a new multiplication by
\[
a \cdot b = \frac{1}{2}(ab+ba)
\]
If $A$ is an associative algebra, then $A^{(+)} = (A, +, \cdot)$ is a Jordan algebra.
\end{example}

Observation: Consider $ a \circ b=ab+ba $, then $ \varphi:  (A, +, \circ) \to (A, +, \cdot) $ s.t. $ \phi(a)=2a $ is an algebraic isomorphism.

check: $ \phi(a \circ b) = \phi((ab+ba)) = 2(ab+ba) = 2ab+2ba = \phi(a) \cdot \phi(b) $.

\begin{example}[Hermitian algebras] 
Suppose $ A $ is an associative algebra with involution $ *$. $ H(A,*) =\{a \in A | a^*=a\}$. Then $ H(A,*) $ is a Jordan subalgebra of $ A^{(+)} $. 
\end{example}

\begin{definition} 
We call Jordan algebra $ J $ special if it can be embedded/ is a subalgebra in $ A^{(+)} $ for some associative algebra $ A $. Otherwise it is called exceptional.
\end{definition}

Obviously, Hermitian algebras are special.

\begin{exercise}
Let $ A $ be an associative algebra, $ A^{(+)} $ is a Jordan algebra. Let $ B=A\oplus A^{op}$ an associative algebra with involution $ * : (a,b)^*=(b,a)$. Prove that $ A^{(+)} \simeq H(B,*) $. 
\end{exercise}

\begin{proof} 
Since $ H(B, *) $ is $ * $-invariant, $ \forall x \in B $ such that $ x=(a,b) $ where $ a \in A $ and $ b \in A^{op} $, $ x^*=x $ shows:
\[
(b,a)=(a,b)
\] 
Thus $ a=b $. i.e.
\[
H(B,*)= \{ (a,a) | a \in A \}
\] 
We define a map $ \phi: A^{(+)} \rightarrow H(B,*) $ s.t. $ a \mapsto (a,a) $. Then obviously $ \phi $ is a linear isomorphism of vector space. We only need to check that $ \phi $ preserves the Jordan product.

Now consider for all $ a,b \in A^+ $ we have:
\[
\phi(a \cdot b)=\phi\left(\frac{1}{2}(ab+ba)\right)=\left(\frac{1}{2}(ab+ba),\frac{1}{2}(ab+ba)\right)
\]
And
\[
\phi(a) \cdot \phi(b)=(a,a)\cdot (b,b)
\]
Remark that the Jordan product on $ H(B,*) $ is inherited from $ B^+ $, that is $ (a_1, b_1) \cdot (a_2,b_2)=\frac{1}{2}((a_1, b_1)(a_2,b_2)+(a_2,b_2)(a_1, b_1)) $. Thus
\[
(a,a)\cdot (b,b)=\frac{1}{2}((a,a)(b,b)+(b,b)(a,a))=\frac{1}{2}((ab,ba)+(ba,ab))=\left(\frac{1}{2}(ab+ba),\frac{1}{2}(ab+ba)\right)
\]   
\end{proof}

\begin{example}[Clifford type Jordan algebra or spin factor] 
 Let $V$ be an $ F $-vector space endowed with a symmetric bilinear form $f: V \times V \to F$. Then the vector space $J(V,f) = V \oplus F \cdot 1$ equipped with the following product
\[
  ( \alpha+ u) \cdot ( \beta+ v) = (\alpha \beta + f(u,v) )\cdot 1 + \alpha v +\beta u \quad \forall  u,v \in V, \alpha, \beta \in F.
\]
is a Jordan Algebra
\end{example}

\begin{example}[Albert algebra] 
Let $\mathbb{O}$ be the algebra of octonions. Algebra with involution $*$ and basis $\{1, e_1, \dots, e_7\}$ such that for $i \neq j$
\begin{align*}
    1^* = 1,  e_i^* = -e_i,  e_i^2 = -1, e_i e_{i+1} = e_{i+3} \pmod 7 ,  e_i  e_j = -e_j e_i
\end{align*}
and if  $ e_i e_j = e_k $ then $e_{\sigma(i)} e_{\sigma(j)} = \operatorname{sgn}(\sigma) e_{\sigma(k)}$ for any $\sigma \in S_3$.

\begin{tikzpicture}[
    scale=3,
    point/.style={circle, fill=black, inner sep=2pt, outer sep=0pt},
    label_style/.style={font=\large\bfseries},
    % Define arrow style to place an arrow in the middle of the path
    arrow_style/.style={
        decoration={
            markings,
            mark=at position 0.6 with {\arrow[scale=1.5]{>}}
        },
        postaction={decorate}
    }
]

% Define vertices of the triangle (Circumradius 2)
\coordinate (e1) at (90:2);    % Top vertex
\coordinate (e2) at (330:2);   % Bottom Right vertex
\coordinate (e6) at (210:2);   % Bottom Left vertex

% Define midpoints (which are also basis elements)
% The inner circle passing through these has radius 1
\coordinate (e4) at ($(e1)!0.5!(e2)$); % Midpoint of 1-2
\coordinate (e5) at ($(e1)!0.5!(e6)$); % Midpoint of 1-6
\coordinate (e7) at ($(e6)!0.5!(e2)$); % Midpoint of 6-2

% Center
\coordinate (e3) at (0,0);

% Draw the inner circle passing through midpoints e4, e5, e7
% Rule: e4 e5 = e7 => Cycle (4, 5, 7) => Counter-Clockwise
\draw[arrow_style, thick] (0,0) circle (1);

% Draw the edges of the triangle with arrows

% Right Edge: Rule e1 e2 = e4 => Cycle (1, 2, 4) => 1 -> 2
\draw[arrow_style, thick] (e1) -- (e2);

% Bottom Edge: Rule e6 e7 = e2 => Cycle (6, 7, 2) => 6 -> 2
\draw[arrow_style, thick] (e6) -- (e2);

% Left Edge: Rule e5 e6 = e1 => Cycle (5, 6, 1) => 1 -> 6
% (Note: 1->5->6 is consistent with e1*e5=e6)
\draw[arrow_style, thick] (e1) -- (e6);

% Draw the medians with arrows

% Vertical Median: Rule e7 e1 = e3 => Cycle (7, 1, 3) => 1 -> 7
% (Note: e1*e3=e7, so 1->3->7)
\draw[arrow_style, thick] (e1) -- (e7);

% Median to Left: Rule e2 e3 = e5 => Cycle (2, 3, 5) => 2 -> 5
\draw[arrow_style, thick] (e2) -- (e5);

% Median to Right: Rule e3 e4 = e6 => Cycle (3, 4, 6) => 6 -> 4
% (Note: e6*e3=e4, so 6->3->4)
\draw[arrow_style, thick] (e6) -- (e4);

% Draw points
\foreach \p in {e1,e2,e3,e4,e5,e6,e7}
    \node[point] at (\p) {};

% Labels
\node[label_style, above] at (e1) {$e_1$};
\node[label_style, below right] at (e2) {$e_2$};
\node[label_style, below left] at (e6) {$e_6$};
\node[label_style, right] at (e4) {$e_4$};
\node[label_style, left] at (e5) {$e_5$};
\node[label_style, below] at (e7) {$e_7$};
\node[label_style, above right, shift={(0.05,0.05)}] at (e3) {$e_3$};

\end{tikzpicture}

Let $\mathbb{A} = H_3(\mathbb{O}, *)$ be an algebra of matrices of order 3 over $\mathbb{O}$ which are hermitian
\[
\mathbb{A} = \left\{ \begin{pmatrix} \alpha & a & b \\ a^* & \beta & c \\ b^* & c^* & \gamma \end{pmatrix} : \alpha, \beta, \gamma \in F, \ a,b,c \in \mathbb{O} \right\}
\]
with operation
\[
x \cdot y = \frac{1}{2}(xy+yx),
\]
$\mathbb{A}$ is a Jordan algebra of dim $27$ (over $F$). This is called Albert algebra since in 1934 Albert proved that $ H_3(\mathbb{O}) $ is indeed exceptional. 

\end{example}

Simple Jordan algebras discovered by Jordan, von Neumann and Wigner are
\begin{enumerate}
    \item $\mathbb{R}_n^{(+)}$
    \item $H(\mathcal{D}_n,*)$ for $\mathcal{D} = \mathbb{R}, \mathbb{C}, \mathbb{H}$
    \item $J(V,f)$
    \item $H_3(\mathbb{O},*)$ 
\end{enumerate}

The algebras $A^+$ and $H(A,*)$ are obviously special. (Remark: the notion here means, given an algebra $ A $, $ A_n $ denotes the algebra of $ n \times n $ matrices with entries in $ A $).

\begin{exercise}
Prove that $ J(V, f) $ is a special Jordan algebra.
\end{exercise}

\begin{definition} 
Let $ V $ be a finite-dimensional vector space over a field $ F $, equipped with a quadratic form
\[
Q : V \to F.
\]
The Clifford algebra $ \mathrm{Cl}(V, Q) $ (also written as $ \mathrm{Cl}_F(V, Q) $) is the associative unital algebra over $ F $ generated by the vector space $ V $, subject to the relation:
\[
v^2 = Q(v) \cdot 1 , \forall  v \in V.
\]
Equivalently, for all $ u, v \in V $,
\[
uv + vu = 2 f(u,v) \cdot 1,
\]
where $ f(u,v) $ is the bilinear form associated with $ Q $:
\[
f(u,v) = \frac{1}{2}(Q(u+v) - Q(u) - Q(v)).
\]
\end{definition}

\begin{proof} 
We want to show that $ J(V,B) $ can be embedded in $ \mathrm{Cl}^+(V, Q) $ where $ Q(v)=f(v,v) $.     

Define a linear map: $ \phi: J(V,f) \rightarrow \mathrm{Cl}^+(V, Q) $ such that:
\[
\alpha \cdot 1_J + v \mapsto \alpha \cdot 1_C + v
\] 
This is obviously injective. We only need to show this map preserve Jordan product. For $ x,y \in J(V,f) $, $ x=\alpha \cdot 1_J + v $ and $ y=\beta \cdot 1_J + w $, we have:
\[
x \circ y = (\alpha \beta + f(v,w)) \cdot 1_J + \alpha w + \beta v 
\]
Hence
\[
\phi(x \circ y) = (\alpha \beta + f(v,w)) \cdot 1_C + \alpha w + \beta v 
\]
On the other hand
\[
\phi(x) \circ \phi(y) = (\alpha \cdot 1_C + v) \circ (\beta \cdot 1_C + w )
\]
This $ \circ $ is inherent from the product in associative algebra $ \mathrm{Cl}(V,Q) $, where $ vu+uv=2f(u,v) \cdot 1_C $. Thus
\[
(\alpha \cdot 1_C + v) \circ (\beta \cdot 1_C + w )=\frac{1}{2}((\alpha \cdot 1_C + v)(\beta \cdot 1_C + w )+(\beta \cdot 1_C + w )(\alpha \cdot 1_C + v))
\]    
Where
\[
RHS=\frac{1}{2}2f(\alpha \cdot 1_C + v,\beta \cdot 1_C + w ) \cdot 1_C=(\alpha \beta f(1_C,1_C)+\alpha f(1_C,w) +\beta f(v,1_C)+f(v,w)) \cdot 1_C
\]
Since $ 1_C=1_C^2=f(1_C,1_C) \cdot 1_C $, $ f(1_C,1_C)=1 $. Similarly $ f(1_C,w)=w $ and $f(v,1_C)=v $. Thus
\[
RHS=(\alpha \beta+\alpha w + \beta v +f(v,w)) \cdot 1_C=(\alpha \beta + f(v,w)) \cdot 1_C + \alpha w + \beta v
\]
Therefore, we have
\[
\phi(x \circ y)=\phi(x) \circ \phi(y)
\]  
\end{proof}

Jordan hoped that by studying finite-dimensional algebras he would find families of simple exceptional algebras $E_n$ parametrized by natural numbers $n$. So That letting $n \to \infty$ to infinity would provide a suitable infinite dimensional exceptional home for quantum mechanics.

These results were deeply disappointing to physicists. There was only one exceptional algebra in the list for $n=3$. For $n>3$ the algebra $H_n(\mathbb{O})$ is not a Jordan algebra at all. For $n=2$ it is a clifford type algebra isomorphic to a $J(V,f)$ and for $n=1$ it is just $\mathbb{R}^+$.

This exceptional algebra $H_3(\mathbb{O})$ was too tiny to provide a home for quantum mechanics and too isolated to give a clue as to the possible existence of infinite-dimensional exceptional algebras. Half a century later the brilliant young Novosibirsk mathematician Efim Zelmanov quashed all remaining hopes for such an exceptional system. In 1979 he showed that even in infinite dimensions there are no simple exceptional Jordan algebras other than Albert algebras.

For mathematicians it turned out to have important applications in Lie theory, Differential Geometry, Projective Geometry.

Observation: Algebra is ring + vector space over a field $ F $.
\[ \alpha(ab) = (\alpha a)b = a(\alpha b) \quad \alpha \in F, a,b \in A. \]
We can also consider algebra over a ring (every ring $A$ is an algebra over $\mathbb{Z}$).

\begin{proposition}
Let $ A $ be simple associative algebra, then $ A^{(+)} $ is a simple Jordan algebra.
\end{proposition}

\begin{definition} 
An algebra $ A $ is called simpe if it does not contain non-trivial ideals. 
\end{definition}

\begin{proof} 
Let $ I $ be a non-zero ideal of $ A^{(+)} $. We use the notation $ x \circ y = xy + yx $. Note that $ I $ is closed under $\circ$.
For any $ x,y \in I $ and $ a \in A $, we have the identity:
\[
[a, x \circ y] = x \circ [a, y] - [a, x] \circ y
\]
Since $ x, y \in I $ and $ I $ is an ideal of $ A^{(+)} $, both $ x \circ [a, y] $ and $ [a, x] \circ y $ are in $ I $. Thus, the commutator $ [a, x \circ y] $ is in $ I $.

Also, since $ x \circ y \in I $, we have $ a \circ (x \circ y) \in I $. Adding these two results:
\[
2a(x \circ y) = [a, x \circ y] + a \circ (x \circ y) \in I.
\]
Now for any $ b \in A $, consider the Jordan product of $ 2a(x \circ y) \in I $ with $ b $:
\[
(2a(x \circ y)) \circ b = 2a(x \circ y)b + 2b a (x \circ y) \in I.
\]
From the previous step (replacing $ a $ with $ ba $), we know $ 2ba(x \circ y) \in I $.
Therefore, $ 2a(x \circ y)b \in I $ for all $ a, b \in A $.
Let $ z = x \circ y $. We have shown that $ 2 A z A \subseteq I $. Since $ \operatorname{char} k \neq 2 $, $ A z A \subseteq I $.

Since $ A $ is simple, for any $ z \in A $, the ideal $ A z A $ is either $ 0 $ (if $ z=0 $) or $ A $ (if $ z \neq 0 $).

Case 1: There exist $ x, y \in I $ such that $ z = x \circ y \neq 0 $.
In this case, $ A z A = A $, so $ A \subseteq I $, which implies $ I = A^{(+)} $.

Case 2: For all $ x, y \in I $, $ x \circ y = 0 $.
This implies $ 2x^2 = x \circ x = 0 $, so $ x^2 = 0 $ for all $ x \in I $.
For any $ a \in A $ and $ x \in I $, we have $ x \circ a \in I $. Thus, $ x \circ (x \circ a) = 0 $.
Expanding this:
\[
x(xa+ax) + (xa+ax)x = x^2 a + xax + xax + ax^2 = 2xax = 0.
\]
So $ xax = 0 $ for all $ a \in A $. This means $ xAx = 0 $.
Consider the ideal generated by $ x $. Since $ A $ is simple, if $ x \neq 0 $, then $ AxA = A $ (or at least $ (AxA)^2 \neq 0 $).
However, $ (AxA)^2 = Ax(AxA)x A = A(xAx)AxA = 0 $. Since $ A $ is simple and not nilpotent (simple algebras are not nilpotent), we must have $ AxA = 0 $, which implies $ x = 0 $.
Thus, in this case, $ I = \{0\} $.

Therefore, the only ideals of $ A^{(+)} $ are $ 0 $ and $ A^{(+)} $ itself.
\end{proof}

\begin{definition} 
We say $ A $ is $ * $-simple if $ A \neq 0 $ and $ A$ does not contain non-trivial ideals $ I $ s.t. $ I^*=I $.
\end{definition}

\begin{theorem}[Herstein] 
Let $ A $ be an associative algebra with involution $ * $. If $ (A,*) $ is $ * $-simple then $ H(A,*) $ is a simple Jordan algebra. 
\end{theorem}

\begin{exercise}
Let $ J(V,f) $ be a Clifford type Jordan algebra. Show that if $ \dim V > 1 $ and $ f $ is non-degenerated, then $ J(V,f) $ is simple.
\end{exercise}

\begin{proof}
Let $J=J(V,f)=\Bbb F1\oplus V$ be the Clifford type Jordan algebra with product
\[
(a1+u)\circ(b1+v)=(ab+f(u,v))1+(av+bu),
\qquad a,b\in\Bbb F,\ u,v\in V,
\]
where $f$ is a symmetric bilinear form on $V$. Assume $\dim V>1$ and $f$ is
non-degenerate. We prove that $J$ is simple.

Let $0\neq I\lhd J$ be a (Jordan) ideal. Choose $0\neq x=a1+u\in I$ with
$a\in\Bbb F$ and $u\in V$.

\noindent
\textbf{Case 1: $u=0$.}
Then $x=a1\in I$ and $a\neq 0$, hence $1\in I$, so $I=J$.

\noindent
\textbf{Case 2: $u\neq 0$.}
If $a=0$, then $x=u\in I$. Since $f$ is non-degenerate, there exists $v\in V$
such that $f(u,v)\neq 0$. Then
\[
u\circ v=f(u,v)1\in I,
\]
so $1\in I$ and hence $I=J$.

It remains to treat the subcase $a\neq 0$. Again pick $v\in V$ with
$f(u,v)\neq 0$ and set
\[
y:=x\circ v=(a1+u)\circ v=f(u,v)1+av\in I.
\]
Since $I$ is an ideal, $y\circ v\in I$. Using $v\circ v=f(v,v)1$, we compute
\[
y\circ v=(f(u,v)1+av)\circ v=f(u,v)v+a(v\circ v)=f(u,v)v+a f(v,v)1\in I.
\]
Also, because $a\neq 0$, we have $\frac{f(u,v)}{a}\,y\in I$, i.e.
\[
\frac{f(u,v)}{a}\,y=\frac{f(u,v)^2}{a}\,1+f(u,v)v\in I.
\]
Subtracting these two elements of $I$ gives a pure scalar:
\[
\Bigl(y\circ v-\frac{f(u,v)}{a}\,y\Bigr)
=\Bigl(a f(v,v)-\frac{f(u,v)^2}{a}\Bigr)1\in I. \tag{$\ast$}
\]
Thus, it suffices to choose $v$ so that the coefficient
$\;a f(v,v)-\frac{f(u,v)^2}{a}\;$ is nonzero.

To achieve this, choose $w\in V$ with $f(u,w)=0$ and $f(w,w)\neq 0$.
(One can do this because $\dim V>1$ and $f$ is non-degenerate: the orthogonal
hyperplane $u^\perp=\{w\in V:\ f(u,w)=0\}$ has dimension $\dim V-1\ge 1$, and
$f$ cannot vanish identically on $u^\perp$ under non-degeneracy, hence there
exists $w\in u^\perp$ with $f(w,w)\neq 0$.)

For $t\in\Bbb F$, set $v_t:=v+t w$. Then $f(u,v_t)=f(u,v)+t f(u,w)=f(u,v)\neq 0$
is independent of $t$, while
\[
f(v_t,v_t)=f(v+tw,v+tw)=f(v,v)+2t f(v,w)+t^2 f(w,w)
\]
is a quadratic polynomial in $t$ whose leading coefficient is $f(w,w)\neq 0$.
Hence the coefficient in $(\ast)$ becomes
\[
\lambda(t):=a f(v_t,v_t)-\frac{f(u,v_t)^2}{a}
=a\bigl(f(v,v)+2t f(v,w)+t^2 f(w,w)\bigr)-\frac{f(u,v)^2}{a},
\]
a nonzero quadratic polynomial (its $t^2$-coefficient is $a f(w,w)\neq 0$).
Therefore $\lambda(t)\neq 0$ for some $t\in\Bbb F$. For this choice of $t$,
equation $(\ast)$ yields $\lambda(t)\,1\in I$, so $1\in I$, and thus $I=J$.

In all cases, every nonzero ideal $I$ equals $J$. Hence $J(V,f)$ is simple.
\end{proof}

\section{Lecture 2}

Here we describe linearizations of Jordan identity. Char of field or ring $\ne 2$, thus exist at least 3 different elements in it. 

This identity
\begin{equation}
a^2 \circ (b \circ a)=(a^2 \circ b) \circ a \tag{1}
\end{equation}
is not linear. Consider $ (a,b,c)=(ab)c - a(bc)$ (from here we ignore $ \circ $ for short-hand) and $ (a^2,b,a)=0 $ the Jordan identity. Replace $ a $ by linear transformation $ a \mapsto a+\lambda c $ we have:
\[
((a+\lambda c)^2 ,b , (a+\lambda c))=(a^2,b,a) +\lambda f(a,b,c)+ \lambda^2 g(a,b,c)+\lambda^3(c^2,b,c)=0
\]    
Where 
\begin{equation}
 f(a,b,c)=2(ac,b,a)+(a^2,b,c) \tag{2}
\end{equation}
and 
\begin{equation*}
g(a,b,c)=2(ac,b,c)+(c^2,b,a)
\end{equation*}
and 
\[
0= ((a+\lambda c)^2 ,b , (a+\lambda c)) =\lambda f(a,b,c)+ \lambda^2 g(a,b,c)
\]
if $ \lambda \neq 0 $. i.e., $f(a,b,c)+ \lambda g(a,b,c)=0 $. Put $ \lambda =1 $ and then put $ \lambda = -1 $ we have 
\[ 
f(a,b,c)= g(a,b,c)=0
\]

Now take $ f(a,b,c)=2(ac,b,a)+(a^2,b,c) $ (if you take $ g $ you will get the same result, since f and g are just switch $ a $ and $ c $) and do transformation $ a \mapsto a+d $.   

Finally we have:
\[
(dc,b,a)+(ac,b,d)+(ad,b,c)=0 \tag{3}
\] 

It is useful to write in operator form (In our class we use right module to write operators). Define
\[
L_x: J\rightarrow J \ \text{s.t.} \ aL_x=xa
\]
Similarly
\[
R_x: J\rightarrow J \ \text{s.t.} \ aR_x=ax
\]
Consider them as elements in $ \operatorname{End}(J) $ (endumorphism as vector space). 

$ R(J) $ is associative algebra generated by $ R_x, x \in J $ and $ L(J) $ is associative algebra generated by $ L_x, x \in J $. For Jordan algebra they coincide since Jordan product is commutative, $ L_x =R_x $. 

$ \operatorname{Mult}(J) $ denotes multiplication algebra. It is a associative algebra generated by $ L_x, R_x, x \in J $. Given $ J $ a Jordan algebra, $ \operatorname{Mult}(J)=R(J)=L(J) $. by (3) we have
\[
((dc)b)a-(dc)(ba)+((ad)b)c-(ad)(bc)+((ac)b)d-(ac)(bd)=0
\]
By the linearization of the Jordan identity. Write them into the form $ R(J) $ for b and d respectively:
\[
b: \ R_{dc}R_a-R_aR_{dc}+R_{ad}R_c-R_cR_{ad}+R_{ac}R_d-R_dR_{ac}=0
\]
i.e.
\[
[R_{dc},R_a]+ [R_{ad},R_c]+[R_{ac},R_d]=0
\]
Similarly write it into action on $ d $
\begin{equation}
d:  R_cR_bR_a+R_aR_bR_c+R_{(ac)b}=R_cR_{ba}+R_aR_{bc}+R_bR_{ac} \tag{**}
\end{equation}

Thus the three identity:

(1) is equal to $ [R_{a^2},R_a] =0$ (I);

(2) is equal to $ [R_{a^2},R_c] +2[R_{ac},R_a]=0$ (II);

(3) is equal to $[R_{dc},R_a]+ [R_{ad},R_c]+[R_{ac},R_d]=0$ (III).

\begin{definition} 
For any algebra $ A $ we call linear operator $ D: A \rightarrow A $ a derivation if 
\begin{equation}
D(ab)=D(a)b+aD(b) \tag{der}
\end{equation}
\end{definition}

\begin{lemma} 
For any algebra $ A $, Let $ D $ be a linear operator $ D: A\rightarrow A $. $ D $ is a derivation if and only if $[R_a,D]=R_{aD}, \forall a \in A$. 
\end{lemma}

\begin{proof} 
By the lebniz law (der), write it into the operator form acts on b $R_aD=R_{aD}+DR_a $, thus $ [R_a,D]=R_{aD} $  
\end{proof}

Remark: in (**), righthand side is symmetric in $ a $ and $ b $. Exchange $ a $ and $ b $ and substitude the left handside (**) we obtain:
\[
R_{(bc)a}-R_{(ac)b}+R_bR_aR_c+R_cR_aR_b-R_aR_bR_c-R_cR_bR_a=0
\]
i.e.
\[
R_{(bc)a}-R_{(ac)b}+[R_b,R_a]R_c+R_c[R_a,R_b]=0
\]
i.e.
\[
R_{(a,c,b)}+[[R_a,R_b],R_c]=0
\]
Or
\[
R_{(a,c,b)}=[R_c,[R_a,R_b]] \tag{IV}
\]
We can write $ (a,c,b)=(ac)b-a(cb)=cR_aR_b-cR_bR_a=c[R_a,R_b]$. By the lemma we have the following corollary:

\begin{corollary}
For any Jordan algebra $ J $, $ a,b \in J $, $ [R_a,R_b] $ is a derivation of $ J $. 
\end{corollary}

\begin{proof} 
$ [R_c,[R_a,R_b]]=R_{(a,c,b)}=R_{c[R_a,R_b]} $. By the lemma $ [R_a,R_b] $ is a derivation.
\end{proof}

\begin{proposition}
For any Jordan algebra $ J=\<a_i | i \in I \> $ 
\[
\operatorname{Mult}(J)=\<R_{a_i}, R_{a_i a_j} \ \text{for} \ a_i, a_j \in J \ \text{where} \ i,j \in I \>
\]
\end{proposition}

\begin{proof} 
By (**), $ R_{(a b) c} $ or $ R_{a (b c)} $ can be reduced into a form $ x \in \< R_{a}, R_{b}, R_{c}, R_{ab}, R_{ac}, R_{bc} \> $. i.e. For a product sequence with length more or equal to $ 2 $, it can always be reduced into a form $ x \in \< R_{a_i}, R_{a_ia_j}| i,j \in I\> $.
\end{proof}

\begin{corollary}
If $ J $ is finitely generated then $ \operatorname{Mult}(J) $ is finitely generated.
\end{corollary}

\begin{corollary}
If $ J=\< a\> $ then $ Mult(J) $ is a commutative associative algebra.
\end{corollary}

\begin{proof} 
As $ Mult(J)=\< R_a, R_{a^2} \> $ and $ [R_a, R_{a^2}]=0 $.  
\end{proof}

\begin{definition} 
An algebra A is called power associative if $ \operatorname{alg}\< a \> $ is associative for any $ a\in A $. $  \operatorname{alg} \<a \> $ this notation means the subalgebra of $ A $ generated by $ a $.    
\end{definition}

\begin{exercise}
Jordan algebras are power associative. (hint: If $ J=\<a \> $ is Jordan then $ \operatorname{Mult}(J) $ is commutative).
\end{exercise}

\begin{proof} 
Let $ A $ be a Jordan algebra. $ \forall a \in A $, $ J=\<a \> $ is a subalgebra of $ A $ generated by $ a $. Hence $ \operatorname{Mult}(J) $ is a commutative associative algebra. $ \forall x,y,z \in J $, $ (x,y,z)=y[R_x, R_z]=0 $. Hence $ J $ is associative. Since $ a \in A $ is arbitrary. $ A $ is power associative.
\end{proof}

\subsection*{Summary of Identities from Lecture 2}

In Lecture 2, the analysis begins with the defining Jordan identity and proceeds through linearization to derive several key relations, particularly in terms of the multiplication operator $R_x$ (where $y R_x = yx$).

\begin{enumerate}
    \item \textbf{The Jordan Identity:}
    The defining axiom for Jordan algebras (assuming commutativity $ab=ba$):
    \[
    (a^2, b, a) = 0 \quad \iff \quad a^2(ba) = (a^2b)a
    \]
    In operator form, this is equivalent to:
    \[
    [R_{a^2}, R_a] = 0 \tag{I}
    \]

    \item \textbf{Partial Linearization:}
    Replacing $a$ with $a + \lambda c$ in the Jordan identity yields coefficients for $\lambda$ and $\lambda^2$, leading to:
    \[
    f(a,b,c) = 2(ac, b, a) + (a^2, b, c) = 0
    \]
    \[
    g(a,b,c) = 2(ac, b, c) + (c^2, b, a) = 0
    \]
    The operator form for $f(a,b,c)=0$ is:
    \[
    [R_{a^2}, R_c] + 2[R_{ac}, R_a] = 0 \tag{II}
    \]

    \item \textbf{Full Linearization:}
    Linearizing $f(a,b,c)$ by replacing $a$ with $a+d$ gives the fully linearized identity:
    \[
    (dc, b, a) + (ac, b, d) + (ad, b, c) = 0
    \]
    In operator form (acting on $b$):
    \[
    [R_{dc}, R_a] + [R_{ad}, R_c] + [R_{ac}, R_d] = 0 \tag{III}
    \]
    Another important operator identity derived from this (acting on $d$) is:
    \[
    R_c R_b R_a + R_a R_b R_c + R_{(ac)b} = R_c R_{ba} + R_a R_{bc} + R_b R_{ac} \tag{**}
    \]

    \item \textbf{Associator and Derivations:}
    By permuting variables in the linearized identities, we obtain a relation linking the associator to the commutator of multiplication operators:
    \[
    R_{(a,c,b)} = [R_c, [R_a, R_b]] \tag{IV}
    \]
    This identity implies that for any $a, b$ in a Jordan algebra, the operator $[R_a, R_b]$ is a \textbf{derivation}.
\end{enumerate}


\section{Lecture 3}

The multiplication algebra $ \operatorname{Mult}(J)=\<R_a: a\in J \> \subseteq \operatorname{End}(J) $ satisfies $ [R_{a^2},R_a] =0$ and $ R_cR_bR_a+R_aR_bR_c+R_{(ac)b}=R_cR_{ba}+R_aR_{bc}+R_bR_{ac} $. 

The universal multiplication envelope is a more universal object. The idea is to forget identities which involve operators and consider an algebra satisfying these relations.  

For $ F $-vector space $ V $ define:
\[
T(V)= F \cdot 1 \oplus V \oplus V \otimes V \oplus \ldots \oplus V^{\otimes n} \oplus \ldots
\] 

In $T(V)$, consider the multiplication given by concatenation. For $u = u_1 \otimes \dots \otimes u_i \in V^{\otimes i}$ and $v = v_1 \otimes \dots \otimes v_j \in V^{\otimes j}$, the product is:  
\[  
uv = u_1 \otimes \dots \otimes u_i \otimes v_1 \otimes \dots \otimes v_j \in V^{\otimes (i+j)}  
\]  

With this multiplication, $T(V)$ is the free associative algebra generated by $V$. It satisfies the following universal property.  

If $A$ is an associative $F$-algebra and $\varphi: V \to A$ is a linear function, then there exists a unique algebra homomorphism $\tilde{\varphi}: T(V) \to A$ such that $\tilde{\varphi}|_V = \varphi$.  

\vspace{1cm}  
\begin{center}  
    \begin{tikzpicture}  
        \node (V) at (0,1.5) {$V$};  
        \node (A) at (4,1.5) {$A$};  
        \node (TV) at (0,0) {$T(V)$};  

        \draw[->] (V) -- node[above] {$\varphi$ (linear)} (A);  
        \draw[->] (V) -- (TV);  
        \draw[->, dashed] (TV) -- node[below right] {$\tilde{\varphi}$ (homo)} (A);  
    \end{tikzpicture}  
\end{center}  
\vspace{1cm} 

If $\{v_1, \dots, v_n, \dots\}$ is a basis of $V$, then $T(V) = F_{\text{free}}\langle v_1, \dots, v_n, \dots \rangle$ is the algebra of non-commutative polynomials in $v_1, \dots, v_n, \dots$.  
  
Let $J$ be a Jordan algebra. We consider the tensor algebra $T(J)$ built on the vector space $J$ and the ideal $I$ of $T(J)$:  
\[
I=\< [a^2,a], a\otimes b\otimes c+c\otimes b\otimes a+(a \circ c)\circ b - a\otimes b \circ c - b\otimes a\circ c - c\otimes a \circ b \>
\]  

For any $ a,b,c \in J $. Denote $U(J) = T(J)/I$. Then $U(J)$ is an associative algebra, which is called the universal multiplication envelope of $J$.  

Observe that $\operatorname{Mult}(J)$ is a homomorphic image of $U(J)$.  
Denote by $R_a$ the element $a+I \in U(J)$ for $a \in J$. We have that $\{R_a : a \in J\}$ satisfies the identities  
\begin{align*}  
    [R_a, R_{a^2}] &= 0 \quad \text{and} \\
    R_{(ac)b} + R_c R_b R_a + R_a R_b R_c &= R_b R_{ac} + R_a R_{bc} + R_c R_{ab}  
\end{align*}  

\begin{remark}[Why $\operatorname{Mult}(J)$ is a homomorphic image of $U(J)$]
Let $J$ be a Jordan algebra over $F$. Consider the tensor algebra
\[
T(J)=\bigoplus_{n\ge 0} J^{\otimes n}
\]
and the ideal $I\lhd T(J)$ generated by the relations
\[
[a^2,a]\quad\text{and}\quad
a\otimes b\otimes c+c\otimes b\otimes a+(a\circ c)\circ b
- a\otimes (b\circ c)- b\otimes (a\circ c)- c\otimes (a\circ b),
\]
for all $a,b,c\in J$ (where $b\circ c\in J\subset T(J)$ is embedded as a tensor
of length $1$). Set
\[
U(J):=T(J)/I,
\]
the \emph{universal multiplication envelope} of $J$.

Assume that the multiplication algebra is defined by
\[
\operatorname{Mult}(J):=\operatorname{alg}\langle R_a,\ R_{a\circ b}\mid a,b\in J\rangle
\subseteq \operatorname{End}_F(J),
\qquad R_a(x):=x\circ a.
\]
Then $\operatorname{Mult}(J)$ is a homomorphic image of $U(J)$ in the following sense:
there exists a surjective $F$-algebra homomorphism
\[
\pi:U(J)\twoheadrightarrow \operatorname{Mult}(J).
\]

\noindent
\textbf{Construction of the map.}
Define an $F$-linear map
\[
\phi:J\to \operatorname{End}_F(J),\qquad \phi(a):=R_a.
\]
By the universal property of the tensor algebra, $\phi$ extends uniquely to an
$F$-algebra homomorphism
\[
\widetilde\phi:T(J)\to \operatorname{End}_F(J)
\]
satisfying $\widetilde\phi|_J=\phi$, hence
\[
\widetilde\phi(a_1\otimes\cdots\otimes a_n)=R_{a_1}\cdots R_{a_n}.
\]
The defining relations of $I$ were chosen precisely so that they hold for the
right multiplication operators on $J$; equivalently,
\[
I\subseteq \ker(\widetilde\phi).
\]
Therefore $\widetilde\phi$ factors through the quotient, yielding an induced
homomorphism
\[
\pi:U(J)=T(J)/I\longrightarrow \operatorname{End}_F(J),
\qquad \pi(\overline{t})=\widetilde\phi(t),
\]
where $\overline{t}$ denotes the class of $t\in T(J)$ modulo $I$.

\noindent
\textbf{Why the image is $\operatorname{Mult}(J)$.}
By construction, $\pi$ contains all right multiplications:
for $a\in J\subset T(J)$,
\[
\pi(\overline{a})=R_a.
\]
Moreover, the cubic relations in $I$ identify (in the quotient $U(J)$) the
operator represented by a tensor word of length $\ge 2$ with expressions
involving Jordan products in $J$ embedded as length-$1$ tensors; consequently,
the image of $\pi$ is closed under passing from words in the generators
to right multiplications by Jordan products. In particular, for any $a,b\in J$,
the induced representation contains $R_{a\circ b}$, and hence contains the
associative algebra generated by $\{R_a, R_{a\circ b}\}$.
Thus
\[
\pi(U(J)) \supseteq \operatorname{alg}\langle R_a,\ R_{a\circ b}\mid a,b\in J\rangle=\operatorname{Mult}(J).
\]
On the other hand, $\pi(U(J))$ is an associative subalgebra of $\operatorname{End}_F(J)$
containing the same generating set, so
\[
\pi(U(J))\subseteq \operatorname{Mult}(J).
\]
Therefore $\pi(U(J))=\operatorname{Mult}(J)$, and $\pi$ is surjective. Equivalently,
\[
\operatorname{Mult}(J)\cong U(J)/\ker(\pi),
\]
i.e.\ $\operatorname{Mult}(J)$ is a homomorphic image of $U(J)$.
\end{remark}

\begin{proposition}
Let $ B $ be a Jordan algebra which contains $ J $, $ \operatorname{Mult}_B(J)=\operatorname{alg}_{\operatorname{End}(B)}\< R_a^B , a \in J \> $. Then $ \operatorname{Mult}_B(J) $ is a homomorphisc image of $ U(J) $. Just consider $ \varphi^B: U(J) \rightarrow \operatorname{Mult}_B(J) $ s.t. $ a+I \mapsto R_a^B $.
\end{proposition}

\begin{proof}[Well-definedness of $\varphi^B$]
Let $J$ be a Jordan algebra over a field $F$, and let $B$ be a Jordan algebra containing $J$ as a subalgebra.
For each $a\in J$, denote by
\[
R_a^B\in \operatorname{End}_F(B),\qquad R_a^B(x):=x\circ_B a
\]
the right multiplication operator in $B$.

Recall that the tensor algebra $T(J)=\bigoplus_{n\ge 0} J^{\otimes n}$ is the free associative $F$-algebra on the
vector space $J$. Consider the $F$-linear map
\[
\psi:J\longrightarrow \operatorname{End}_F(B),\qquad a\longmapsto R_a^B.
\]
By the universal property of $T(J)$, there exists a unique associative algebra homomorphism
\[
\widetilde{\psi}:T(J)\longrightarrow \operatorname{End}_F(B)
\]
such that $\widetilde{\psi}(a)=R_a^B$ for all $a\in J\subset T(J)$.

Let $I\triangleleft T(J)$ be the (two-sided) ideal defining the universal multiplication envelope
\[
U(J):=T(J)/I,
\]
where $I$ is generated by the standard operator identities that hold for right multiplications
in every Jordan algebra (e.g.\ the relations encoding $[R_{a^2},R_a]=0$ and the polarized
Jordan identity, as in the definition of $U(J)$).

\noindent\textbf{Claim.} $I\subseteq \ker(\widetilde{\psi})$.

\noindent\emph{Justification.}
Each generator $r\in I$ is a noncommutative polynomial expression in the symbols $a\in J$
that becomes $0$ after interpreting each $a$ as the operator $R_a$ in any Jordan algebra.
Since $B$ is a Jordan algebra and each $R_a^B$ is a right multiplication in $B$, the same
operator identities hold in $\operatorname{End}_F(B)$; hence $\widetilde{\psi}(r)=0$ for every generator $r$ of $I$.
By ideality and multiplicativity of $\widetilde{\psi}$, it follows that $\widetilde{\psi}(I)=0$, i.e.\ $I\subseteq \ker(\widetilde{\psi})$.

Therefore, $\widetilde{\psi}$ factors through the quotient $T(J)/I$:
there exists a unique associative algebra homomorphism
\[
\varphi^B:U(J)=T(J)/I\longrightarrow \operatorname{End}_F(B)
\]
such that $\varphi^B\circ \pi=\widetilde{\psi}$, where $\pi:T(J)\to T(J)/I$ is the canonical projection.
Equivalently,
\[
\varphi^B(a+I)=R_a^B\qquad(a\in J).
\]

\noindent\textbf{Well-definedness.}
If $x+I=y+I$ in $U(J)$, then $x-y\in I\subseteq \ker(\widetilde{\psi})$, hence
\[
\widetilde{\psi}(x)=\widetilde{\psi}(y).
\]
Applying $\varphi^B\circ\pi=\widetilde{\psi}$ gives
\[
\varphi^B(x+I)=\varphi^B(y+I),
\]
so $\varphi^B$ is well-defined.

Finally, by construction the image of $\varphi^B$ is the associative subalgebra of $\operatorname{End}_F(B)$
generated by $\{R_a^B\mid a\in J\}$, i.e.
\[
\operatorname{Im}(\varphi^B)=\operatorname{alg}_{\operatorname{End}(B)}\langle R_a^B\mid a\in J\rangle
=\operatorname{Mult}_B(J).
\]
Thus $\varphi^B:U(J)\twoheadrightarrow \operatorname{Mult}_B(J)$ is a surjective homomorphism.
\end{proof}

\begin{exercise}
Prove that there exists $ J \subseteq B $ such that $ \operatorname{Mult}_B(J) \simeq U(J) $.  
\end{exercise}

\begin{proof}
Let $J$ be a Jordan algebra over a field $F$ with $\operatorname{char}F\neq 2$, and let
\[
U(J)=T(J)/I
\]
be its universal multiplication envelope. Denote by
\[
\iota:J\longrightarrow U(J),\qquad a\longmapsto a+I
\]
the canonical linear map, and write $1$ for the unit of the associative algebra $U(J)$.

\noindent\textbf{Step 1: A faithful $U(J)$-module.}
Let $M:=U(J)$ viewed as the \emph{right regular} $U(J)$-module via
\[
m\cdot u:=mu\qquad(m,u\in U(J)).
\]
Equivalently, let $\rho^{\mathrm r}:U(J)\to \operatorname{End}_F(M)$ be the right regular representation
$\rho^{\mathrm r}(u)(m)=mu$.

\noindent\textbf{Step 2: The split null extension.}
On the vector space $B:=J\oplus M$ define a bilinear, commutative product by
\begin{equation}\label{eq:split-null}
(x+m)\circ (y+n)
:= (x\circ_J y)\;+\; m\cdot \iota(y)\;+\; n\cdot \iota(x),
\qquad (x,y\in J,\ m,n\in M),
\end{equation}
and declare $M\circ M=0$ (which is already encoded by \eqref{eq:split-null}).

Then $J$ embeds as a subalgebra of $B$ via $x\mapsto (x,0)$, since
\[
(x,0)\circ (y,0)=(x\circ_J y,\,0).
\]

\noindent\textbf{Step 3: $B$ is a Jordan algebra.}
By the defining property of $U(J)$, the elements $\iota(a)$ ($a\in J$) satisfy in $U(J)$ precisely the
(operator) relations imposed by the Jordan identities for right multiplications (e.g.\ the relations
encoding $[R_{a^2},R_a]=0$ and the polarized Jordan identity, as in the definition of $U(J)$).
Hence the operators $\rho^{\mathrm r}(\iota(a))\in \operatorname{End}(M)$ satisfy the same relations, so $M$ is a
(right) Jordan $J$-module via
\[
m\cdot a:=m\cdot \iota(a)\qquad (m\in M,\ a\in J).
\]
Therefore the split null extension $B=J\oplus M$ with product \eqref{eq:split-null} is a Jordan algebra.

\noindent\textbf{Step 4: Construct the map $\Phi$ and prove it is an isomorphism.}
Consider the associative subalgebra
\[
\operatorname{Mult}_B(J):=\operatorname{alg}_{\operatorname{End}(B)}\langle R^B_{(a,0)}\mid a\in J\rangle
\subseteq \operatorname{End}(B),
\]
where $R^B_{(a,0)}$ denotes right multiplication by $(a,0)\in B$.
By the universal property of $U(J)$ (equivalently, by Proposition~3.1 applied to this $B$), there is a unique
associative algebra homomorphism
\[
\Phi:U(J)\longrightarrow \operatorname{Mult}_B(J)
\]
such that
\[
\Phi(\iota(a))=R^B_{(a,0)}\qquad (a\in J).
\]
It is surjective by definition of $\operatorname{Mult}_B(J)$.

To prove injectivity, restrict to the invariant subspace $M\subseteq B$ (identified with $\{0\}\oplus M$).
For $m\in M$ and $a\in J$, using \eqref{eq:split-null} we compute
\[
R^B_{(a,0)}(0,m)=(0,m)\circ (a,0)=(0,\,m\cdot \iota(a)).
\]
Thus the action of $\Phi(\iota(a))$ on $M$ coincides with the right regular action of $\iota(a)$ on $M$.
By multiplicativity, for every $u\in U(J)$ we have
\[
\Phi(u)\big|_{M}=\rho^{\mathrm r}(u).
\]
Hence, if $\Phi(u)=0$, then $\rho^{\mathrm r}(u)=0$, so $m u=0$ for all $m\in M$; in particular
\[
u=1\cdot u=0.
\]
Therefore $\ker\Phi=0$ and $\Phi$ is injective.

Consequently $\Phi$ is an isomorphism, and for the Jordan algebra $B=J\oplus U(J)$ constructed above we obtain
\[
\operatorname{Mult}_B(J)\cong U(J),
\]
with $J$ embedded in $B$ as the subalgebra $J\oplus 0$.
\end{proof}

Now we introduce nilpotent, nil and solvable algebras.

Let $A$ be an algebra. If $B, C \subseteq A$ are subspaces, we denote by  
\[ BC = \left\{ \sum_{i=1}^n b_i c_i : b_i \in B, c_i \in C, n \in \mathbb{Z}, n>0 \right\}. \]  

In this way we may consider, $A^2, A^3, A A^2, \dots$. If $A$ is an associative algebra, we call $A$ nilpotent if there exists some integer $n>0$ such that $A^n=0$, where $A^n$ consists of sums of products of $n$ elements of $A$. If $A$ is not associative, we may have for example $A^2 A=0$ but $A A^2 \neq 0$.  
This way we have different types of nilpotency. The most important are nilpotency and solvability.  

Define $A^1 = A$ and  
\[ A^n = \sum_{i=1}^{n-1} A^i A^{n-i} \quad \text{if } n > 1 \]  
and  
\[ A^{(0)} = A \quad \text{and} \quad A^{(n+1)} = (A^{(n)})^2 \quad \text{if } n > 0. \]  
Observe that  
\[ A^{n+1} \subseteq A^n \quad \text{and} \quad A^{(n)} \subseteq A^{(n-1)} \]  
for any $n > 0$.  

\begin{definition} 
An algebra $ A $ is called nipotent if there exists $ n \in \mathbb{N} $ such that $ A^n =0 $. The least $ n $ with $ A^n =0 $ is called the nilpotancy index of $ A $. A is called solvable if $ \exists n \in \mathbb{Z}_{>0} $ such that $ A^{(n )}=0 $. The least $ n $ is called the solvability index of $ A $. A is called a nil algebra if $ \operatorname{alg}\< a \> $ is nilpotent for any $ a\in A $.         
\end{definition}

Any nilpotent algebra is nil.

Any nilpotent algebra is solvable since $ A^{(n )} \subseteq A^{2^n} , n\ge 1$.

In the case of associative algebras, concepts of solvable and nilpotent algebras coincide (but a nil algebra is not necessarily nilpotent). But for non-associative algebras, this is not the case.  
Studying relations between these different classes is one of the key questions for non-associative algebras.  

\begin{theorem}  
Let A be a finite-dimensional associative algebra. Then there exists an ideal in A, denoted by $\operatorname{Nil} A$, which contains all nil ideals of $A$ and the following conditions are satisfied:  
\begin{enumerate}  
    \item $\operatorname{Nil} A$ is nilpotent.  
    \item $A / \operatorname{Nil} A \simeq A_1 \oplus \dots \oplus A_m$, and each $A_i$ is simple.  
    \item Each simple algebra is isomorphic to $M_n(D)$, where $D$ is a division algebra for some $n \ge 1$.  
    \item If $A / \operatorname{Nil} A$ is separable (in particular, $ \operatorname{char}F=0 $ and $ A $ is of finite dimensional, in this case $ F $ is perfect and $ \operatorname{Nil}(A/\operatorname{Nil}A)=0 $, $ S $ is reduced) then $A \simeq S \oplus \operatorname{Nil} A$, where $S \simeq A / \operatorname{Nil} A$.  
\end{enumerate}  
The ideal $\operatorname{Nil} A$ is denoted the nil radical of A.
\end{theorem}  

Then in the case of associative algebras of finite dimension can be studied separately: semi-simple and nilpotent.  
Wedderburn-Molien theorem was generalized by Artin for the case of the rings which satisfy chain conditions on ideals. But in general, associative algebras do not contain a maximal nilpotent ideal (when $ \dim A = \infty $), such that the quotient by this ideal does not have non-trivial nilpotent ideals.  

The idea to generalize a notion of nilpotent radical:  
To give an ideal $R(A)$ of an associative algebra $A$, satisfying a certain property, containing all ideals in $A$ with this property in a way that $A/R(A)$ does not contain non-zero ideals satisfying this property in order to describe algebras $R(A)$ and $A/R(A)$ as ``radical algebras'' and ``semi-simple algebras''.  

As a consequence of Molien-Wedderburn we have that any finite-dimensional associative nil algebra is nilpotent.  
Next, we will show the same for Jordan algebras.  

\begin{remark}  
If $A$ is a power associative algebra and $a \in A$ then $\operatorname{alg}\langle a \rangle$ is nilpotent $\iff \exists n \text{ s.t. } a^n=0$.  
\end{remark} 

\begin{lemma} 
Let $ A $ be a power associative algebra. $ I \triangleleft A $. If $ I $ is nil ideal and $ A/I $ is nil algebra then $ A $ is a nil algebra.     
\end{lemma}

\begin{proof} 
Let $ \bar{a} \in A/I $, there exists $ n>0 $ s.t. $ (\bar{a})^n=0 $. Then $ a^n \in I $ and $ (a^n)^m=0 $ for some $ m > 0 $.
\end{proof}

\begin{lemma} 
$ I,J $ are to nil ideals in power associative algebra $ A $, then $ I+J $ is nil.   
\end{lemma}

\begin{proof} 
$ I $ is nil, then $ \frac{I}{I \cap J} $ is nil. By $\frac{I}{I \cap J} \simeq \frac{I+J}{J}  $. $ \frac{I+J}{J} $ is nil. And $ J $ is nil. By former lemma $ I+J $ is nil.      
\end{proof}

\begin{corollary}
$ I_1 + \ldots + I_n $ is nil as long as $ I_j $ is nil for all $ j $.
\end{corollary}

\begin{proposition}
Let $ A $ be power associative algebra. Then there exists maximal nil ideal (nil radical) $ \operatorname{Nil}A $ (any nil ideal in A is contained in $ \operatorname{Nil}A $ and $ \operatorname{Nil}(A/ \operatorname{Nil}A)=0 $ ).
\end{proposition}

\begin{proof} 
Set $ N = \sum I $ where $ I  $ nil ideal in $ A $. We will show that $ N $ is in fact the nil radical $ \operatorname{Nil}A $. $ N $ is nil since any $ a \in N $ can be writen as a finite sum of element of $ I_1 + \ldots +I_n $ which is by corollary nil. By definition of $ N $ it contains all nil ideals of $ A $. Suppose I is some nil ideal in $ A/\operatorname{Nil}A $. By corespondence theorem we have $ J \triangleleft A $ and $ \operatorname{Nil}A \subseteq J $. $ J / \operatorname{Nil}A=I$ is nil. Then $ J $ is nil. Since $ J \subseteq \operatorname{Nil}A $ we have $ J=\operatorname{Nil}A $ and $ J / \operatorname{Nil}A =0 $. Thus $ \operatorname{Nil}(A/ \operatorname{Nil}A)=0 $                  
\end{proof}

As a consequence we showed that the nil radical exists for any associative $ A $ or Jordan algebra $ J $.

\section{Lecture 4}

\begin{lemma} 
Let $ A $ be a commutative algebra. Then $ A $ is nilpotent if and only if $ \operatorname{Mult}A $ is nilpotent.   
\end{lemma}

\begin{proof} 
($ \Rightarrow $ ). If $ A $ is nilpotent, then $\exists n, \forall a,b_i \in A $, $ A^n \ni (ab_1)b_2 \ldots b_{n-1} =a R_{b_1}R_{b_2}\ldots R_{b_{n-1}} =0 $. Thus $ \operatorname{Mult}A $ is nilpotent.

($ \Leftarrow $ ). Let us show that any element of $ A^{2^n} $ can be writen as sum of $ aR_{b_1}R_{b_2}\ldots R_{b_{n}} , a,b_i \in A$.

We prove it by induction on $ n $. For $ n=1 $, $ ab \in A $, $ ab=aR_b $.

Suppose it for $ k=n-1 $ is true. Consider element of $ A^{2^n} $. Take $ a_1,a_2 \in A^{2^{n-1}} $, regard $ a_1 \in A $ and $ a_2 \in A^{2^{n-1}} $, and let $ a_2=a_1^{\prime }R_{b_1}R_{b_2}\ldots R_{b_{n-1}}$ then $ A^{2^n}\ni a_1 (a_1^{\prime }R_{b_1}R_{b_2}\ldots R_{b_{n-1}}) =(a_1^{\prime }R_{b_1}R_{b_2}\ldots R_{b_{n-1}}) a_1= (a_1^{\prime }R_{b_1}R_{b_2}\ldots R_{b_{n-1}}) R_{a_1}$ use the indution step we finished. (Remark, this is not obvious since our $ A $ is just asssume to be commutative but no need associative so abc may not equal to $ a(bc) $, i.e., if $ bc=0 $ then (ab)c may not 0).     
\end{proof}

\begin{exercise}
For $ J $ is a Jordan algebra, prove if $ (\operatorname{Mult}(J))^{n}=0 $ then $ J^{2n-1}=0 $.   
\end{exercise}

\begin{tcolorbox}[colback=gray!20, colframe=black] 
Note: I think it should be changed into: if $ J $ is unital, then if $ (\operatorname{Mult}(J))^{n}=0 $ then $ J^{2n-1}=0 $; if $ J $ is not necessary unital, then if $ (\operatorname{Mult}(J))^{n}=0 $ then $ J^{2n}=0 $.
\end{tcolorbox}

\begin{proof}
Consider a word $ x_1, \ldots x_{2n-1} \in J^{2n-1} $. It may have arbitrary order of operations. For example, when $ n=2 $, $ abcd $ or $ ab(cd) $ or $ a((bc)d) $ and so on. All we know is the ``length'' of this word is $ 2n-1 $. If $ 1 \in J $, we rewrite it into:
\[
1 R_{x_1, \ldots x_{2n-1}}
\]      
Now since $R_{x_1, \ldots x_{2n-1}} \in \operatorname{Mult}J  $ and $  \operatorname{Mult}J  $ is generated by $ R_a $ where $ a $ is a word of length 1 or 2. We have 
\[
R_{x_1, \ldots x_{2n-1}}=\sum_i \alpha_i R_{x_{i_1}}, \ldots R_{x_{i_{k_i}}}
\]    
And $ k_i \ge \lceil\frac{2n-1}{2}\rceil=n $. Hence if $ (\operatorname{Mult}(J))^{n}=0 $ then $ J^{2n-1}=0 $ when $ J $ is unital. 

If $ 1 \notin J $ then:
\[
xR_{x_1, \ldots x_{2n-1}}
\] 
is 0 forall $ x \in J $, which means that if $ (\operatorname{Mult}(J))^{n}=0 $ then $ J^{2n}=0 $. 
\end{proof}

\begin{lemma} 
Let $ J =B+Fv$ be Jordan algebra, $ B $ is an ideal, $ J^2 \subseteq B $. If $ U(B)^+ $ is nilpotent then $ U(J)^+ $ is nilpotent. Recall $ U(J)^+=T(R_a)/ \<\text{the idal generated by }[R_{a^2}, R_a], R_a R_b R_c +R_c R_b R_a +R_{(ac)b}-R_aR_{bc}-R_bR_{ac}-R_cR_{ba} \> $. $ U(B)^+=U_J(B)^+=\operatorname{alg}_{U(J)^+}\< R_b, b\in B \> $. (Remark: From here $ U(J)^+ $ is the associative algebra generated by $ \<R_a, a \in J \> $, in other words $ U(J)=F\cdot 1\oplus U(J)^+ $).      
\end{lemma}

\begin{proof} 
$ U(B)^+ $ is nilpotent. We check that 
\[
(U(J)^+)^{3m} \subseteq (U(B)^+)^{m}+(U(B)^+)^{m}U(J)^+
\] 
as subset for any $ m \ge 1 $ by induction.

Consider $ R_aR_bR_c \in U(J)^+  $.
\begin{itemize}
\item if $ a \in B $ then $ R_aR_bR_c \in U(B)^+U(J)^+ \subset U(B)^++U(B)^+U(J)^+ $.
\item if $ a=v $, $ b,c \in B $. $ R_vR_bR_c =-R_cR_bR_v-R_{(vc)b}+R_vR_{bc}+R_bR_{vc}+R_cR_{bv} \in  U(B)^++U(B)^+U(J)^+ $. (For $ R_vR_{bc} $ consider $ [R_{bc}, R_v]+[R_{vb}, R_c]+[R_{vc},R_b]=0 $).
\item if $ a=b=v $, $ c \in B $. $ R_vR_vR_c =-R_cR_vR_v-R_{(vc)v}+R_vR_{vc}+R_vR_{vc}+R_cR_{vv} \in  U(B)^++U(B)^+U(J)^+ $. (For $ R_vR_{vc} $ consider $ [R_{vc}, R_v]+[R_{vv}, R_c]+[R_{vc},R_v]=0 $ in which we have $ R_{v^2}R_c $, do not forget $ J^2 \in B $).
\item if $ a=c=v $, $ b \in B $. $ R_vR_bR_v =-R_vR_bR_v-R_{(vv)b}+R_vR_{bv}+R_bR_{vv}+R_cR_{bv} \in  U(B)^++U(B)^+U(J)^+ $. (For $ R_vR_{bv} $ consider $ [R_{bv}, R_v]+[R_{vb}, R_v]+[R_{vv},R_b]=0 $ in which $ R_vR_{vb} $ consider $ 2[R_v,R_{vb}] +[R_b,R_{v^2}]=0$).  
\item if $ a=b=c=v $. $ 2R_vR_vR_v =-R_{(vv)v}+R_vR_{vv}+R_vR_{vv}+R_vR_{vv} \in  U(B)^++U(B)^+U(J)^+ $. ($ R_vR_{v^2}=R_{v^2}R_v $).
\end{itemize}
Therefore, $ (U(J)^+)^{3} \subseteq U(B)^++U(B)^+U(J)^+ $.  

Suppose it works for $ m $, then for $ m+1 $, $ {(U(J)^+)}^{3m+3}={(U(J)^+)}^{3m}{(U(J)^+)}^{3} \subseteq {((U(B)^+)}^{m}+{(U(B)^+)}^{m}{(U(J)^+)}(U(B)^+ +U(B)^+U(J)^+) \subseteq {U(B)^+}^{m+1}+{U(B)^+}^{m+1}U(J)^+$. 
\end{proof}

\begin{theorem}[Albert]  
Let $J$ be a Jordan algebra of finite dimension. If $J$ is nil, then the universal multiplication enveloping algebra $U(J)^+$ is nilpotent.  
\end{theorem}  

\begin{proof} 
Let $B \subseteq J$ be a maximal subalgebra with the property that $U(B)^+$ is nilpotent. Such a subalgebra exists since $J$ has finite dimension and the set $S$ of subalgebras $A$ of $J$ with $U(A)^+$ nilpotent is not empty (since $\{0\} \in S$).  

Suppose $B \neq J$. Observe that the multiplication algebra $\operatorname{Mult}_J(B)$ is a homomorphic image of $U(B)^+$ and thus is nilpotent. Then we have the following chain of subspaces, which must terminate at $\{0\}$ due to the nilpotency of $\operatorname{Mult}_J(B)$:  
\[ J \supset JB \supseteq (JB)B \supseteq \dots \supseteq ((\dots(JB)B)\dots)B = \{0\} \]  
Since $B \neq J$, this implies there exists an element $a \in J$ such that $a \notin B$ and $aB \subseteq B$.  

We have two following possibilities for $a^2$:  
\begin{enumerate}  
    \item $a^2 B \subseteq B$.  
    \item There exists $b \in B$ such that $a^2 b \notin B$.  
\end{enumerate}  

In the first case, the subalgebra $B$ is invariant under action by $R_a$ and $R_{a^2}$. Since $R_{a^k}$ belongs to $\operatorname{alg}\langle R_a, R_{a^2} \rangle$ for each $k \ge 1$, then $a^k B \subseteq B$ for any $k \ge 1$. As $a \notin B$ and $J$ is nil, there exists an integer $m > 1$ such that $a^m \in B$ but $a^{m-1} \notin B$. Put $v = a^{m-1}$. Then $v^2 = a^{2m-2} \in B$ and $vB = a^{m-1}B \subseteq B$. $B_1 = B + Fv$ satisfies Lemma about $J = B + Fv$. Then $U(B_1)^+$ is nilpotent, which leads to contradiction since $B$ was chosen maximal with respect to property being nilpotent. Thus case 1 does not happen. 

Next we show that if $aB \subseteq B$ then  
\begin{enumerate}
    \item $(a^2B)B \subseteq B$;  
    \item $(a, a^2B, B) \subseteq B$ and  
    \item $(a^2B)(a^2B) \subseteq B$  
\end{enumerate}  

Denote $x \equiv y \iff x - y \in B$. Let $b, b_1 \in B$, then using that the associator is a derivation in the first argument (For $a, b \in J, [R_a, R_b]$ is a derivation of $J$) we have  
\begin{align*}  
    (a^2b)b_1 &= (a,a,b)b_1 + (a(ab))b_1 \equiv (a,a,b)b_1 \\
    &= (a, ab_1, b) - a(a, b_1, b) \in B, \quad \text{then (i)}  
\end{align*}  

For the second inclusion we use Jordan identity and (i)  
\begin{align*}  
    (a, a^2b, b_1) &= (a(a^2b))b_1 - a((a^2b)b_1) \\
    &= (a^2(ab))b_1 - a((a^2b)b_1) \in B.  
\end{align*}  

For the third formula we have  
\begin{align*}  
    (a^2b)(a^2b_1) &= (a(ab))(a^2b_1) + (a,a,b)(a^2b_1) \\
    &\equiv (a,a,b)(a^2b_1) \quad (\text{using (i)}) \\
    &= (a, a(a^2b_1), b) - (a, a^2b_1, b)a \quad (\text{by derivation}) \\
    &\in B \quad \text{by (ii)}  
\end{align*}  

To prove case 2; there exists $b \in B$ such that $a^2b \notin B$. Let $a_1 = a^2b$. Then $a_1B \subseteq B$ (by (i)) and $a_1^2 \in B$ by (iii). Such element is then satisfying case 1 which is impossible.  
$\implies B=J$ or equivalently $U(J)^+$ is nilpotent.  
\end{proof}

\section{Lecture 5}

\begin{corollary}
Let $ J $ be finite demensional Jordan algebra and $ a \in J $ is nilpotent then $ R_a $ is nilpotent in $ U(J) $.   
\end{corollary}

\begin{proof} 
$ a \in J $ is nilpotent, $ J^{\prime} =\operatorname{alg}_J\< a\>=\mathbb{F}a+\mathbb{F}a^2 + \ldots +\mathbb{F}a^n$. $ J^{\prime } $ is finite dimensional and nil. By Albert theorem $ U(J^{\prime }) $ is nilpotent. Which means $ \operatorname{Mult}_J(J^{\prime}) $ is nilpotent. Thus $ R_a $ is nilpotent element.
\end{proof}

\begin{corollary}
Every finite dimentional nil Jordan algebra is nilpotent.
\end{corollary}

\begin{proof} 
By Albert theorem $ U(J) $ is nilpotent, hence $ \operatorname{Mult}(J) $ is nilpotent. By lemma 4.1 $ J $ is nilpotent.  
\end{proof}

\begin{corollary}
Every finite dimentional solvable Jordan algebra is nilpotent.
\end{corollary}

\begin{proof} 
By the definition of solvability, $ J $ is nil. Hence nilpotent by the previous corollary.
\end{proof}

\begin{corollary}  
If $J$ is a Jordan algebra of finite dimension then $\operatorname{Nil}(J)$ is nilpotent and $\operatorname{Nil}(J/\operatorname{Nil}(J)) = \{0\}$.  
\end{corollary}  

\begin{proof} 
$ \operatorname{Nil}(J) $ is nilpotent since $ \operatorname{Nil}(J) $ is a finite dimensional nil Jordan algebra. And $\operatorname{Nil}(J/\operatorname{Nil}(J)) = \{0\}$ follows from the proof of Proposition 3.9. 
\end{proof}

\begin{exercise}  
Let $A$ be a finite dimensional algebra. Show that $A$ has a solvable ideal $\operatorname{Sol} A$ which contains all solvable ideals in $ A $. Moreover $ \operatorname{Sol}(A/\operatorname{Sol} A)=0 $. Hence there exists a solvable radical for any finite dimensional algebra.  
\end{exercise} 

To prove this exercise, we need some lemmas, just as we did in the proof of Proposition 3.9.

Recall: For an ideal $I$ of an algebra $A$, set $I^{(1)} = I$ and $I^{(n+1)} = I^{(n)} I^{(n)}$ (the ideal generated by all products of two elements of $I^{(n)}$). An ideal $I$ is solvable if $I^{(m)} = 0$ for some $m$.  

\begin{lemma}  
If $I, J \trianglelefteq A$ are ideals, then  
\[ (I+J)^2 \subseteq I^2 + J^2 + (I \cap J). \]  
\end{lemma}  
\begin{proof}  
For $x \in I+J$ and $y \in I+J$, bilinearity gives $xy$ as a sum of terms in $I^2, J^2, IJ, JI$. But $IJ \subseteq I \cap J$ and $JI \subseteq I \cap J$.
\end{proof} 

\begin{lemma}  
If $I$ and $J$ are solvable ideals, then $I+J$ is a solvable ideal.  
\end{lemma}  
\begin{proof}  
Let the derived lengths be $s, t \ge 0$ with $I^{(s)} = 0$ and $J^{(t)} = 0$. By Lemma 5.6,  
\[ (I+J)^{(2)} = (I+J)^2 \subseteq I^2 + J^2 + (I \cap J) = I^{(2)} + J^{(2)} + (I \cap J). \]  
The three ideals on the right are solvable: $I^{(2)}$ and $J^{(2)}$ have derived lengths at most $s-1$ and $t-1$ respectively, and $I \cap J$ is contained in both $I$ and $J$, hence it is solvable. By induction on the sum of derived lengths $s+t$, their sum is solvable. Therefore, $I+J$ is solvable.  
\end{proof}  

 
Let $\operatorname{Sol} A$ be the sum of all solvable ideals of $A$. Since $A$ is finite-dimensional, $ \operatorname{Sol}A $ exists. By repeated application of Lemma 5.7, $\operatorname{Sol} A$ is itself a solvable ideal. By its construction, it contains every solvable ideal of $A$.  

For any ideal $K \trianglelefteq A$ and any $n \ge 0$, one has $(B/K)^{(n)} = (B^{(n)} + K)/K$ for any ideal $B \supseteq K$. If $B/K$ is a solvable ideal of $A/K$, then $(B/K)^{(m)} = 0$ for some $m$, which implies $B^{(m)} \subseteq K$. If $K$ is solvable, then $B^{(m)}$ is solvable, which implies that $B$ is solvable.   

Apply this with $K = \operatorname{Sol} A$. If $\operatorname{Sol}(A/\operatorname{Sol} A) \neq 0$, let's pick a non-zero solvable ideal $B/ \operatorname{Sol} A \trianglelefteq A/\operatorname{Sol} A$. Then by the argument above, $B$ must be a solvable ideal of $A$. By the definition of $\operatorname{Sol} A$, $B$ must be contained in $\operatorname{Sol} A$. This implies $B/\operatorname{Sol} A = 0$, a contradiction. Thus, $\operatorname{Sol}(A/\operatorname{Sol} A) = 0$.  

\begin{remark}
By proposition, any power-associative algebra (of any dimension) admits a nil radical. In the case of any algebra, it is possible to consider nilpotent elements and nil algebras. We say that an element $a$ in an algebra $A$ is nilpotent if there exists a non-associative monomial $v(x)$ such that $v(a)=0$. In this case we have a nil ideal $\operatorname{Nil} A$ containing all nil ideals of $A$ such that $\operatorname{Nil}(A/\operatorname{Nil} A) = 0$.  
\end{remark}    

Now we talk about non-degenerated Jordan algebra.

We saw that Jordan algebras have a nil radical. If $J$ is finite-dimension, then $\operatorname{Nil}(J)$ is nilpotent. We would like to prove a Molien-Wedderburn theorem for Jordan algebras of finite dimension with zero nil radical.  

We prove that every Jordan algebra of finite dimension with $\operatorname{Nil}(J) = 0$ is a direct sum of simple Jordan algebras, assuming that every Jordan algebra has an identity element.  

Let $J$ be a Jordan algebra. For $b \in J$, define an operator $U_b$ by:  
\[ U_b = 2R_b^2 - R_{b^2}  \]  
That is, for any $a \in J$:  
\[ aU_b = 2(ab)b - ab^2 \]  

\begin{remark}  
If $J = A^{(+)}$, we have $aU_b = bab$ for all $a, b \in J$.  
\begin{align*}  
    4aU_b &= 2((ab+ba)b + b(ab+ba)) - 2(ab^2+b^2a) \\
    &= 2(ab^2 + bab + bab + b^2a) - 2ab^2 - 2b^2a \\
    &= 4bab  
\end{align*}  
\end{remark}  

Given a Jordan algebra $J$ and elements $a,b,x \in J$, we introduce the following notation  
\[ \{axa\} = xU_a \]  
and define  
\[ U_{a,b} = \frac{1}{2}(U_{a+b} - U_a - U_b) \]  

\begin{exercise}
Verify $ U_{a,b}=R_aR_b+R_bR_a-R_{a  b} $. 
\end{exercise}

\begin{proof} 
\[  
U_{a+b}  
= 2R_{a+b}^2 - R_{(a+b)^2}  
= 2(R_a+R_b)^2 - R_{a^2 + 2ab + b^2}.  
\]  
Expanding gives  
\[  
U_{a+b}  
= 2(R_a^2 + R_aR_b + R_bR_a + R_b^2) - (R_{a^2} + 2R_{ab} + R_{b^2}).  
\]  
Grouping terms,  
\[  
U_{a+b}  
= (2R_a^2 - R_{a^2}) + (2R_b^2 - R_{b^2})  
  + 2\bigl(R_aR_b + R_bR_a - R_{ab}\bigr)  
= U_a + U_b + 2\bigl(R_aR_b + R_bR_a - R_{ab}\bigr).  
\]  
Therefore,  
\[  
U_{a,b} = \frac{1}{2}\bigl(U_{a+b} - U_a - U_b\bigr)  
= R_aR_b + R_bR_a - R_{ab}.  
\]  
\end{proof}

\begin{remark}
If $J = A^{(+)}$, we have $ xU_{a,b}=\frac{1}{2}(axb+bxa) $.
\end{remark}

we denote by  
\[ \{axb\} = xU_{a,b} \]  
$\{axb\}$ is called a triple product of the Jordan algebra. 

\begin{corollary}
$ \{axb\}=\{bxa\} $ since $ U_{a,b}=U_{b,a} $. 
\end{corollary}

\begin{corollary}
$ xU_{a+b}=\{(a+b)x(a+b)\}=\{axa\}+\{axb\} +\{bxa\}+\{bxb\}=xU_a+xU_b+2U_{a,b}$. 
\end{corollary}

\begin{definition}  
An element of a Jordan algebra $J$, $a \in J$, is called an absolute zero divisor if $U_a=0$ which means $ \forall b \in J $, $ bU_a=0 $. We call a Jordan algebra $J$ non-degenerate if $J$ does not contain non-zero absolute zero divisors.  
\end{definition}  

\begin{proposition}
Let $ J $ be a Jordan algebra, then $ U_{aU_b}=U_bU_aU_b, U_{a^2}=U_a^2 , \forall a,b \in J$.  
\end{proposition}

\begin{proof} 
We show this property for special Jordan algebras. Further, we will see that this is sufficient. Let $J \subseteq A^{(+)}$, with $A$ an associative algebra. Then  
\[ x U_{aU_b} = x U_{bab} = babxbab = b(a(bxb)a)b = x U_b U_a U_b \]  
for any $a,b,x \in J$.

Further, 
\[ bU_{a^2} = a^2ba^2 = a(aba)a = bU_a^2 \quad \text{for every } a,b \in J. \]  
\end{proof}

\begin{theorem}[Shirshov-Cohn]  
Every Jordan algebra generated by two elements is special.  
\end{theorem} 

We will prove this theorem later.

\begin{definition} 
$ z(J)=\< a \in J | U_a=0  \> $ vector space generated by absolute zero divisors.  
\end{definition}

\begin{definition}  
The Unital hull of an algebra $A$ is the algebra $A^\# = A \oplus F\cdot 1$, with multiplication defined as $a \cdot b = ab$ for $a,b \in A$, and $a \cdot 1 = 1 \cdot a = a$, $ 1 \cdot 1=1 $ for any $a \in A^\#$.  
\end{definition}

\begin{lemma} 
$ J $ a Jordan algebra, $ a $ an absolute zero divisor, then for any $ b \in J^{\sharp} $, $ aU_b $ is an absolute zero divisor in $ J $ as well.    
\end{lemma}

\begin{proof} 
Since $ J $ is ideal in $ J^{\sharp} $, $ aU_b=2(a \cdot b)\cdot b - a \cdot b^2 \in J $. $ \forall x \in J $, $ x U_{aU_b} =x U_b U_a U_b \in J U_a U_b =0$.     
\end{proof}

\begin{corollary}
$ z(J) $ is an ideal in $ J $.  
\end{corollary}

\begin{proof} 
It is enough to prove that if $a \in J$ is an absolute zero divisor and $b \in J$, then $ab \in z(J)$.  
By the previous lemma, for $1+b \in J^\#$, we know that $aU_{1+b} \in z(J)$.  
We have the expansion $aU_{1+b} = aU_1+aU_b+2aU_{1,b}$.  
Since $z(J)$ is a vector space and $a \in z(J)$ (by hypothesis) and $aU_b \in z(J)$ (by the previous lemma), it follows that  
\[ aU_{1,b}= a(R_1R_b + R_bR_1 - R_{1b})=ab \in z(J). \]  
Therefore, $ab \in z(J)$.  
\end{proof}

\begin{lemma}  
$z(J) \subseteq \operatorname{Nil} J$.  
\end{lemma}  

\begin{proof} 
Let $a \in z(J)$. We show that $a$ is nilpotent. 

By definition, $a = a_1 + \dots + a_n$, with each $a_i$ being an absolute zero divisor.  

If $n=1$, then $a$ is an absolute zero divisor. We have $a^3 = a U_a = 0$, so $a$ is nilpotent.  

We proceed by induction on $n$. Consider $a = b+z$, where $b$ is a sum of $n-1$ absolute zero divisors (and is nilpotent by the inductive hypothesis), and $z$ is an absolute zero divisor.  

By the Shirshov-Cohn Theorem, the subalgebra $J_0 = \operatorname{alg}\langle b, z \rangle$ is special. Thus, there exists an associative algebra $A$ such that $J_0 \subseteq A^{(+)}$. In $A$, we have the relations $zxz = 0$ for all $x \in J_0$ (since $U_z=0$ in $J_0$), which implies $z^3= 0$ and $zb^iz=0$ for all $i \ge 0$.  

Let $m$ be an integer such that $b^m=0$. Then we have  
\[ a^{2m+1}=(b+z)^{2m+1} = \sum_{i+j=2m} b^i z b^j + \sum_{i+j=2m-1}b^i z^2 b^j=0 \]  
\end{proof}

\begin{lemma}  
Let $J$ be a Jordan algebra. If $I$ is an ideal of $J$, then $I^3$ is an ideal in $J$.  
\end{lemma} 

\begin{proof} 
$ a,b,c \in I $, $ x \in J $. $ ((ab)c)x=(ab)(cx)+(ab,c,x) $ the frst term is already in $ I^3 $. $ (ab,c,x)=-(ax,c,b)-(bx,c,a) \in I^3$.     
\end{proof}

\begin{exercise}
Give an example of Jordan algebra $ J $, $ I \trianglelefteq J $ where $ I^3 $ is not an ideal in $ J $.    
\end{exercise}

\begin{theorem} 
If $ J $ is finite dimensional Jordan algebra. Then $ \operatorname{Nil}J=0 $ iff $ z(J)=0 $.   
\end{theorem}

\begin{proof}
$ \operatorname{Nil}J=0 $ then $ z(J)=0 $ since $ z(J) \subseteq \operatorname{Nil}J $. Conversely If $ z(J)=0 $, suppose $ N=\operatorname{Nil}J \neq 0$. $ N^{\< 1\>} \supseteq N^{\< 2\>} \supseteq \ldots \supseteq N^{\< i\>}=0$ where $ N^{\< 1\>}=N $ and $ N^{\< i+1\>}=(N^{\< i\>})^3 $. Since $ N^{\< k \>}  \trianglelefteq J$ we have $ I \triangleleft J $ with $ I \neq 0 $ and $ I^3=0 $. Let $ 0 \neq a \in I^2 $, then $ a^2 \in I^2 I^2 \subseteq I^3 =0 $. $ \forall x \in J $, $ xU_a=2(xa)a-a^2x=0 $ (note: $ (xa)a \in I I^2\subseteq I^3=0 $ ), $ a $ is an absolute zero divisor. $ a \in z(J) $. So $ a=0 $ and $ I^2=0 $. But $ \forall a \in I $, $ xU_a =2(xa)a-a^2x  =0 $ shows $ I=0 $. Controdiction.    
\end{proof}

\begin{remark}  
Observe that $J$ is always non-degenerate if $\operatorname{Nil} J = 0$. The condition that $J$ is finite-dimensional was used to prove that the nilradical is nilpotent, which in turn was used to find a non-trivial ideal whose cube is zero. There exist special, non-degenerate, infinite-dimensional Jordan algebras which contain non-trivial nil ideals (that are not nilpotent).  
\end{remark} 

\begin{corollary}
$ J $ finite dimensional Jordan algebra with $ \operatorname{Nil}J=0 $. Then for any ideal $ I \triangleleft J $ we have $ \operatorname{Nil}I=0 $.    
\end{corollary}

\begin{proof} 
Suppose $ \operatorname{Nil}I\neq 0 $. By theorem before, $ z(I)\neq 0 $.$\exists a \in z(I), a \neq 0 $, $ IU_a=0 $. From the other side $ z(J)=0 $, $ JU_a \neq 0 $. Choose $ x \in J $ and $ b = xU_a \neq 0 $. $ JU_b = JU_{xU_a}=JU_aU_xU_a \in IU_a =0$. Hence $ U_b =0  $ in $ J $ and $ b=0 $. Controdiction.           
\end{proof}

\begin{exercise}
J finite dimensional Jordan algebra, $ I \triangleleft J $ then $ \operatorname{Nil}I =I \cap \operatorname{Nil}J $.  
\end{exercise}

\begin{proof} 
($\supseteq$) We have $I \cap \operatorname{Nil} J \subseteq \operatorname{Nil} I$ since $I \cap \operatorname{Nil} J$ is a nil ideal of $I$.  

($\subseteq$) For the other inclusion, consider $\bar{J} = J/\operatorname{Nil} J$. Then we have $\operatorname{Nil} \bar{J} = 0$.  
Let $\bar{I} = (I + \operatorname{Nil} J) / \operatorname{Nil} J$, which is an ideal of $\bar{J}$.  
By the previous corollary, since $\operatorname{Nil} \bar{J} = 0$, we must have $\operatorname{Nil} \bar{I} = 0$.  

The canonical projection of $\operatorname{Nil} I$ into $\bar{J}$ is $(\operatorname{Nil} I + \operatorname{Nil} J) / \operatorname{Nil} J$. This is a nil ideal of $\bar{I}$. Since $\operatorname{Nil} \bar{I} = 0$, we must have  
\[ (\operatorname{Nil} I + \operatorname{Nil} J) / \operatorname{Nil} J = \{0\} \]  
This implies $\operatorname{Nil} I \subseteq \operatorname{Nil} J$.  
Since by definition $\operatorname{Nil} I \subseteq I$, we can conclude that $\operatorname{Nil} I \subseteq I \cap \operatorname{Nil} J$.  
\end{proof}

\section{Lecture 6}

\begin{definition} 
Let $ A $ be an algebra,
\begin{itemize}
\item associative center of $ A $: $ N(A) =\{ a \in A | (a,A,A)=(A,a,A)=(A,A,a)=0  \}$.
\item commutative center of $ A $: $ K(A)=\{ a \in A | [a,A]=0  \}$.
\item center of $ A $: $ Z(A)=N(A) \cap K(A) $.
\end{itemize}
\end{definition}

\begin{lemma} 
Let $ J $ be a Jordan algebra, let $ I \triangleleft J $ and $ e \in I $ be identity element in $ I $. Then $ e \in Z(J) $ and $ J =I \oplus K $ for some ideal $ K \triangleleft J $.
\end{lemma}

\begin{proof} 
$ e \in Z(J) $: $ x \in I, a \in J, e \in I $. $ (x,e,a)=(x,e^2,a) $. Consider the linearization of Jordan $ (b,cd,a)=c(b,d,a)+(b,c,a)d $ we have $ (x,e^2,a)=2e(x,e,a)=2(x,e,a)$. Hence $ (x,e,a)=0 $. Let $ a,b \in J $, $ (e,a,b)=(e^2,a,b) $, consider another linearization $ (ac,b,d)+(ad,b,c)+(cd,b,a)=0 $ we have $ (e,a,b)=(e^2,a,b)=-(eb,a,e)-(eb,a,e)=-2(eb,a,e) =-2(x,a,e)$ where $ x=eb $. $ (x,a,e)=xa-x(ae)=(xe)a-x(ae)=(x,e,a)=0  $ hence $ (e,J,J)=0 $. For every commutative or anticommutative algebra $ (a,b,a)=0 $. linearizing: $ a \mapsto a+c $. $ (a,b,c)=-(c,b,a) $ hence $ (J,J,e)=0 $. Finally in commutative algebra we have $ (a,b,c)+(b,c,a)+(c,a,b)=0 $. Therefore $ (J,e,J)=0 $ 
\begin{exercise}
Prove: in commutative algebra we have $ (a,b,c)+(b,c,a)+(c,a,b)=0 $.
\end{exercise}

\begin{proof} 
$ (a,b,c)+(b,c,a)+(c,a,b)=(ab)c-a(bc)+(bc)a-b(ca)+(ca)b -c(ab)$. Since commutivity, $ (ab)c=c(ab) $, $ a(bc)=(bc)a $, $ b(ca)=(ca)b $, $ (a,b,c)+(b,c,a)+(c,a,b)=0 $.     
\end{proof}

Consider $ K=\{x-xe | x\in J  \} $ we claim $ K \triangleleft J $. For $ x,y \in J $, $ (x-xe)y=xy-(xe)y=xy-x(ey)=xy-(xy)e \in K $. Since $ a=ae+(a-ae) \in I + K$, $ I \cap K=0 $: $ \kappa=ae=b-be $ then $ \kappa =ae=(ae)e =be-(be)e=be-b(ee)=be-be=0$.       
\end{proof}

\begin{definition} 
Let $ J $ be a Jordan algebra. We call $ J $ nil-semisimple if $ \operatorname{Nil}J=0 $.   
\end{definition}

\begin{theorem} 
Suppose that every simple Jordan algebra contains unity element. Then each fin-dim Jordan algebra $ J $ with $ \operatorname{Nil}J=0 $ has unity element and can be writen as a direct sum of simple algebras.   
\end{theorem}

\begin{proof} 
Induction on $ \operatorname{dim}J $. If $ J $ is simple nothing to prove. Let $ I \triangleleft J $ non-trivial ideal. It follows $ \operatorname{Nil}I=0 $. $ \operatorname{dim}I < \operatorname{dim}J $. $ I=I_1 \oplus \ldots \oplus I_k $ where $ I_i $ are simple. And $ I $ has unity element. By last lemma there exists $ K $, $ J=I\oplus K=I_1 \oplus \ldots \oplus I_k \oplus K_1 \oplus \ldots \oplus K_l $ with all summands simple.   
\end{proof}

Next we will show that every simple Jordan algebra contains unity element.

\begin{definition} 
An idempotent element $ e $ is an element such that $ e^2=e, e \neq 0 $. For a set of idempotent element $ \{e_i\} $, $ e_i $ is called orthogonal idempotent if $ e_ie_j=0 $ for all $ i \neq j $.
\end{definition}

\begin{definition} 
Let $ J $ be a Jordan algebra and $ e \in J $ an idempotent element. Consider Jordan identity:
\[
R_cR_bR_a+R_aR_bR_c+R_{(ac)b}=R_cR_{ba}+R_aR_{bc}+R_bR_{ac}
\]    
Replace $ a,b,c $ by $ e $ we have:
\[
R_e+2R_e^3-3R_e^2=R_e(2R_e-1)(R_e-1)=0
\]  
Then consider $ e $ acts on $ J $ we have decomposation of $ J $ as vector space
 \[ J=J_0(e) \oplus J_{\frac{1}{2}}(e) \oplus J_1(e) \] 
 Where $ J_i(e)=\{ a \in J | ea=ia , i=0,\frac{1}{2}, 1\} $. Such decomposition is called a Pierce decomposition with respect to idempotent $ e $.
\end{definition}

\begin{theorem}[Pierce decomposation in Jordan algebra] 
Let $ J $ be Jordan algebra, e idempostent. $ J=J_0 \oplus J_{\frac{1}{2}} \oplus J_1 $ where $ J_i=\{ a \in J | ea=ia , i=0,\frac{1}{2}, 1\} $. We have the following: $ J_1^2 \subseteq J_1 $, $ J_0^2 \subseteq J_0 $, $ J_i J_{\frac{1}{2}} \subseteq J_{\frac{1}{2}},i=0,1 $, $ J_{\frac{1}{2}}^2 \subseteq J_0\oplus J_1 $, $ J_1J_0=0 $.      
\end{theorem}

\begin{proof} 
We have decomposation $ J=J_0 \oplus J_{\frac{1}{2}} \oplus J_1  $, take $ a \in J_0 \cup J_1 $, $ x \in J $. $ (e,x,a)=(e^2,x,a)=-2(ea,x,e) $. If $ a \in J_0 $, $ -2(ea,x,e)=0 $ hence $ (e,x,a)=0 $. If $ a \in J_1 $, $ -2(ea,x,e)=-2(a,x,e)=2(e,x,a) $. Hence $ (e,x,a)=0 $.  

Take $ a_i, b_i \in J_i $, $ i=0,1 $. $ (a,J,e)=0, \forall a \in J_0 \cup J_1 $. $ (a_1b_1)e=a_1(b_1e)+(a_1,b_1,e)=a_1b_1 $, Hence $ J_1^2 \subseteq J_1 $. $ (a_0b_0)e=a_0(b_0e)+(a_0,b_0,e)=0 $, Hence $ J_0^2 \subseteq J_0 $, $ (a_1b_0)e=a_1(b_0e)+(a_1,b_0,e)=0 ,a_1b_0 \in J_0$, on the other hand $ (a_1b_0)e=(b_0a_1)e=b_0(a_1e)+(b_0,a,e)=b_0a_1=a_1b_0, a_1b_0 \in J_1 $. Therefore $ a_1b_0 \in J_0 \cap J_1 =0 $, i.e. $ J_1J_0=0 $.  

Let $ a \in J_0 \cup J_1 $ and $ b \in J_{\frac{1}{2}} $, $ (ab)e=a(be)+(a,b,e) =\frac{1}{2}ab$. Hence $ J_0J_{\frac{1}{2}} \subseteq J_{\frac{1}{2}} $ and $ J_1J_{\frac{1}{2}} \subseteq J_{\frac{1}{2}} $.

Let  $ a=a_0+a_{\frac{1}{2}}+a_1 $, $ (a,e,e)=(a_{\frac{1}{2}},e,e) =(a_{\frac{1}{2}}e)e-a_{\frac{1}{2}}e^2=-\frac{1}{4}a_{\frac{1}{2}} \neq 0$. But $ (a_{\frac{1}{2}} b_{\frac{1}{2}},e,e)=-(a_{\frac{1}{2}}e,e,b_{\frac{1}{2}}) -(b_{\frac{1}{2}}e,e,a_{\frac{1}{2}}) =0$. Hence $ J_{\frac{1}{2}} ^2 \cap J_{\frac{1}{2}} =0 $ and $ J_{\frac{1}{2}} ^2 \subseteq J_0 \oplus J_1 $.  
\end{proof}

\begin{proposition}
Let $ A $ be power associative algebra, finite dimensional, not nil. Then such algebra contains idempotent.  
\end{proposition}

\begin{proof} 
Let $ a \in A $ non-nilpotent. take $ \{a,a^2, \ldots\} $. Since $ A $ is finite dimensional, this set is linear dependent. we have $ a^k+\alpha_{k+1}a^{k+1}+ \ldots + \alpha_{k+m}a^{k+m} =0$. since $ a^i \neq 0 $  at least one coefficent not zero. Rewrite it as $ a^k=a^k(-\alpha_{k+1}a- \ldots -\alpha_{k+m}a^m) =a^kg(a)$. Hence $ a^k=a^k(g(a))^k =a^k g_1(a) a^k$. Then $ e=a^kg_1(a) $ is non-zero and idempotent ($ (a^kg_1(a))^2=a^kg_1(a)a^kg_1(a)=a^kg_1(a) $).        
\end{proof}

\section{Lecture 7}

\begin{lemma}  
Let $J$ be a Jordan algebra and let $e$ be an idempotent in $J$. Then $U_e$ is a projection onto $J_1$ with kernel $J_{1/2} \oplus J_0$, where $J = J_1 \oplus J_{1/2} \oplus J_0$ is the Peirce decomposition of $J$ with respect to $e$.  
\end{lemma}  
\begin{proof}  
First, we show that the image of $U_e$ is $J_1$. Let $a \in \mathrm{Im}(U_e)$, so $a = bU_e$ for some $b \in J$. By definition, $bU_e = 2(eb)e - e^2b = 2(be)e - be$.  
We check if $ae = a$:  
\begin{align*}  
    ae &= (2(be)e - be)e = 2((be)e)e - (be)e  
\end{align*}  
$ 2((be)e)e=2(be)(ee)+2(be,e,e) $ where $ (be,e,e)=-(e,e,b)-(be,e,e) $ hence $ (be,e,e)=-\frac{1}{2}(e,e,b) $. We have $2(be)(ee)+2(be,e,e)= 2(be)e-(e,e,b)=2(be)e-eb+e(eb)=3(be)e - be $. Thus,    
\begin{align*}  
    ae &= (3(be)e - be) - (be)e \\
    &= 2(be)e - be \\
    &= a  
\end{align*}  
Since $ae=a$, $a$ is in $J_1$. Thus, $\mathrm{Im}(U_e) \subseteq J_1$.  

Conversely, let $a \in J_1$. By definition, $ae = a$. We apply $U_e$:  
\[ aU_e = 2(ea)e - e^2a = 2ae - ea = 2a - a = a. \]  
This shows that $a \in \mathrm{Im}(U_e)$. Thus, $J_1 \subseteq \mathrm{Im}(U_e)$.  
Combining both inclusions, we have $\mathrm{Im}(U_e) = J_1$.  

Moreover, for any $a \in J$, let $b=aU_e$. We have shown $b \in J_1$. Applying $U_e$ again, $bU_e = b$. So, $aU_eU_e = aU_e$, which means $U_e^2 = U_e$. Thus, $U_e$ is a projection operator.  

Next, we find the kernel of $U_e$. Let $t \in J_0$. By definition, $te=0$. Then $tU_e = 2(te)e - e^2t = 0$. So, $J_0 \subseteq \ker(U_e)$.  

Let $t \in J_{1/2}$. By definition, $te = \frac{1}{2}t$. Then  
\[ tU_e = 2(te)e - e^2t = 0. \]  
So, $J_{1/2} \subseteq \ker(U_e)$.  

Therefore, $J_0 \oplus J_{1/2} \subseteq \ker(U_e)$. Since the image is $J_1$ and $J=J_1 \oplus J_{1/2} \oplus J_0$, the kernel must be exactly $J_0 \oplus J_{1/2}$.  
\end{proof}  

\begin{lemma}  
Let $J$ be a finite dimensional Jordan algebra. If $\operatorname{Nil} J = 0$, then $\operatorname{Nil} J_1 = 0$ for any idempotent $e \in J$.  
\end{lemma}  

\begin{proof}  
If $\operatorname{Nil}(J)=0$, then $J$ is non-degenerate. We will show that $J_1$ is non-degenerate as well.  
Let $a \in J_1$ be an absolute zero divisor, i.e., $J_1U_a=0$.  
Since $a \in J_1$, we have $aU_e=a$ by the previous lemma. Then, using the fundamental identity $U_{yU_x} = U_x U_y U_x$:  
\[ U_a = U_{aU_e} = U_e U_a U_e \]  
Therefore, $ \forall x \in J $   
\[
    xU_a = xU_e U_a U_e\in J_1 U_a U_e =0 
 \]
This implies $a$ is an absolute zero divisor in $J$. Since $J$ is non-degenerate, $a=0$.  
\end{proof} 

\begin{definition} 
An idempotent $e \in J$ is called a principal idempotent if there are no idempotents $f$ such that $ef=0$.  
\end{definition}  

\begin{remark}  
If $J$ is a finite-dimensional Jordan algebra, then an idempotent $e \in J$ is a principal idempotent if and only if $J_0(e)$ is nil.  
(This is because an idempotent $f$ is orthogonal to $e$ if and only if $f \in J_0(e)$, and if $ J_0(e) $ have no idempotent then it must be nil).  
\end{remark}  

\begin{lemma}  
Let $J$ be a finite-dimensional, non-degenerate Jordan algebra. If $e$ is a principal idempotent in $J$, then $e$ is the unity of $J$.  
\end{lemma}  

\begin{proof} 
In this case $J_0(e)$ is a nil subalgebra. Since $J$ is non-degenerate, it does not have non-zero nil ideals. A nil subalgebra of a non-degenerate Jordan algebra must be zero, so $J_0(e)=0$. This implies that the Peirce decomposition of $J$ is $J = J_1 \oplus J_{1/2}$. We need to show that $J_{1/2}=0$.  

\begin{tcolorbox}[colback=gray!20, colframe=black] 
Note: $ J_0 $ is just a subalgebra but not an ideal, why we can say: ``A nil subalgebra of a non-degenerate Jordan algebra must be zero''? In fact we need a discuss of ``trace form'' (if you know something about Lie algebra this is similarly to define a Killing form in Lie algebra). This is natural: for a Jordan algebra $ J $ and its multiplication algebra $ \operatorname{Mult}J $, we know that all multiplications $ R_x, x\in J $ is a linear map of the vector space $ J $. Hence we define a symmetric bilinear form $ (x,y)=\operatorname{tr}(R_xR_y) , x,y \in J$ on $ J $. We call it the trace form of $ J $. By the property of trace and $ R_{(x,y,z)}=[R_z,[R_x,R_y]] $, $ (x,yz)=(xy,z), \forall x,y,z \in J $. And based on this property the radical of this form is an ideal of $ J $. Furthermore we can prove when $ \operatorname{char}=0 $, it is the nil radical $ \operatorname{Nil}J $. And $ \forall x \in J_{\frac{1}{2}}, \operatorname{tr}R_x=0 $. Hence $ \forall x,y \in J$, write $ x=x_1+x_{\frac{1}{2}}+x_0 $ and $ y=y_1+y_{\frac{1}{2}}+y_0 $, $(x,y)=\operatorname{tr}R_{x_1y_1+z_0} =\operatorname{tr}R_{x_1y_1}+\operatorname{tr}R_{z_0}$ where $ z_0 $ is an element in $ J_0 $. Now $ J_0 $ is nil. $\operatorname{tr}R_{z_0}=0, \forall z_0 \in J_0  $. If $ x \in J_0 $, then $ \forall y \in J $, $ (x,y)= \operatorname{tr}R_{x_1y_1}=0$ (since $ x_1=0 $ ). Hence $ x $ is in the radical of this trace form. Thus $ J_0 $ is in fact in the nil radical of $ J $. Since $ J $ is non-degenerated, $ J_0=0 $.                           
\end{tcolorbox}

Let $t \in J_{1/2}$. Let us prove that $ t^2 $ is absolute zero divisor in $ J $. Since $U_e$ is a projection on $J_1$, we have: $tU_e = 0$. And $J_1 U_t = J_1 U_t U_e = J U_e U_t U_e = JU_{tU_e}=0$. In which $ J_1 U_t = J_1 U_t U_e $ since $ J_1 U_t \in J_1 $. In particular, $U_{t^2} = U_t U_t$ (Proposition 5.15). This gives $J_1 U_t^2 = (J_1 U_t) U_t = 0$. 

Observe that $t^2 U_e = t^2$ (since $t^2 \in J_1 \oplus J_0$ and $J_0=0$, thus $t^2 \in J_1$). $J_{1/2} U_t^2 = J_{1/2} U_t U_t = J_{1/2} U_e U_t^2 U_e = 0$. Thus we conclude that $t^2$ is an absolute zero divisor of $J$. Since $J$ is non-degenerate, $t^2=0$. Then for each $t_1, t_2 \in J_{1/2}$, $2t_1t_2 = (t_1+t_2)^2 - t_1^2 - t_2^2 = 0.$ 

Then $J_{1/2}^2 = 0$. We have that $J_{1/2}$ is a nilpotent ideal in $J$. Since $J$ is non-degenerate, its nilradical is $0$, so $J_{1/2}=0$. Consequently, $J=J_1$ and $e$ is the unity element in $J$. 
\end{proof}

\begin{corollary}  
Each finite-dimensional Jordan algebra with nilradical equal to zero has a unity element.  
\end{corollary}  

\begin{proof}  
It is enough to show that $J$ has a principal idempotent. We know from proposition 6.9 that a non-nil algebra contains an idempotent $e$.  
If $e$ is not a principal idempotent, then $J_0(e)$ is not nil and there exists an idempotent $f \in J_0(e)$. Let $e_1 = e+f$. We have $e_1^2 = e_1$ and $J_1(e+f) \supsetneq J_1(e)$.  
If $J_0(e_1)$ is not nil, we can continue the process. Since $\dim J < \infty$, we will eventually find an idempotent $\tilde{e} \in J$ such that $J_0(\tilde{e})$ is nil, which means $\tilde{e}$ is principal. By the previous lemma, $\tilde{e}$ is the unity of $J$.  
\end{proof}  

\begin{corollary}  
Each finite-dimensional Jordan algebra with $\mathrm{Nil}(J)=0$ has a unity and is isomorphic to a direct sum of simple algebras.  
\end{corollary}  

\begin{proof}  
Each simple algebra has $\mathrm{Nil}(J)=0$ (otherwise it is nilpotent). And Theorem 6.5 follows.  
\end{proof}  

\section{Lecture 8}

We look at how many pairwise orthogonal idempotents $J$ contains.  

\begin{definition}  
\begin{itemize}  
    \item An idempotent is called primitive if it cannot be written as $e = e_1 + e_2$ for non-trivial orthogonal idempotents $e_i$.  
    \item An idempotent $e \in J$ is called absolutely primitive if for each $a \in J_1(e)$, we have $a = \alpha e + n$, for some $\alpha \in F$ and nilpotent $n$.  
\end{itemize}  
\end{definition}  

\begin{lemma}  
Every absolutely primitive idempotent is primitive.  
\end{lemma}  

\begin{proof}  
Let $J$ be a Jordan algebra and $e \in J$ an absolutely primitive idempotent. Observe that if $e = e_1 + e_2$ with $e_1 e_2 = 0$, then $e_1, e_2 \in J_1(e)$, since $e_1 e = e_1$ and $e_2 e = e_2$ (note: $ (e-e_1)e_1=e_2e_1=0 $ hence $ e_1=ee_1 $, similarly for $ e_2 $). Now, if $f \in J_1(e)$ is an idempotent, we write $f = \alpha e + n$ where $0 \neq \alpha \in F$ and $n$ is nilpotent. We have $ \alpha e + n = f = f^2 = \alpha^2 e + 2\alpha n + n^2 $ which implies $ (\alpha - \alpha^2)e = (2\alpha - 1)n + n^2 $ Then $\alpha^2=\alpha$, which implies $\alpha=0$ or $\alpha=1$. Given $\alpha \neq 0$. And $n^2=(1-2\alpha)n$. Since $\alpha \neq 0$, we obtain $\alpha=1$. Then $n^2=-n \implies n=0$ and $f=e$. $f$ is the unique idempotent in $J_1(e)$. $\implies e$ is primitive.   
\end{proof} 

\begin{lemma}  
If $F$ is algebraically closed, then any primitive idempotent in a finite-dimensional Jordan algebra is absolutely primitive.  
\end{lemma}  

\begin{proof}  
Let $e \in J$ be a primitive idempotent and let $a \in J_1(e)$. Since $J$ is finite-dimensional, there exists a minimal polynomial $m_a \in F[x]$ such that $m_a(a)=0$. We will prove that $m_a$ is a power of an irreducible polynomial. In fact, if $m_a = f_1 f_2$ with $\gcd(f_1, f_2)=1$, then by Bezout's theorem there exist $h_1, h_2 \in F[x]$ such that 
\[ f_1 h_1 + f_2 h_2 = 1 \quad (*) \] 
Besides, we may choose such polynomials that $\deg h_1 < \deg f_2$ and $\deg h_2 < \deg f_1$ (by Euclid's Algorithm). Let $\phi_a : F[x] \to \mathrm{alg}\langle e, a \rangle \subseteq J_1(e)$ be an algebra homomorphism with $\phi_a(x) = a$ and $\phi_a(1) = e$. Applying $\phi_a$ to (*) we obtain 
\[ f_1(a)h_1(a) + f_2(a)h_2(a) = e \quad (**) \] 
Let $e_1 = f_1(a)h_1(a)$ and $e_2 = f_2(a)h_2(a)$. Multiplying (**) by $e_i$ we obtain $e_i^2 = e_i$ (since $e$ is the unity element in $J_1(e)$ and $f_1 f_2$ annihilates $a$). Also we have $e_1 e_2 = 0$.  

Moreover, since $\deg f_i, \deg h_i < \deg m_a$, both elements are non-zero, and then $e_1$ and $e_2$ are orthogonal idempotents, which is a contradiction. Then, $m_a = p^k$, where $p \in F[x]$ is irreducible. Since $F$ is algebraically closed, $p = x-\alpha$ for some $\alpha \in F$, and $0=m_a(a)=(a-\alpha e)^k$. Then if we let $n = a-\alpha e$, we have $a=\alpha e + n$, where $n$ is nilpotent of index at most $k$.  
\end{proof} 

\begin{lemma}  
Let $J$ be a finite-dimensional Jordan algebra with a unity element $1$. Then we may write $1 = e_1 + \dots + e_m$, where the $e_i$ are pairwise orthogonal primitive idempotents.  
\end{lemma}  

\begin{proof}  
If $1$ is primitive, there is nothing to prove. If not, we can write $1 = e_1 + e_2$, where $e_1, e_2$ are orthogonal idempotents. If $e_1$ is not primitive, we can write $e_1 = e_1' + e_1''$ with $e_1', e_1''$ being primitive idempotents. We have $e_1', e_1'' \in J_1(e_1)$ and $e_2 \in J_0(e_1)$, so they are orthogonal. Then $e_1', e_1'', e_2$ are orthogonal. We can continue this process, which must end, since the number of pairwise orthogonal idempotents is bounded by $\dim J$.  
\end{proof}  

From now on, we consider algebraically closed fields.  

\begin{definition}  
Let $J$ be a finite-dimensional Jordan algebra with $\mathrm{Nil}(J)=0$. The maximal number of primitive idempotents in a decomposition of the unity element is called the capacity of $J$.  
\end{definition}  

In the conditions of last definition, we write $1=e_1+\dots+e_n$. If $F$ is algebraically closed, we have $J_1(e_i) = Fe_i + N_i$, where $N_i$ is a set of nilpotent elements and $\mathrm{Nil}\,J_1(e_i) = 0$, for all $i=1,\dots,n$.  

\begin{definition}[Full Peirce decomposition] 
Let $J$ be a Jordan algebra with unity element.  
Let $e_1, \dots, e_n$ be non-zero idempotents, pairwise orthogonal with $1=e_1+\dots+e_n$.  
Then  
\[ J = \bigoplus_{i=1}^n J_{ii} \oplus \bigoplus_{1 \le i < j \le n} J_{ij} \]  
\[ J_{ii} = J_1(e_i) \text{ and } J_{ij} = J_{1/2}(e_i) \cap J_{1/2}(e_j) \text{ for } 1 \le i < j \le n. \]  
Moreover,  
\begin{enumerate}  
    \item[i] $J_{ii}^2 \subseteq J_{ii}$, $J_{ij}^2 \subseteq J_{ii} + J_{jj}$. $J_{ii}J_{ij} \subseteq J_{ij}$, $J_{ii}J_{jj} = 0$ for $i \ne j$.  
    \item[ii] $J_{ij}J_{jk} \subseteq J_{ik}$, $J_{ij}J_{kk} =0$ for $i,j,k$ different. 
    \item[iii] $J_{ij}J_{kl} = J_{ii}J_{kl} = 0$ for $i,j,k,l$ different.  
\end{enumerate}  
\end{definition}

\begin{proof} 
Let $i \ne j$ and put $e = e_i + e_j$, $B = JU_e$. Then $B = J_{ii} + J_{ij} + J_{jj}$ is the Peirce decomposition of $B$ relative to $e_i, e_j$. We have $J_{ii} = B_1(e_i)$, $J_{jj} = B_0(e_i) = J_{jj} = B_1(e_j)$, $J_{ij} = B_{1/2}(e_i) = J_{ji} = B_{1/2}(e_j)$. Hence (i) follows from theorem 6.8. 

Next let $i,j,k$ be distinct and put $f = e_i + e_j + e_k$, $C = JU_f$. Then $C = J_{ii} + J_{jj} + J_{kk} + J_{ij} + J_{ik} + J_{jk}$ is the Peirce decomposition of $C$ relative to $e_i, e_j, e_k$ ($J_{pq} = C_{pq}$). Also let $e = e_i + e_j$. Then $e$ is an idempotent contained in $C$ and $C_1(e) = J_{ii} + J_{ij} + J_{jj}$, $C_0(e) = J_{kk}$ and $C_{1/2}(e) = J_{ik} + J_{jk}$. Since $C_0(e)C_1(e) = 0$ we have $J_{jj}J_{kk} = 0$ and $ J_{ij}J_{kk}=0 $. Since $C_{1/2}(e)C_1(e) \subseteq C_{1/2}(e)$ we have $J_{ij}J_{jk} \subseteq J_{ik} + J_{jk}$ and since $J_{ij}=J_{ji}$, $J_{jk}=J_{kj}$ we have also $J_{ij}J_{jk} = J_{kj}J_{ji} \subseteq J_{ki} + J_{ji}$. Hence $J_{ij}J_{jk} \subseteq (J_{ik} + J_{jk}) \cap (J_{ki} + J_{ji}) = J_{ik}$. 

It remains to prove (iii). For this we take $g = e_i + e_j + e_k + e_l$ and $D = JU_g$. Then $D = J_{ii} + J_{jj} + J_{kk} + J_{ll} + J_{ij} + J_{ik} + J_{il} + J_{jk} + J_{jl} + J_{kl}$. Also $e = e_i+e_j$ is an idempotent contained in $D$ and $J_{ij} \subseteq D_1(e)$ while $J_{kl} \subseteq D_0(e)$. Hence $D_0(e)D_1(e)=0$ gives $J_{ij}J_{kl}=0$.  
\end{proof}

\begin{remark}
Why it is enough to consider algebraically closed fields? Let $A$ be a simple algebra with unity over $F$, central over $F$ ($Z(A)=F$). Then for each field $K \supset F$ we have that the algebra $A_K = K \otimes_F A$ is simple as well and is central over $K$ with $\dim_K A_K = \dim_F A$.  
\end{remark}

\section{Lecture 9}

\begin{definition}  
Let $J$ be a Jordan algebra with unity element. An element $a \in J$ is called invertible if there exists $b \in J$ such that $ab=1$ and $a^2b=a$. In this case, $b$ is called the inverse of $a$ and is denoted by $a^{-1}$.  
\end{definition} 

\begin{lemma}  
Let $J$ be a Jordan algebra with unity element, then  
\begin{enumerate} 
    \item If $b = a^{-1}$ then $a = b^{-1}$.  
    \item $a$ is invertible $\iff$ $U_a$ is invertible. Moreover, $U_a^{-1} = U_{a^{-1}}$.  
    \item $a, b$ are invertible $\iff \{aba\}$ is invertible.  
\end{enumerate}  
\end{lemma} 

\begin{proof}  
\noindent(i) Suppose $b$ is the inverse of $a$, then  
\[ a^2b^2 = (a^2b)b - (a^2,b,b) = ab + 2(ab,b,a) = 1 + 2(1,b,a) = 1 \]  
Then,  
\[ b^2a = b^2(a^2b) = (b^2a^2)b = b. \]  
Thus $a$ is the inverse of $b$. \\
(ii) Let $a$ be invertible with inverse $b$. Then:  
\begin{align*}  
      bU_a &= 2(ba)a - ba^2 = 2a - a = a \\
      b^2U_a &= 2(b^2a)a - b^2a^2 = 2ba - 1 = 1 
\end{align*}  
Observe that $U_1 = \mathrm{Id}_{J}$, the identity operator in $J$. Then  
\[ \mathrm{Id}_{J} = U_1 = U_{b^2U_a(b^2)} = U_a U_{b^2} U_a. \]  
This means that $U_a$ is invertible.  
In this case we have  
\[ U_a U_b = U_{b U_a} U_b = U_a U_b U_a U_b \]  
Since $U_a, U_b$ are invertible we conclude that $U_a U_{a^{-1}} = \mathrm{Id}_{J}$. Analogously, $U_{a^{-1}} U_a = \mathrm{Id}_{J}$.  
Consequently, $(U_a)^{-1} = U_{a^{-1}}$.   

Conversely, suppose that $U_a$ is invertible and let $b = a U_a^{-1}$. Then $b U_a = a$.  
Since $[U_a, R_a]=0$ we also have $[U_a^{-1}, R_a]=0$. ($ [U_a^{-1}, R_a]=U_a^{-1}[R_a, U_a] U_a^{-1}=0$). Then  
\begin{align*}  
    ba &= a U_a^{-1} R_a = a R_a U_a^{-1} = a^2 U_a^{-1} \\
       &= 1 U_a U_a^{-1} = 1  
\end{align*}  
and  
\[ a = b U_a = 2(ba)a - ba^2 = 2a - ba^2. \]  
Thus $ba^2=a \implies a$ is invertible.  

(iii) By (ii), $\{aba\}$ is invertible $\iff U_{\{aba\}}$ is invertible. We have $U_{\{aba\}} = U_a U_b U_a$. Then $\{aba\}$ is invertible $\iff U_a$ and $U_b$ are invertible. By (ii) it is equivalent to the fact that $a$ and $b$ are invertible.  
\end{proof} 

\begin{theorem}[Albert-Jacobson-McCrimmon]  
Let $J$ be a Jordan algebra. If the unit element of $J$ is an absolutely primitive idempotent (every element $a \in J$ has the form $a = \alpha 1 + n$ with nilpotent $n$), then the set $N = \{n \in J \mid n \text{ is nilpotent}\}$ is an ideal in $J$.  
\end{theorem}  

\begin{lemma}[McCrimmon]  
Let $z \in J$ be nilpotent. Then there exists a nilpotent $w \in J$ such that $1 - 4z = (1 - 2w)^2$.  
\end{lemma} 

\begin{proof}  
We show that there exists $w \in \mathrm{alg}(z) = \mathrm{vect}\{z, \dots, z^k\}$ where $z$ has nilpotency index $k+1$. Writing $w = \alpha_1 z + \alpha_2 z^2 + \dots + \alpha_k z^k$, with $\alpha_i \in F$. We want $(1 - 2\sum_{i=1}^k \alpha_i z^i)^2 = 1-4z$. Since  
\[ (1-2\sum_{i=1}^k \alpha_i z^i)^2 = 1 - 4\sum_{i=1}^k \alpha_i z^i + 4\sum_{i,j=1}^k \alpha_i \alpha_j z^{i+j} \]  
and $\{z, z^2, \dots, z^k\}$ is linearly independent, the coefficients $\alpha_i$ have to satisfy  
\[ \begin{cases}  
\alpha_1 = 1 \\
\alpha_2 - \alpha_1^2 = 0 \\
\vdots \\
\alpha_k - \sum_{\substack{i+j=k \\ i,j<k}} \alpha_i \alpha_j = 0  
\end{cases} \]  
Since it is a triangular system, it has a solution, thus such a $w$ exists.  
\end{proof} 

\begin{proof}[Proof of the Theorem]  
It is enough to show that $N$ is a subspace.  
Since $n \in N$ implies $n^2 \in N$ and then for $n_1, n_2 \in N$, $2n_1n_2 = (n_1+n_2)^2 - n_1^2 - n_2^2 \in N$.  
And thus for $n \in N$ and $a = \alpha 1 + n'$ with $\alpha \in F$ and $n' \in N$, we have $an = \alpha n + n'n \in N$.  

We have to show that $N$ is a subspace.  
We suppose that there exist $n_1, n_2 \in N$ such that $n_1+n_2 = \alpha 1 + n$ with $\alpha \ne 0$ and $n \in N$.  
We may suppose without loss of generality that $n_1+n_2 = 1-4n$ with nilpotent $n$.  
By the previous lemma, there exists $w \in N$ such that $(1-2w)^2 = 1-4n$. Moreover, $1-2w$ is invertible, since $2w$ is nilpotent.  
Then we have that $U_{1-2w}$ is invertible (by properties of invertible).  

Using $U_{1-2w}^{-1} = U_{(1-2w)^{-1}}$ we obtain  
\[ 1 \cdot U_{1-2w} = (1-2w)^2 = 1-4n = n_1+n_2 \]  
Put $t=(1-2w)^{-1}$ we get  
\[ 1 = n_1 U_t + n_2 U_t \]  
Since $m_1 = n_1 U_t$ and $m_2 = n_2 U_t$ are nilpotent we obtain that $1-m_1 = m_2$ is invertible and nilpotent, contradiction.  
Thus $n_1+n_2$ are nilpotent for any $n_1, n_2 \in N$. 
\end{proof}  

\section{Lecture 10}

Let $J$ be a finite-dimensional simple Jordan algebra. By decomposition into primitive idempotents $1=e_1+\dots+e_n$, $e_i$ orthogonal and primitive. 

If $F$ is algebraically closed, then these idempotents are absolutely primitive, i.e., $J_1(e_i) = Fe_i + N_i$ where $N_i$ is a nilpotent ideal in $J_1(e_i)$ (by Albert-Jacob-McCrimmon Thm). Then by lemma from previous lecture (Let J be a finite dimensional Jordan algebra, if Nil$J=0 \implies$ Nil$J_1=0$). $N_i = \mathrm{Nil}\,J_1(e_i)=0$, thus $J_1(e_i) = Fe_i$. Consequently the Peirce decomposition gives  
\[ J = \bigoplus_{i=1}^n Fe_i \oplus \bigoplus_{i<j} J_{ij}  \]  
When a Jordan algebra has the Peirce decomposition in this form we call such an algebra reduced.  

\begin{corollary}  
Let $J$ be a finite dimensional simple reduced Jordan algebra. If $J$ has capacity 1, then $J \cong F$.  
\end{corollary}  

\begin{proposition}  
Let $J$ be a finite dimensional  simple reduced Jordan algebra. If the capacity of $J$ is 2, then $J$ is isomorphic to a Jordan algebra of Clifford type $J(V,f)$.  
\end{proposition} 

\begin{proof}  
We have $J = Fe_1 \oplus J_{12} \oplus Fe_2$, with $1 = e_1 + e_2$. Let $V = J_{12} + F(e_1-e_2)$. For $u,v \in V$, we want to show that $uv = \alpha(e_1+e_2) = \alpha 1$ for some $\alpha \in F$. (Actually, it is enough to show this for $u^2$, since $2uv = (u+v)^2 - u^2 - v^2 = (\alpha_1 - \alpha_2 - \alpha_3)1$ where $(u+v)^2 = \alpha_1 1$, $v^2 = \alpha_2 1$, and $u^2 = \alpha_3 1$).  

Observe that $(e_1-e_2)^2 = e_1+e_2=1$ and $u(e_1-e_2) = 0$ for $u \in J_{12}$. Moreover, $J_{12}^2 \subseteq Fe_1 + Fe_2$, and for $u \in J_{12}$, this means $u^2 = \alpha e_1 + \beta e_2$ for some $\alpha, \beta \in F$.  

We will show that $\alpha=\beta$. We get:  
\[ (\alpha e_1)u = (u^2 e_1)u = u^2(e_1 u) \quad \text{and} \quad (\beta e_2)u = (u^2 e_2)u = u^2(e_2 u) \]  
Since $u \in J_{12}$, we have $e_1 u = \frac{1}{2}u$ and $e_2 u = \frac{1}{2}u$. Thus,  
\[ \alpha(\frac{1}{2}u) = u^2(\frac{1}{2}u) \quad \text{and} \quad \beta(\frac{1}{2}u) = u^2(\frac{1}{2}u) \]  
This implies $\frac{1}{2}\alpha u = \frac{1}{2}\beta u$, so $\frac{1}{2}(\alpha-\beta)u = 0$. If $u \ne 0$, then $\alpha = \beta$. (If $u=0$, then $u^2=0$, so $\alpha=\beta=0$). Then we have shown that $uv \in F \cdot 1$ for all $u,v \in V$.  

Define $f: V \times V \to F$ by $f(u,v) = uv$ where $uv = \alpha \cdot 1$, for all $u,v \in V$. Then $f$ is a symmetric bilinear form. To show that $f$ is non-degenerate, take $x \in V$ such that $f(x,v)=0$ for all $v \in V$. Let $x = y + \gamma(e_1-e_2)$ with $y \in J_{12}$. $f(x, x) = \gamma^2 = 0$, so $x=y \in J_{12}$. Then $f(x, z) = 0$ for all $z \in J_{12}$, which means $xz=0$. So we have $x \in J_{12}$ and $x J_{12} = 0$. Also $x e_1 = \frac{1}{2}x$ and $x e_2 = \frac{1}{2}x$. This means $F x$ is an ideal in $J = Fe_1 \oplus J_{12} \oplus Fe_2$. Since $J$ is a simple Jordan algebra, this implies $x=0$. Thus $f(x,V)=0$ implies $x=0$, so $f$ is non-degenerate. And $ J=F1+V \simeq J(V,F) $.  
\end{proof}  

\begin{remark}  
Let $J$ be a reduced Jordan algebra. Using an argument of the Proposition, we may show that if $i \ne j$, $x \in J_{ij}$ and $x \cdot J_{ij}=0$, then $x=0$. Indeed, the condition $x \cdot J_{ij}=0$ implies that $Fx$ is an ideal in $Fe_i + J_{ij} + Fe_j$ and $(Fx)^2=0$. Then $Fx \subseteq \mathrm{Nil}(J_1(e_i+e_j))=0$, since $\mathrm{Nil}(J)=0$.  
\end{remark}  

We describe now the subspaces $J_{ij}$ for a simple Jordan algebra with capacity $\ge 2$, where $F=\bar{F}$.  

\begin{lemma}  
Let $J$ be a simple, finite-dimensional, reduced Jordan algebra, and $1=e_1+\dots+e_n$ a decomposition into absolutely primitive, pairwise orthogonal idempotents. Then $J_{ij} \ne 0$ for all distinct $i,j \in \{1, \dots, n\}$.  
\end{lemma} 

\begin{proof}  
First, observe that if $J_{ij} \ne 0$ then $J_{ij}^2 \ne 0$. Indeed, if $J_{ij}^2=0$, then $J_{ij}$ is a nil ideal of $Fe_i+J_{ij}+Fe_j$, so $J_{ij} \subseteq \mathrm{Nil}(Fe_i+J_{ij}+Fe_j) = 0$. Next, $Fe_i+J_{ij}+Fe_j$ is a Jordan algebra of Clifford type. Thus if $J_{ij} \ne 0$, there exists an $a_{ij} \in J_{ij}$ such that  
\[ a_{ij}^2 = \alpha(e_i+e_j) \quad \text{with } \alpha \ne 0. \]  

Let $i,j,k$ be distinct. Observe that  
\begin{equation*}  
    a_{ij}  b_{jk} = 2(a_{ij} b_{jk})a_{ij}  
\end{equation*}  
That is, for all $a_{ij} \in J_{ij}$ and $b_{jk} \in J_{jk}$.  
By the Jordan identity, we have
\begin{align*}  
    a_{ij}^2 b_{jk} &= -2(a_{ij}^2, e_k, b_{jk}) = 4(a_{ij}b_{jk}, e_k, a_{ij}) \\
    &= 2(a_{ij}b_{jk})a_{ij}, \quad \text{knowing that } a_{ij}b_{jk} \in J_{ik}.  
\end{align*}  

Next we will show that if $J_{ij} \ne 0$ and $J_{jk} \ne 0$ then $J_{ik} \ne 0$. 

By $ a_{ij}  b_{jk} = 2(a_{ij} b_{jk})a_{ij}  $ there exist $a_{ij} \in J_{ij}$ and $b_{jk} \in J_{jk}$ such that  
$a_{ij}^2 = \alpha(e_i+e_j)$ and $b_{jk}^2 = \beta(e_j+e_k)$.  

Then, by linearization of the Jordan identity,  
\begin{align*}  
    2(a_{ij}b_{jk})^2 &= -a_{ij}^2 b_{jk}^2 + (a_{ij}^2 b_{jk})b_{jk} + 2((a_{ij}b_{jk})b_{jk})a_{ij} \\
    &= -\alpha\beta(e_i+e_j)(e_j+e_k) + \alpha((e_i+e_j)b_{jk})b_{jk} + (a_{ij}b_{jk}^2)a_{ij} \\
    &= -\alpha\beta e_j + \frac{1}{2}\alpha b_{jk}^2 + \frac{1}{2}\beta a_{ij}^2  \\
    &= \frac{1}{2}\alpha\beta(e_i+e_k) \ne 0.  
\end{align*}  
Thus $J_{ik} \ne 0$.  

\begin{tcolorbox}[colback=gray!20, colframe=black] 
Note: To get  
\[ 2(a_{ij}b_{jk})^2 = -a_{ij}^2 b_{jk}^2 + (a_{ij}^2 b_{jk})b_{jk} + 2((a_{ij}b_{jk})b_{jk})a_{ij}, \]  
Consider 
\[ (a_{ij}^2, b_{jk}, b_{jk}) = 2(a_{ij}, b_{jk}, a_{ij}b_{jk}). \]  
Expand both associators:  
\begin{itemize}  
    \item $(a_{ij}^2, b_{jk}, b_{jk}) = (a_{ij}^2 b_{jk})b_{jk} - a_{ij}^2 b_{jk}^2$  
    \item $(a_{ij}, b_{jk}, a_{ij}b_{jk}) = (a_{ij}b_{jk})^2 - a_{ij}(b_{jk}(a_{ij}b_{jk}))$  
\end{itemize}  
Hence  
\[ (a_{ij}^2 b_{jk})b_{jk} - a_{ij}^2 b_{jk}^2 = 2\Big((a_{ij}b_{jk})^2 - a_{ij}\big(b_{jk}(a_{ij}b_{jk})\big)\Big). \]  
Rearrange and use commutativity to write $a_{ij}\big(b_{jk}(a_{ij}b_{jk})\big) = \big((a_{ij}b_{jk})b_{jk}\big)a_{ij}$, and $ a_{ij} b_{jk}^2=2(b_{jk}a_{ij} )b_{jk}=2(a_{ij}b_{jk} )b_{jk}$ which yields the displayed equality.  
\end{tcolorbox}

We have shown that the relation $\sim$ on $\{1, \dots, n\}$ defined in the following way  
\[ i \sim j \quad \text{if} \quad J_{ij} \ne 0 \]  
is an equivalence relation (if $ J_{ij}\neq 0 $ and $ J_{jk} \neq 0 $ then $ J_{ik} \neq 0 $).  

Let $S_1, \dots, S_k$ be the distinct equivalence classes of $\{1, \dots, n\}$ with respect to $\sim$.  
Defining $E_i = \sum_{\sigma \in S_i} e_\sigma$ for $i=1, \dots, k$, we get that the $E_i$ are pairwise orthogonal idempotents, $1 = E_1 + \dots + E_k$ and $J = \bigoplus_{i=1}^k J_i$.  

\[ J_1(E_t) = \sum_{i \in S_t} Fe_i + \sum_{i,j \in S_t, i<j} J_{ij} \]  
for any $t=1, \dots, k$. Thus,  
\[ J = J_1(E_1) \oplus \dots \oplus J_1(E_k) \]  
and each $J_1(E_t)$ is an ideal of $J$. \\
Since $J$ is a simple algebra, $k=1 \implies$ all $J_{ij} \ne 0$. 
\end{proof} 

\section{Lecture 11}

\begin{definition}  
For $J$ defined in the previous Lemma 10.4, we say that $e_i$ and $e_j$ are connected if $\exists a_{ij} \in J_{ij}$ is invertible in $ J_1(e_i+e_j) $. $e_i$ and $e_j$ are strictly connected if $\exists a_{ij} \in J_{ij}$ with $a_{ij}^2 = e_i+e_j$.  
\end{definition} 

\begin{corollary}  
For $J$ satisfying the conditions of Lemma 10.4 and $ F=\bar{F} $, for any $i \ne j \in \{1, \dots, n\}$, the idempotents $e_i$ and $e_j$ are strictly connected.  
\end{corollary} 

\begin{proof}  
We have seen in the proof of Lemma 10.4, that $\exists a_{ij} \in J_{ij}$ such that $a_{ij}^2 = \alpha(e_i+e_j)$, with $\alpha \in F, \alpha \ne 0$. Since $F = \bar{F}$, $\exists \beta \in F$ with $\beta^2 = \alpha$. Then $(\beta^{-1}a_{ij})^2 = \beta^{-2}\alpha(e_i+e_j) = e_i+e_j$.  
\end{proof} 

\begin{lemma}  
Let $j$ be a Jordan algebra satisfying Lemma 10.4. Then there exists $e_{ij} \in J_{ij}, i\neq j$ such that  
\begin{enumerate}  
    \item[(i)] $e_{ij}^2 = e_i + e_j$, and  
    \item[(ii)] $2 e_{ij} e_{jk} = e_{ik}$ for each $k \neq i,j$.  
\end{enumerate}  
\end{lemma}

\begin{proof}  
By Corollary 11.2, $e_1$ and $e_{i}, i\neq 1$ are strictly connected $\implies \exists e_{1i} \in J_{1i}$ with $e_{1i}^2 = e_1+e_{i}$. Define $e_{ij} = 2e_{1i}e_{1j}$ for all $i,j > 1, i \neq j$. Then  
\begin{align*}  
    e_{ij}^2  &= 4(e_{1i} e_{1j})^2 = -2e_{1i}^2 e_{1j}^2 + 2(e_{1i}^2 e_{1j})e_{1j} + 4((e_{1i} e_{1j})e_{1j})e_{1i} \\
    &= -2e_{1i}^2 e_{1j}^2 + 2(e_{1i}^2 e_{1j})e_{1j} + 2(e_{1j}^2 e_{1i})e_{1i}  \\
    &= -2(e_1+e_i)(e_1+e_j)+2((e_1+e_i)e_{1j})e_{1j}+2((e_1+e_j)e_{1i})e_{1i} \\
    &= -2e_1+ e_{1j}^2+ e_{1i}^2 = e_i+e_j.
\end{align*}  
For distinct $i,j,k$:  
\begin{itemize}
\item If $ j=1 $, $ 2 e_{i1} e_{1k} =2 e_{1i} e_{1k}= e_{ik} $ by defintion.
\item If $ i=1 $ (or $ k=1 $), $ 2 e_{1j} e_{jk} =2 e_{1j} (2 e_{1j} e_{1k})= 4(e_{1k} e_{1j})e_{1j}=2e_{1k}e_{1j}^2=e_{1k}$.
\item If $ i,j,k $ are distinct from 1, $ 2 e_{ij} e_{jk} = 2 (2 e_{1i} e_{1j}) (2 e_{1j} e_{1k}) = 8(e_{1i} e_{1j})(e_{1j}e_{1k})=-4e_{1j}^2(e_{1i}e_{1k})+4(e_{1j}^2e_{1i})e_{1k}+8((e_{1j}e_{1k})e_{1i})e_{1j}=-2(e_1+e_j)e_{ik}+4((e_1+e_j)e_{1i})e_{1k}+4(e_{jk}e_{1i})e_{1j}=2e_{1i}e_{1k}=e_{ik}$.
\end{itemize}
\end{proof} 

Let $J$ be a finite-dimensional, simple Jordan algebra of capacity $> 2$ and $F = \bar{F}$. Let $e_1, \dots, e_n$ be primitive idempotents, pairwise orthogonal, $1 = e_1 + \dots + e_n$. By Proposition about simple Jordan of capacity 2 for $i \ne j$, $J_1(e_i+e_j) = Fe_i + J_{ij} + Fe_j$ are subalgebras of Clifford type.  

Let $f_{ij}: J_{ij} \to F$ such that $a_{ij}^2 = f_{ij}(a_{ij})(e_i+e_j)$, for any $a_{ij} \in J_{ij}$.  
Then $f$ is quadratic, since $f_{ij}(\alpha a_{ij}) = \alpha^2 f_{ij}(a_{ij})$, $\forall \alpha \in F, \forall a_{ij} \in J_{ij}$.  

Let $\bar{f}_{ij}$ be a bilinear form associated to quadratic $f_{ij}$.  
\[ \bar{f}_{ij}(a_{ij}, b_{ij}) = \frac{1}{2} (f_{ij}(a_{ij}+b_{ij}) - f_{ij}(a_{ij}) - f_{ij}(b_{ij})) \]  
then  
\[ \bar{f}_{ij}(a_{ij}, b_{ij})(e_i+e_j) = \frac{1}{2}((a_{ij}+b_{ij})^2 - a_{ij}^2 - b_{ij}^2) = a_{ij}b_{ij}. \]  
Thus $\bar{f}_{ij}$ is non-degenerate: since for $a_{ij} \in J_{ij}$ such that $\bar{f}_{ij}(a_{ij}, J_{ij})=0$ we get $a_{ij}J_{ij}=0$ and next $F a_{ij} \triangleleft J_1(e_i+e_j)$ with $(F a_{ij})^2=0$, then $a_{ij}=0$.  

\begin{lemma}  
Let $J$ be a fin-dim simple Jordan algebra, $F=\bar{F}$ of capacity $n>2$ ($1=e_1+\dots+e_n$). Let $i,j,k$ be all distinct. Let $f_{ij}, f_{ik}, f_{jk}$ be quadratic forms for $J_1(e_i+e_j)$, $J_1(e_i+e_k)$ and $J_1(e_j+e_k)$.  
Then $f_{ik}(2a_{ij}b_{jk}) = f_{ij}(a_{ij})f_{jk}(b_{jk})$, $\forall a_{ij} \in J_{ij}, b_{jk} \in J_{jk}$.  
\end{lemma}  

\begin{proof}  
\begin{align*}  
    2(a_{ij}b_{jk})^2&= -a_{ij}^2 b_{jk}^2 + (a_{ij}^2 b_{jk})b_{jk} + 2((a_{ij}b_{jk})b_{jk})a_{ij}  \\
    &= -a_{ij}^2 b_{jk}^2 + (a_{ij}^2 b_{jk})b_{jk} + (b_{jk}^2 a_{ij})a_{ij}  \\
    &= -f_{ij}(a_{ij})(e_i+e_j)f_{jk}(b_{jk})(e_j+e_k) + (f_{ij}(a_{ij})(e_i+e_j)b_{jk})b_{jk}+ (f_{jk}(b_{jk})(e_j+e_k)a_{ij})a_{ij} \\
    &= -f_{ij}(a_{ij})f_{jk}(b_{jk})e_j  + \frac{1}{2} f_{ij}(a_{ij})f_{jk}(b_{jk})(e_j+e_k) + \frac{1}{2} f_{jk}(b_{jk})f_{ij}(a_{ij})(e_i+e_j).  \\
     &= f_{ij}(a_{ij})f_{jk}(b_{jk})\left(-e_j + \frac{1}{2}(e_j+e_k) + \frac{1}{2}(e_i+e_j)\right) = \frac{1}{2}f_{ij}(a_{ij})f_{jk}(b_{jk})(e_i+e_k). 
\end{align*}  
From the other side, since $ a_{ij}b_{jk} \in J_{ik}$, we get $2(a_{ij}b_{jk})^2 = 2f_{ik}(a_{ij}b_{jk})(e_i+e_k) = \frac{1}{2}f_{ik}(2a_{ij}b_{jk})(e_i+e_k)$ since $ f_{ik} $ is quadratic.  
\end{proof}

\begin{lemma}  
Under the assumption of Lemma 11.4, all subspaces $J_{ij}$ with $i \ne j$ are isomorphic as vector spaces.  
\end{lemma}  

\begin{proof}  
The linear mapping $\varphi_{ijk}: J_{ij} \to J_{ik}$ defined by $\varphi_{ijk}(a_{ij}) = a_{ij} e_{jk}$ is injective. Indeed, suppose $a_{ij} e_{jk} = 0$ for some $a_{ij} \in J_{ij}$.  
\[ 0 = 2(a_{ij} e_{jk}) e_{jk} = e_{jk}^2 a_{ij} = (e_j+e_k) a_{ij} = \frac{1}{2} a_{ij}. \]  
Then $\varphi_{ijk}$ is injective for all $i,j,k$. Analogously, $\varphi_{kji}: J_{kj}=J_{jk} \to J_{ji}=J_{ij}$ is injective. Then we get the following inequality $\dim J_{ij} \le \dim J_{ik} \le \dim J_{ji} = \dim J_{ij}$. This implies all $J_{ij}$ have the same dimension and are therefore isomorphic.  
\end{proof} 

\begin{definition}  
We call an algebra $A$ a composition algebra if $A$ has unity and a non-degenerate quadratic form $n$ in $A$ such that $n(xy)=n(x)n(y)$, for all $x, y \in A$.  
\end{definition}  

\begin{remark}  
Any composition algebra $A$ is quadratic, that is $\forall a \in A, a^2 = t(a)a - n(a)1$, with $t(a), n(a) \in F$.  
It also has an involution $a \to \bar{a}$ such that $a+\bar{a} = t(a)1$ and $a\bar{a}=\bar{a}a=n(a)1$, $\forall a \in A$.  
\end{remark}

\begin{proposition}  
Let $J$ be a Jordan algebra satisfying the conditions of Lemma 11.4 with $e_{ij}$ defined as in Lemma 11.3. Consider $D=J_{12}$ with the following multiplication:  
\[ a*b = 8(ae_{23})(be_{13}) \quad \text{for } a,b \in D \]  
and denote $f_{12}$ by $n$. Then  
\[ n(a*b)=n(a)n(b) \quad \text{and} \quad a*e_{12} = e_{12}*a = a. \]  
$\forall a,b \in D$. Then $D$ is a composition algebra with quadratic form $n$ and unity $e_{12}$.  
\end{proposition}  

\begin{proof} 
We already show that $ \bar{f_{12}} $ is non-degenerated.  
\begin{align*}  
n(a*b)  
  &= f_{12}\big(8\,(a e_{23})(b e_{13})\big)  
   = f_{13}\big(2 a e_{23}\big)\,f_{23}\big(2 b e_{13}\big)  
   = f_{12}(a)\,f_{12}(b)  
   = n(a)\,n(b) \\
\text{Moreover}\quad  
a*e_{12}  
  &= 8\,(a e_{23})(e_{12} e_{13})  
   = 4\,(a e_{23}) e_{23}  
   = 2 a e_{23}^{2}  
   = 2 a(e_{2}+e_{3})  
   = a, \\
e_{12}*a  
  &= 8\big((e_{12} e_{23})(a e_{13})\big)  
   = 4\,e_{13}(a e_{13})  
   = 2ae_{13}^2=a,  
\end{align*}  
 D is a composition algebra.  
\end{proof}

We have shown that $J_{12} = D$ is a composition algebra with unit element $e_{12}$ and quadratic form $n = f_{12}$.  
We will show that $t(a)(e_1+e_2)=2ae_{12}, \forall a \in D$.  
Let $a \in D$.   
\begin{align*}  
    a*a &= 8(ae_{23})(ae_{13}) = -4a^2(e_{23}e_{13}) + 4((ae_{23})a)e_{13} + 4(a(ae_{13}))e_{23}+4a(a(e_{13}e_{23})) \\
    &= -2n(a)(e_1+e_2)e_{12} + 2(a^2 e_{23})e_{13} + 2(a^2e_{13})e_{23}+2a(ae_{12}) \\
    &= -2n(a)e_{12} + 2(n(a)(e_1+e_2)e_{23})e_{13}+2(n(a)(e_1+e_2)e_{13})e_{23} + 2(ae_{12})a \\
    &= n(a)(-2e_{12} + \frac{1}{2}e_{12} + \frac{1}{2}e_{12}) + 2(ae_{12})a = 2(ae_{12})a - n(a)e_{12} \\
    &= 2(ae_{12})a - n(a)1_D  
\end{align*}  
As $D$ is quadratic we have $2ae_{12} \in F1$ and then it follows that $t(a) (e_1+e_2)= 2ae_{12}$.  

Additionally we have $\bar{a} = t(a)1_D - a = 2(ae_{12})e_{12} - a, \forall a \in D$.  


\section{Lecture 12}

Let $A$ be an algebra with involution $a \to \bar{a}$. We may extend this involution to a matrix algebra $M_n(A)$  
\[ \overline{\{a_{ij}\}} = \{\bar{a}_{ji}\}, \quad \text{for any } \{a_{ij}\} \in M_n(A). \]  

We call a matrix $x \in M_n(A)$ Hermitian if $\bar{x}=x$. The set of Hermitian matrices is denoted by $H_n(A)$. In fact $H_n(A)$ is a subalgebra of $M_n(A)^+$ since for any $x,y \in H_n(A)$, $x \cdot y = \frac{1}{2}(xy+yx) \in H_n(A)$.  

For $a \in A$ and $i,j=1, \dots, n$ define $a[ij] = ae_{ij} + \bar{a} e_{ji}$. (here $e_{ij}$ denotes elementary matrix with only 1 on entry $(i,j)$ while other entries zero)

we have $a[ij] \in H_n(A)$ and $a[ji] = \bar{a}[ij]$. $x \in H_n(A)$ then $x$ has a form  
\[  
x =   
\begin{pmatrix}  
\alpha_1 & a_{12} & \dots & a_{1n} \\
\bar{a}_{12} & \alpha_2 & & \\
\vdots & & \ddots & \\
\bar{a}_{1n} & & & \alpha_n  
\end{pmatrix}  
\]  
with $\bar{\alpha}_i = \alpha_i, i=1,\dots,n$. For each $i=1,\dots,n$ set $e_i=e_{ii}$, then $e_1, \dots, e_n$ are pairwise orthogonal idempotents $1=e_1+\dots+e_n$. In the full Peirce decomposition  
\[ H_n(A) = \sum_{i=1}^n H_n(A)_{ii} + \sum_{1 \le i < j \le n} H_n(A)_{ij} \]  
For distinct $i,j$ we have $H_n(A)_{ii} = \{a[ii] : a\in A \}$ and $H_n(A)_{ij} = \{a[ij] : a\in A \}$. Indeed, $a[ii] = (a+\bar{a})e_{ii}$ and  
\begin{align*}  
    a[ij] \cdot e_i &= \frac{1}{2}((ae_{ij}+\bar{a}e_{ji})e_{ii} + e_{ii}(ae_{ij}+\bar{a}e_{ji})) \\
    &= \frac{1}{2}(\bar{a}e_{ji}+ae_{ij}) = \frac{1}{2}a[ij].  
\end{align*} 

\begin{lemma}  
Let $i,j,k \in \{1,\dots,n\}$ distinct. Then  
\begin{enumerate}  
    \item[(i)] $2a[ij] \cdot b[jk] = (ab)[ik]$;  
    \item[(ii)] $2a[ii] \cdot b[ik] = (a+\bar{a})b[ik]$;  
    \item[(iii)] $2a[ij] \cdot b[jj] = a(b+\bar{b})[ij]$;  
    \item[(iv)] $2a[ij] \cdot b[ij] = (a\bar{b})[ii] + (\bar{a}b)[jj]$;  
    \item[(v)] $2a[rs] \cdot b[kl] = 0$ if $\{r,s\} \cap \{k,l\} = \emptyset$.  
\end{enumerate}  
\end{lemma}  
\begin{proof}  
The proof is straightforward. 
\begin{enumerate}  
    \item[(i)] $((ae_{ij}+\bar{a}e_{ji})(be_{jk}+\bar{b}e_{kj}) + (be_{jk}+\bar{b}e_{kj})(ae_{ij}+\bar{a}e_{ji}))= (ab)e_{ik} + (\bar{b}\bar{a})e_{ki} = (ab)[ik] $;  
    \item[(ii)] $((ae_{ii}+\bar{a}e_{ii})(be_{ik}+\bar{b}e_{ki}) + (be_{ik}+\bar{b}e_{ki})(ae_{ii}+\bar{a}e_{ii}))= (ab)e_{ik} +(\bar{a}b)e_{ik} + (\bar{b}a)e_{ki} +(\bar{b}\bar{a})e_{ki}= (a+\bar{a})b[ik] $;    
    \item[(iii)] $((ae_{ij}+\bar{a}e_{ji})(be_{jj}+\bar{b}e_{jj}) + (be_{jj}+\bar{b}e_{jj})(ae_{ij}+\bar{a}e_{ji}))= (ab)e_{ij} +(a\bar{b})e_{ij} + (b\bar{a})e_{ji} +(\bar{b}\bar{a})e_{ji}= a(b+\bar{b})[ij] $; 
    \item[(iv)] $((ae_{ij}+\bar{a}e_{ji})(be_{ij}+\bar{b}e_{ji}) + (be_{ij}+\bar{b}e_{ji})(ae_{ij}+\bar{a}e_{ji})) = a\bar{b}e_{ii} + \bar{a}b e_{jj} + b\bar{a}e_{ii} + \bar{b}ae_{jj} = (a\bar{b})[ii]+(\bar{a}b)[jj]  $;
    \item[(v)] $ ((ae_{rs}+\bar{a}e_{sr})(be_{kl}+\bar{b}e_{lk}) + (be_{kl}+\bar{b}e_{lk})(ae_{rs}+\bar{a}e_{sr}))=0 $.  
\end{enumerate}
\end{proof}  

\begin{definition}  
An algebra $A$ is called alternative if $(a,a,b)=0$ and $(a,b,b)=0$ for all $a,b \in A$.  
\end{definition}  

\begin{theorem}  
Let $n \ge 3$. If $H_n(A)$ is a Jordan algebra, then $A$ is an alternative algebra (associative, if $n \ge 4$) and for any $a \in A$, $a+\bar{a} \in N(A)$, the associative center of A.  
\end{theorem}  

\begin{proof}  
For $i,j \in \{1, \dots, n\}$ we have shown that $(H_n(A))_{ij} = \{a[ij]: a \in A\}$. Then we also have $(J_{ij}, J_{jk}, J_{kl})=0$ for $\{i,j,k,l\}$ distinct, since $(a_{ij}, b_{jk}, c_{kl}) = 2(a_{ij} \cdot e_i, b_{jk}, c_{kl}) = -2(a_{ij} c_{kl}, e_k, b_{jk}, e_i)  -2(e_i c_{kl}, b_{jk}, a_{ij}) = 0$.  

Thus for $a,b,c \in A$ and $\{i,j,k,l\}$ distinct, $(a[ij], b[jk], c[kl])=0$. In particular $4(a[ij] \cdot b[jk]) \cdot c[kl] = 2(ab)[ik] \cdot c[kl] = ((ab)c)[il]$. And $4 a[ij] \cdot (b[jk] \cdot c[kl]) = (a(bc))[il]$. This implies $(ab)c = a(bc)$. If $n > 3$, A is associative.  

Using the above for $l=k$ we obtain $(J_{ij}, J_{jk}, J_{ik}) = 0$. Then for $a,b,c \in A$ and $k \ne i,j$ we have $(a[ij], b[jk], c[kk])=0$. Rewriting both sides we get: $4(a[ij] \cdot b[jk]) \cdot c[kk] = 2(ab)[ik] \cdot c[kk] = ((ab)(c+\bar{c}))[ik]$. $4 a[ij] \cdot (b[jk] \cdot c[kk]) = 2a[ij] \cdot b(c+\bar{c})[jk] = (a(b(c+\bar{c})))[ik]$. This implies $(a,b,c+\bar{c})=0$, and analogously $(a+\bar{a},b,c)=0$.  

Finally, $(a_{ij}, e_{ji}, b_{jk}) = (a_{ij}, b_{ji}, e_{jk}) = b_{ij}(a_{ij}, e_i, e_{jk}) + (a_{ij}, b_{ij}, e_{jk})e_i = 0$. Thus $(a[ij], b[ji], c[jk]) = 0$. Rewriting this gives $(a, b+\bar{b}, c)=0$. Consequently $a+\bar{a} \in N(A)$.   
\end{proof}  

\begin{exercise}
Prove that if $ n=3 $ then $ A $ is alternative.  
\end{exercise}

Let $J$ be finite-dimensional simple Jordan algebra with capacity $> 2$ and $F=\bar{F}$. $e_1, \dots, e_n$ primitive, pairwise orthogonal idempotent, $e_1+\dots+e_n=1$. $\exists e_{ij} \in J_{ij}$ such that $e_{ij}^2=e_i+e_j$ and $2e_{ij} \cdot e_{jk} = e_{ik}$. $D=J_{12}$ is a composition algebra with trace $tr(a)$ and involution $a \to \bar{a}$.  

We show that $J \cong H_n(D)$. To do this we will define $a_{ij} \in J_{ij}$ such that $\varphi: a_{ij} \to a[ij]$, $a \in D$ defines isomorphism of Jordan algebras.  

For each $a \in D$ and $i,j=1, \dots, n$  
\begin{itemize}  
    \item for $i \ne 2$ and $j \ne 1$: \quad $a_{1j} = 2ae_{2j}$, \quad $a_{j1} = \bar{a}_{1j}$.  
    \item $a_{i2}=2ae_{i1}$, \quad $a_{2i}=\bar{a}_{i2}$  
    \item for $i=1, \dots, n$: \quad $a_{ii} = tr(a)e_i$  
    \item for $i,j > 2$, distinct: \quad $a_{ij} = 2a_{i2}e_{2j} = 4(ae_{i1})e_{2j}$ 
\end{itemize}

\begin{lemma}  
\begin{enumerate}  
    \item[(i)] For any $i \neq j$, $a_{ji} = \overline{a_{ij}}$.  
    \item[(ii)] $2 a_{ij} b_{jk} = (a * b)_{ik}$;   
    \item[(iii)] $2 a_{ii} b_{ij} = \operatorname{tr}(a) b_{ij}$; for $a,b \in D$.  
    \item[(iv)] $2 a_{ij} b_{jj} = \operatorname{tr}(b) a_{ij}$;  
    \item[(v)] $2 a_{ij} b_{ji} = \operatorname{tr}(a * b)(e_i+e_j) = (a*  b)_{ii} + (b* a)_{jj}$.  
\end{enumerate}  
Where $ i,j,k $ distinct and $ a,b \in D =J_{12} $.  
\end{lemma} 

\begin{theorem}  
Let $J$ be a finite-dimensional simple reduced Jordan algebra of capacity $n \ge 3$, with $n$ idempotents pairwise orthogonal, strictly connected.  
Then $J \simeq H_n(A, *)$, with $A$ a composition algebra, which is associative if $n>3$.  
\end{theorem}  

\begin{proof}  
It is enough to show that $\varphi: a_{ij} \to a[ij]$, $a \in A$ defines a bijection between vector spaces $J$ and $H_n(A, *)$. Moreover by Lemma 12.5, $\varphi$ is an isomorphism of Jordan algebras.  
\end{proof} 

\section{Lecture 13}

Let $X$ be a set of variables. Denote by $F\langle X \rangle$ free associative algebra generated by $X$.  

The subalgebra of $F\langle X \rangle^{(-)}$ generated by $X$ is the free Lie algebra generated by $X$, we will denote it as $\operatorname{Lie}\langle X \rangle$. Its elements are called Lie elements. These are at least two ways to check whether a given element of $F\langle X \rangle$ is Lie element or not.  

Denote by $[x_1,\dots,x_n]$ the element $[[...[[x_1,x_2],x_3],...x_n]]$. The Dynkin-Specht-Wever criterion guarantees that for $f \in F\langle X \rangle$ is homogeneous of degree $n$ then $f$ is Lie element $\iff [f]=nf$.  

The other criterion considers the coalgebra structure of $F\langle X \rangle$, which is given by a comultiplication $\Delta: F\langle X \rangle \to F\langle X \rangle \otimes F\langle X \rangle$, in which each $x \in X$ is primitive that is, $\Delta(x)=1\otimes x + x\otimes 1$. The Friedrichs criterion says that $f \in F\langle X \rangle$ is Lie element $\iff f$ is primitive.  

Analogously, the subalgebra of $F\langle X \rangle^{(+)}$ generated by $X$ is denoted by $SJ\langle X \rangle$ and its elements are called Jordan elements. Contrary to the case of Lie algebras, $SJ\langle X \rangle$ is not free Jordan algebra (since it is not closed under homomorphic images) and we do not know criterion to decide if element $f \in F\langle X \rangle$ is Jordan element. Nevertheless, $SJ\langle X \rangle$ has the following universal property of free algebras.  

\begin{proposition}Let $J$ be a special Jordan algebra. Then any mapping $\varphi: X \to J$ can be uniquely extended to a homomorphism $\hat{\psi}: SJ\langle X \rangle \to J$ and $\hat{\psi}$ is unique, i.e. the diagram  
\[  
\begin{tikzcd}  
X \arrow[r, "\varphi"] \arrow[d] & J \\
SJ\langle X \rangle \arrow[ru, "\exists ! \hat{\psi}"'] &  
\end{tikzcd}  
\]  
commute.  
\end{proposition}
\begin{proof} 
Let $J$ be special, therefore there exists $A$, an associative algebra, such that $J \hookrightarrow A^{(+)}$. Therefore we have $\varphi: X \to J \hookrightarrow A^{(+)}$.  

Since $F\langle X \rangle$ is a free associative algebra, there exists $\hat{\varphi}: F\langle X \rangle \to A$ an associative homomorphism which extends $\varphi$. Then $\hat{\varphi}: F\langle X \rangle^{(+)} \to A^{(+)}$ is a homomorphism of Jordan algebras and taking the restriction to $SJ\langle X \rangle$, we obtain $\hat{\varphi}: SJ\langle X \rangle \to A^{(+)}$ is a homomorphism of Jordan algebras. Moreover, since $\hat{\varphi}(x) = \varphi(x) \in J$, we obtain $\operatorname{Im}(\hat{\varphi}) \subseteq J$.  
\end{proof}

\begin{remark} The class of special Jordan algebras is closed under direct sums and subalgebras. But we will see that they are not closed under homomorphic images. Thus, not forming a variety of algebras. This class forms what we call quasi-variety, which consists of algebras satisfying quasi-identities (a quasi-identity is a relation of type $w=0 \implies v=0$, for example $x^2 = 0 \implies x=0$ define the quasi-variety of algebras without nilpotent elements). Each quasi-variety contains free algebras.  
\end{remark}
Let $J$ be a Jordan algebra. A Jordan homomorphism $\varphi: J \to A^{(+)}$, where $A$ is an associative algebra, is called specialization of $J$. Among all specializations of $J$ we have one which is universal.  

\begin{definition} [Special universal algebra]  
An associative algebra $S(J)$ with a specialization $i: J \to S(J)^{(+)}$ is called special universal algebra of $J$ if for any homomorphism $\varphi: J \to A^{(+)}$ with $A$ an associative algebra there exists a unique $\tilde{\varphi}: S(J) \to A$ homomorphism of associative algebras such that $\varphi = \tilde{\varphi} \circ i$.  
\[  
\begin{tikzcd}  
J \arrow[r, "\varphi"] \arrow[d, "i"] & A^{(+)} \\
S(J) \arrow[ru, "\exists ! \tilde{\varphi}"'] &  
\end{tikzcd}  
\]  
\end{definition}
\begin{proposition} Let $J$ be a Jordan algebra. Consider the tensor algebra $T(J)$ generated by a vector space $J$:  
\[  
T(J) = F \cdot 1 \oplus J \oplus J \otimes J \oplus J \otimes J \otimes J \oplus \dots  
\]  
and an ideal $I$ of $T(J)$ generated by $a \otimes b+b\otimes a - a \circ b$, $a \in J$.  
Then $S(J) = T(J)/I$.  
\end{proposition}

\begin{proof} 
Let $i_J : J \to T(J)/I$ be a mapping defined by $i(a) = a+I, \forall a \in J$. Let $\varphi: J \to A^{(+)}$ be a homomorphism of Jordan algebras, with A being associative algebra. In particular, $\varphi$ is linear. Since $T(J)$ is a free associative algebra over $J$, there exists a unique homomorphism $\tilde{\varphi}$ of associative algebras $\tilde{\varphi}: T(J) \to A$ such that $\varphi = \tilde{\varphi} \circ i$. Observe that $I \subset ker \tilde{\varphi}$, $\tilde{\varphi}(a \otimes b+b\otimes a - a \circ b) = 0$. Thus there exists $\bar{\varphi}: T(J)/I \to A$ such that $\varphi = \bar{\varphi} \circ i$.  
\end{proof}

\begin{exercise}
    Let $J$ be Jordan algebra. Show that $J$ special Jordan algebra $\Leftrightarrow i: J \to S(J)$ is injective.  
\end{exercise}
\begin{exercise}
    Show that $S(SJ\langle X \rangle) = F\langle X \rangle$.  
\end{exercise}
\begin{exercise}
    Let $J$ be Jordan algebra of dimension $n$. Show that dim $S(J) \leq 2^n$.  
\end{exercise}
\begin{remark}If $X = \{x\}$ then $SJ\langle X \rangle = F\langle X \rangle$. When $|X| > 1$ this is no longer true. For if $X = \{x, y\}$ then $x^2, y^2, xy+yx \in SJ\langle X \rangle$ but $[x,y] \notin SJ\langle X \rangle$.  
\end{remark}

The free associative algebra $F\langle X \rangle$, $X = \{x_1, x_2, ..., x_n, ...\}$ has natural involution * defined by the property that $x_i^* = x_i$. $(x_{i_1}...x_{i_n})^* = x_{i_n}...x_{i_1}, \forall x_{i_1}, ..., x_{i_n} \in X$.  

Denote $H\langle X \rangle = H(F\langle X \rangle, *) = \{f \in F\langle X \rangle : f^* = f\}$.  

Then $X \subseteq H\langle X \rangle$ and $H\langle X \rangle$ is a Jordan subalgebra of $F\langle X \rangle^{(+)}$. Then $SJ\langle X \rangle \subseteq H\langle X \rangle$.  

And all Jordan element in $F\langle X \rangle$ are symmetric. For any $f \in F\langle X \rangle$ we will denote ${f} = f + f^*$.  

Let $ k=\{x_1x_2x_3x_4\} $ we show that $ k \notin SJ\<x_1, x_2, x_3, x_4 \> $.

In fact, suppose $k =  j(x_1, x_2, x_3, x_4)  \in SJ\langle x_1, x_2, x_3, x_4 \rangle$. Consider the Grassmann algebra $G: \operatorname{alg} < 1, e_i |e_i e_j = - e_j e_i, i=1,2,\ldots >$. $G$ has basis $1, e_{i_1}e_{i_2}\ldots e_{i_k}$ with $i_1 < i_2< \ldots i_k$. 

Let $X = \{x_1, x_2, x_3, x_4\}$ and define $\varphi: X \to G$ by $\varphi(x_i) = e_i, i=1,...,4$. Then there exists unique homomorphism of associative algebras $\tilde{\varphi}: F\langle X \rangle \to G$ such that $\tilde{\varphi}|_X = \varphi$. Then we have $\tilde{\varphi}(k) = e_1e_2e_3e_4 + e_4e_3e_2e_1 = 2e_1e_2e_3e_4$.  

From the other side $\varphi(k) = \tilde{\varphi}(j( x_1, x_2 , x_3, x_4) ) = j(\tilde{\varphi}(x_1),\tilde{\varphi}(x_2), \tilde{\varphi}(x_3), \tilde{\varphi}(x_4))=0$, since $e_i \cdot e_j = 0$. Thus $ k $ can not be an element of $SJ\langle x_1, x_2, x_3, x_4 \rangle$.

\begin{definition}
An element $\{x_{i_1} x_{i_2} x_{i_3} x_{i_4}\} \in H\langle X \rangle$ is called a tetrad.  
\end{definition}

\begin{theorem}[Cohn]
    Jordan algebra $H\langle X \rangle$ is generated by $X \cup \{x_{i_1} x_{i_2} x_{i_3} x_{i_4}| x_{i_k} \in X\}$.  
\end{theorem}

\begin{proof}
Let $H_0\langle X \rangle$ be a subalgebra of $F\langle X \rangle^{(+)}$ generated by $X \cup \{x_{i_1} x_{i_2} x_{i_3}x_{i_4} | x_{i_k} \in X\}$. We need to show that for any monomial $u \in F\langle X \rangle$ we have $\{u\} \in H_0\langle X \rangle$. We prove it using induction in $\deg u$ of $u$.  

We have $\{x_1 x_2\} = x_1 x_2 + x_2 x_1 = x_1 \circ x_2$ and $2\{x_1 x_2 x_3\} = 2(x_1 \circ x_2) \circ x_3 + (x_2 \circ x_3) \circ x_1 - (x_1 \circ x_3) \circ x_2$.  

Then of $\deg u \le 4$, we have $\{u\} \in H_0\langle X \rangle$. 

Suppose $n > 4$ we write $u \equiv 0$ if $u \in H_0\langle X \rangle$. We observe that  
$0 \equiv x_1 \circ \{x_2 x_3 ... x_n\} = x_1 x_2 ... x_n + x_1 x_n \ldots x_2 +x_2 ... x_n x_1+x_n \ldots x_2x_1 = \{x_1...x_n\} + \{x_2 ... x_n x_1\}$.

Then $\{x_1 x_2 ... x_n\} = -\{x_2 ... x_n x_1\} \equiv (-1)^n \{x_1 x_2 ... x_n\}$.  

Then if $n$ is odd, $\{x_1 ... x_n\} = 0$.  

If $n$ is even, we write $n = 2k$, then $ 0 \equiv \{x_1x_2\} \circ\{x_3\ldots x_n\}=(x_1x_2+x_2x_1)(x_3\ldots x_n+x_n\ldots x_3)+(x_3\ldots x_n+x_n\ldots x_3)(x_1x_2+x_2x_1)=\{x_1\ldots x_n\}+\{x_2x_1x_3\ldots x_n\} +\{x_1x_2x_n\ldots x_3\}+\{x_2x_1x_n\ldots x_3\}$ 

Observe that $\{x_1 x_2 x_n ... x_3\} = \{x_3 ... x_n x_2 x_1\} = \{x_2 x_1 x_3... x_n \}$.  

Since $x_3 ... x_n$ has even length. Also $\{x_2 x_1 x_n  ...  x_3\} = \{x_3...x_n x_1x_2\}=\{x_1\ldots x_n\}$. 

Concluding $0 \equiv \{x_1 x_2 x_3 ... x_n\} + \{x_2 x_1 x_3 ... x_n\}$.  

and therefore $\{x_1 x_2 ... x_n\} = \{x_n ... x_2 x_1 \} \equiv (-1)^{\frac{n(n-1)}{2}} \{x_1  ... x_n\}$.  

Finally we need to check the case when $n = 4k$:  
$0 = \{x_1 x_2 x_3 x_4 \} \circ \{x_5 ... x_n\} = x_1 ... x_4 x_5 ... x_n +x_4 ... x_1 x_5 ... x_n+x_1 ... x_4 x_n ... x_5 +x_4 ... x_1 x_n ... x_5  x_5 ... x_n x_1 ... x_4 +x_5 ... x_n x_4 ... x_1+x_n ... x_5 x_1 ... x_4 +x_n ... x_5 x_4 ... x_1=\{x_1\ldots x_n\}+\{x_4 ... x_1 x_5 ... x_n\}+\{x_1 ... x_4 x_n ... x_5\}+\{ x_5 ... x_n x_1 ... x_4 \}=2(\{x_1\ldots x_n\}+\{x_4 ... x_1 x_5 ... x_n\})=4\{x_1\ldots x_n\}  $.
\end{proof}

\begin{corollary}
    If $|X| \leq 3$, $SJ(X) = H\langle X \rangle$.  
\end{corollary}

\begin{remark}
    There is direct proof given of Corollary given by Shirshov. He described an algorithm to write symmetric element in 3 letters as Jordan elements.  
\end{remark}

\section{Lecture 14}

\begin{theorem}[Cohn Criterium] 
If $J$ is special Jordan algebra. If $I \triangleleft J$, denote $\hat{I} = ideal_{S(J)} \langle I \rangle$.  
Then $J/I$ is special $\Leftrightarrow \hat{I} \cap J = I$.  
\end{theorem}
\begin{proof}
Let $A$ be associative algebra such that $J/I \hookrightarrow A^{(+)}$. Then we have specialization $J \to J/I \hookrightarrow A^{(+)}$ then there exists unique homomorphism of associative algebras $\tilde{\varphi}: S(J) \to A, \varphi = \tilde{\varphi} \circ i$ and $i$ is injective then $I = Ker \varphi = \varphi^{-1} \hat{I}$.  
$J \xrightarrow{i} S(J) \xrightarrow{\tilde{\varphi}} A$ We will show that $\hat{I} \cap J = I$. The inclusion $I \subseteq \hat{I} \cap J$ is clear, from the other side, $\hat{I} \subseteq \operatorname{Ker} \tilde{\varphi}$. Thus $\hat{I} \cap J \subseteq J \cap \operatorname{Ker} \tilde{\varphi} = I$.  

Conversely if $\hat{I} \cap J = I$, let $A = S(J)/\hat{I}$. Consider canonical projection of $S(J)$ to $A$. Since $\varphi: S(J) \to A$ is homomorphism of associative algebras, $\varphi_J: J \to A^{(+)}$ is a specialization then $\operatorname{Ker} \varphi_J = J \cap\operatorname{Ker} \varphi = J \cap \hat{I} = I$. Therefore $J/I = \operatorname{Im} \varphi_J \hookrightarrow A^{(+)}$.  
$\Rightarrow J/I$ is special.  
\end{proof}

\begin{corollary} 
    Let $ SJ\langle X \rangle/I=J$, $ J $ special Jordan algebra. Then $S(J) = F\langle X \rangle/\hat{I}$, where $\hat{I} = ideal_{F\langle X \rangle}\langle I \rangle$.  
\end{corollary}

\begin{definition} 
When $I\neq \hat{I } \cap SJ\langle X \rangle$, the elements $k \in \hat{I} \cap SJ\langle X \rangle$, $k \notin I$, are called Cohn elements for ideal $I$.  
\end{definition}

\begin{theorem} 
Let $I = ideal_{SJ\langle x, y, x \rangle} \langle x \cdot y \rangle$. Then $SJ\langle x, y, z \rangle / I$ is exceptional algebra.  
\end{theorem}

\begin{proof} 
 Let $k = \{(x \cdot y)xyz\} \in F\langle x, y, z \rangle$. Then $k \in \hat{I}$. Moreover $k^* = k$ then $k \in H\langle x, y, z \rangle = SJ\langle x, y, z \rangle$.  
$\Rightarrow k \in \hat{I} \cap SJ\langle x, y, z \rangle$. We will show that $k \notin I$ (then $I \neq \hat{I} \cap SJ\langle x, y, z \rangle $, by Cohn criterion $J$ is not special).  

Suppose that $k \in I$, then there exists $j(t, x, y, z) \in SJ \langle t, x, y, z \rangle$, such that all monomials in $j$ contain $t$ and $j(x\cdot y,x, y, z) = k = \{(x \cdot y)xyz\}$ and we may suppose that $J$ is homogeneous since $k$ is multihomogeneous of degree $(2, 2, 1)$.  

In $j(t,x, y, z)$ we have $\deg_z j = 1$ and $\deg_t j = 1$ or $2$. 

If $\deg_t j = 2$ then $j =j(t,z)= \alpha t^2 \cdot z + \beta (t \cdot z) \cdot t$. Then $\alpha (x \cdot y)^2 \cdot z + \beta ((x \cdot y) \cdot z) \cdot (x \cdot y) = k = \{(x \cdot y)xyz\}$. In the right hand side of the identity there are only monomials starting with $z$ and ending in $z$. Thus $\beta = 0$. Again comparing monomials, we conclude that $\alpha = 0$. 

Then $\deg_t j = 1$ and therefore degree of $j$ is $ (1, 1, 1, 1)$.  
since $j(t, x, y, z)^* = j(t, x, y, z)$ we have $j \in H \langle t, x, y, z \rangle$, meaning that $j$ is a linear combinations of thetrads in $t, x, y$ and $z$. And letter $t$ has to appear as first or last letter in each term.  
$j(t, x, y, z) = \alpha_1 \{t xyz\} + \alpha_2 \{t yxz\} + \alpha_3 \{x t y z\} + \alpha_4 \{x y t z\} + \alpha_5 \{ytxz\} + \alpha_6 \{y x t z\}$.  

Recall that $J(x \cdot y, x,y, z) = \{(x \cdot y)xyz\}$.  

Substituting and comparing the coefficients of the monomials $x^2y^2 z$, $y^2 xz$, $ xy^2yz $, $yx^2yz$, $xyxyz$ we obtain:  

$0 = \frac{1}{2} \alpha_3 + \frac{1}{2} \alpha_4$, $0 = \frac{1}{2} \alpha_5 + \frac{1}{2} \alpha_6$, $0 = \frac{1}{2} \alpha_2 + \frac{1}{2} \alpha_4$, $\frac{1}{2} = \frac{1}{2} \alpha_1 + \frac{1}{2} \alpha_6$, $\frac{1}{2} = \frac{1}{2} \alpha_1 + \frac{1}{3} \alpha_3 + \frac{1}{2} \alpha_4$.  

Then we conclude that $\alpha_1 = 1$, $\alpha_6 = \alpha_5 = 0$  
$\alpha_3 = \alpha_2 = - \alpha_4$ and thus $j = \{txyz\} + \alpha_2 \{tyxz\} + \alpha_2 \{xtyz\} - \alpha_2 \{xytz\} $.  

For a coefficient of $yxyxz$, we get $0 =  \frac{\alpha_2}{2}$.  
Consequently, $j = \{txyz\}$, which is not Jordan element which leads to a contradiction. Thus $k \notin I$.  
\end{proof}

\begin{theorem} 
The free Jordan algebra in two generators $\operatorname{Jord}[x, y]$ is special. It imples $ \operatorname{Jord}[x,y] \simeq SJ\< x,y> $ 
\end{theorem}
We are not going to prove it. We will describe free Jordan algebras.

Let $T\{V\}$ be the free non-associative algebra over vector space $V$:  
\[T\{V\} = F1 \oplus V \oplus (V \otimes V) \oplus ((V \otimes V) \otimes V) \oplus (V \otimes (V \otimes V)) \oplus ...\]  
on where $V^{\otimes n} = \underset{i+j = n}{\oplus} V^{\otimes i} \otimes V^{\otimes j}$. The algebra $T(V)$ has the following universal properties.  

Let $ A $ be $F$-algebra and $\varphi: V \to A$ be linear mapping. Then there exists unique homomorphism $\tilde{\varphi}: T\{V\} \to A$ such that $\tilde{\varphi}|_V = \varphi$.  

$\tilde{\varphi}$ is defined recursively starting with $\tilde{\varphi}(0) = \varphi(0)$. And if $\tilde{\varphi}$ is defined for $V^{\otimes i}$ $i < n$, for $u \in V^{\otimes n}$ we write $u = \sum_i u_{ij} \otimes w_{ik}$, $u_{ij} \in V^{\otimes i}$, $w_{ik} \in V^{\otimes k}, j+k=n$.  
we define $\tilde{\varphi}(u) = \sum_i \tilde{\varphi}(u_{ij}) \tilde{\varphi}(w_{ik})$.  

If we fix a basis $e_i$, $i \in I$ in $V$, we may consider $T\{V\}$ as non-associative non-commutative algebra in variables $X = \{e_i : i \in I\}$. To make parallel with $F \langle X \rangle$ denote $T\{V\}$ by $F\{X\}$ ($F\langle X \rangle$ associative free, $F(X)$ associative commutative).  

Now let $I_{Jord} = ideal_{F\{X\}} \langle (a^2 b, a), [a, b]: a, b \in F\{X\} \rangle$.  

Then a quotient algebra $F\{X\} / I_{Jord} = Jord [X]$ is the free Jordan algebra over $V$.  

Open problem: To determine basis for Jord[X] also for $SJ \langle X \rangle$, for $|X| > 3$. Find an algorithm to write each element of $Jord[X]$ as linear combination of elements of this basis.  

\begin{theorem}[Cohn] 
Homomorphic image of $SJ\langle x,y \rangle$ is special.
\end{theorem}

\begin{proof} 
Let $I \triangleleft SJ\langle x,y \rangle$ and $\hat{I} = \text{ideal}_{F\langle X \rangle} \< I\>$. We prove that $I = \hat{I} \cap SJ\langle x,y \rangle $. By Cohn criterion $J =SJ\langle x,y \rangle / I$ is special. We will show that $I$ does not contain Cohn elements.

Let $k \in \hat{I} \cap SJ \langle x, y \rangle$. Then $k = \sum_n u_n i_n v_n$, where $i_n \in I$, $v_n, u_n$ are monomials in $F\langle x, y \rangle$ and $k^* = k$. Then  
$k = \sum_n u_n i_n v_n = \sum_n v_n^* i_n u_n^* = \frac{1}{2} \sum_n u_n i_n v_n + v_n^* i_n u_n^*$.  

For $u, v \in F \langle x, y \rangle$, consider in $F \langle x, y, z \rangle$ the element $r = uzv + v^* zu^*$. We have $r^* = r$.  
Thus, $r \in H \langle x, y, z \rangle = SJ \langle x, y, z \rangle$. Then there exists $j(x, y, z)$ such that $r = j(x, y, z)$. Therefore $uzv + v^* zu^* = j(x, y, z)$ and consequently for every $i \in I$, $j(x, y, i) = ui v + v^* i u^* \in I$. Then all summands in $k$ belong to $I$.  
\end{proof}

\begin{theorem}[Shirshov-Cohn] 
Every Jordan algebra generated by two elements is special. 
\end{theorem}

Consider $\pi$: Jord $[X] \to SJ \langle X \rangle$ surjective homomorphism with $\pi(x_i) = x_i$ $\forall x_i \in X$.  
Observe that $f \in \operatorname{Ker} \pi$ if $f \neq 0$ in Jord $[X]$ but $f = 0$ for any special Jordan algebra. Non-zero elements of $\operatorname{Ker} \pi$ are called s-identities.  

It is known that $\operatorname{Ker} \pi \neq 0$ for $|X| \ge 3$.  
We know there exists $f(x, y, z) \in \operatorname{Jord} [x, y, z]$ with $\deg_x f = \deg_y f = 3$ and $\deg_z f = 2$ in $\operatorname{Ker} \pi$. An example of such identity was obtained by Glennie in 1966 using computer. The Glennie identity denoted by $G_8$ is a minimal example. (since any polynomial of deg $\le 7$ in 3 variables does not belong to ker $\pi$).  

\begin{remark}
Open problems related to s-identities.  

1. Describe generators of $\operatorname{Ker} \pi$ as T-ideal  

2. To decide whether $\operatorname{Ker} \pi$ is finitely generated.  
\end{remark}

Next we describe Glennie identity and show that Albert algebra does not satisfy $G_8$.  

Consider $f = [[x, y]^3, z] \in F \langle x, y, z \rangle$, then $f^* = f$ ($([x, y]^3)^* = - [[x, y]^3]$ and if $u^* = -u$ \& $a^* = a$ then $[u, a]^* = [u, a]$). Therefore there exists $j(x, y, z) \in SJ \langle x, y, z \rangle$ such that $[[x, y]^3, z] = j(x, y, z)$.  

Observe that in any associative algebra we have $[a, b^2] = 2b \cdot [a, b] $ ($ \cdot $ here is $ a \cdot b=\frac{1}{2}(ab+ba) $), then $j(x, y, z^2) = [[x, y]^3, z^2] = 2z \cdot [[x, y]^3 z] = 2z \cdot j(x, y, z)$. Define $Sh(x, y, z) = j(x, y, z^2) - 2z j(x, y, z)$. 

In $SJ \langle x, y, z \rangle$ we have $Sh(x, y, z) = 0$. $Sh(x, y, z)$ is a scalar multiple of $G_8$. To write it down one has to write $[[x, y]^3, z]$, $ [[x, y]^3, z^2] $ as Jordan element.  

Finally, $Sh(x, y, z)$ is not an identity for exceptional Jordan algebras.  
In $H_3(\mathbb{D})$ the elements $x = \begin{bmatrix} 0 & 1 & 0 \\ 1 & 0 & 0 \\ 0 & 0 & 0 \end{bmatrix}$ $y = \begin{bmatrix} 0 & 0 & 0 \\ 0 & 0 & 1\\ 0 & 1 & 0 \end{bmatrix}$ and $z = \begin{bmatrix} 0 & a & c \\ \bar{a} & 0 & \bar{b} \\ \bar{c} & \bar{b} & 0 \end{bmatrix}$, with $a, b, c \in \mathbb{D}$ with the property that $(a, b, c) \neq 0$.  

In this way we have $Sh(x, y, z) \neq 0$ and consequently $H_3(\mathbb{D})$ is not homomorphic image of special algebra.  

\begin{definition} 
 An element $f \in F\langle X \rangle$ is called tetrads eater if $\{abcf\} \in SJ\langle X \rangle$, for all $a, b, c \in SJ\langle X \rangle$.  
\end{definition}

We have seen that $SJ\langle X \rangle \subseteq H\langle X \rangle \subsetneq F\langle X \rangle$ and $H\langle X \rangle = \operatorname{alg}_{F\langle X \rangle}^{(+)} \<X \cup \{x_{i_1}x_{i_2}x_{i_3}x_{i_4}, x_{i_j} \in X\}$.  

Given $u, v \in F\langle X \rangle$, we write $u \equiv v$ if $u - v \in SJ\langle X \rangle$.  

We have $\{xyzt\} \equiv -\{yxzt\} \equiv \{yxtz\} \equiv ...$ for $x, y, z, t \in X$.  

Observe that an skew symmetry continues if we take elements from $SJ\langle X \rangle$.  

Next we construct the following tetrads. 

\begin{lemma} 
For any $x, y, z, t \in X$ we have $\{x^2 yzt\} \equiv x \circ \{xyzt\}$.  
\end{lemma}

\begin{proof} 
 For any $x, y, z, t \in X$ we have $4\{x^2 yzt\} \equiv x \circ x\{yzt\}$.  
$x \circ \{xyzt\} = x^2 yzt +xtzyx + xyztx +  tzyx^2 = \{x^2yzt\} + x\{tzy\}x $  
\end{proof}

\begin{remark}
The result of the above lemma works for $x, y, z, t \in SJ\langle X \rangle$.
\end{remark}

\begin{lemma} 
For all $x, y, z, t \in SJ\langle X \rangle$, we have $\{xyzt\} \equiv \frac{1}{4} [x, y] \circ [z, t]$. 
\end{lemma}

\begin{proof} 
We have $\{xyzt\} \equiv -\{yxzt\} \equiv \frac{1}{2} \{[x, y] z t\} \equiv -\frac{1}{2}\{[x, y] tz\}= \frac{1}{4}\{[x, y] [z,t]\} = \frac{1}{4} ([x, y][z, t] + [z, t][x, y])= \frac{1}{4} [x, y] \circ [z, t]$.  
\end{proof}

\begin{corollary}
    $x, y, z \in SJ\langle X \rangle$ we have $[x, y] \circ [x, z] \equiv 0$. 
\end{corollary}

Denote by $(x, y, z)_+$ the associator of $x, y, z$ with operation $x \circ y$ in $SJ\langle X \rangle$.  
For any $a, b, c \in SJ\langle X \rangle$ we have $0 \equiv (a, b, c)_+ = (a \circ b) \circ c - a \circ (b \circ c) = [[c, a], b]$   

\begin{lemma} 
    For all $x, y, z, t \in SJ\langle X \rangle$ we have $\{[x, y]^2 xzt\} \equiv 0$.   
\end{lemma} 

\begin{proof} 
Let $v = [x, y]$. Then $\{[x, y]^2 xzt\} = -\{[x, y]^2 zxt\} \equiv -\frac{1}{4} [[x, y]^2, z] \circ [x, t] \equiv -\frac{1}{4} [v^2 z] \circ [x, t] = -\frac{1}{4} ([v, z] \circ v) \circ [x, t] = -\frac{1}{4} ([v, z] \circ v \circ[x, t]) + ([v,  z], v, [x, t])$.  

Then $[v, z] \equiv 0$ and $v \circ [x, t] \equiv 0$.  

Since $j = [v, z] \in SJ\langle X \rangle$ we have $([v, z], v, [x, t])_+ = (j, v, [x, t])_+ = [[[x, t], j], v] = [j_1, [x, y]] \equiv 0$. $j_1 \in SJ\langle X \rangle$.  
\end{proof}

\begin{lemma}
Let $a, b \in SJ\langle X \rangle$ such that $\{abxy\} \equiv 0$ for all $x, y \in SJ\langle X \rangle$. Then $[a, b]^2$ is a tetrad eater.
\end{lemma}

\begin{proof}
Linearize the identity $\{[x, y]^2 x z t\} \equiv 0$ in $x$ and then in $y$. We obtain
\[
\{[x, y]^2 u z t\} + \{([x, y] \circ [u, y]) x z t\} \equiv 0,
\]
and
\[
\{([x, y] \circ [x, u]) w z t\} + \{([x, y] \circ [w, u]) x z t\} + \{([w, y] \circ [x, u]) x z t\} \equiv 0.
\]
Using the hypothesis $\{abxy\} \equiv 0$ and substituting $a, b$ into these linearizations yields
$\{[a, b]^2 x z t\} \equiv 0$ for all $x, z, t$. Hence $[a, b]^2$ is a tetrad eater.
\end{proof}

\begin{corollary}
For all $x, y \in SJ\langle X \rangle$, the element $\left[[x, y]^2, x\right]^2$ is a tetrad eater.
\end{corollary}

\begin{proof}
Apply the previous lemma with $a = [x, y]^2$ and $b = x$. By the lemma above, $\{[x, y]^2 x z t\} \equiv 0$ for all $z, t$, so the hypothesis is satisfied.
\end{proof}

\begin{remark}
Tetrad eaters were introduced by E. Zelmanov, but this one is due to V. Skosyrski. This is the smallest degree known tetrad eater.
\end{remark}

\begin{definition}
Let
\[
N_t = \{ f \in SJ\langle X \rangle : \{fabc\} \in SJ\langle X \rangle \text{ for all } a, b, c \in SJ\langle X \rangle \}.
\]
We call $N_t$ the set of tetrad eaters of $SJ\langle X \rangle$.
\end{definition}

\begin{theorem}
Let $n \in N_t$, $n \neq 0$. Then $SJ\langle X \rangle / \operatorname{ideal}_{SJ\langle X \rangle}\langle n \rangle$ is exceptional.
\end{theorem}

\begin{proof}
Let $x, y, z \in X$ not appearing in $n$. Consider
\[
k = \{nxyz\} \in SJ\langle X \rangle \cap \hat{I}, \quad I = \operatorname{ideal}_{SJ\langle X \rangle}\langle n \rangle,
\]
where $\hat{I} = \operatorname{ideal}_{F\langle X \rangle}\langle I \rangle$.
We claim $k \notin I$. If $k \in I$, then $k = j(n, x, y, z)$ for some Jordan polynomial $j(t, x, y, z)$ in which every monomial involves $t$. Since $x, y, z$ do not occur in $n$, the only possible multilinear Jordan element is $\alpha \{txyz\}$, which is not Jordan. Hence $k \notin I$, so by the Cohn criterion the quotient is exceptional.
\end{proof}

Simple special Jordan algebras.

We will show that $J = H(A, *)$ for some associative simple algebra with involution $(A, *)$ or $J$ is a PI Jordan algebra.

Let $J$ be a Jordan algebra. For $a, b, c \in J$, denote $(a, x, b)_+$ by $x D_{a, b}$. We have that $D_{a, b}$ is a derivation of $J$.

\begin{lemma}
If $n \in N_t$, then $(a, n, b)_+ \in N_t$ for all $a, b \in SJ\langle X \rangle$.
\end{lemma}

\begin{proof}
Since $SJ\langle X \rangle \subseteq F\langle X \rangle$ and $a, b \in SJ\langle X \rangle$, in $F\langle X \rangle$ (with Jordan product) we have
\[
x D_{a, b} = (a, x, b)_+ = [x, [a, b]],
\]
so $D_{a, b}$ is a derivation. Hence
\[
\begin{aligned}
0 \equiv \{nxyz\} D_{a, b} &=
(n D_{a, b}) x y z + n(x D_{a, b}) y z + n x (y D_{a, b}) z + n x y (z D_{a, b}) \\
&\quad + (z D_{a, b}) y x n + z (y D_{a, b}) x n + z y (x D_{a, b}) n + z y x (n D_{a, b}).
\end{aligned}
\]
All terms except the first are in $SJ\langle X \rangle$, so $\{(n D_{a, b}) x y z\} \equiv 0$. Therefore $(a, n, b)_+ = n D_{a, b} \in N_t$.
\end{proof}

\begin{lemma}
For all $m, n \in N_t$ and all $a \in SJ\langle X \rangle$, the elements $n \circ m$ and $(n, a, m)_+$ belong to $N_t$.
\end{lemma}

\begin{proof}
Linearizing the identity $\{x^2 y z t\} \equiv x \circ \{x y z t\}$ gives
\[
\{(n \circ m) x y z\} = n \circ \{m x y z\} + m \circ \{n x y z\} \equiv 0,
\]
so $n \circ m \in N_t$. For $(n, a, m)_+$ we apply the same identity repeatedly:
\[
\{((a \circ n) \circ m) y z t\} \equiv m \circ \{(a \circ n) y z t\} + (a \circ n) \circ \{m y z t\} \equiv \{(a \circ n)(y \circ m) z t\} - y \circ \{(a \circ n) m z t\},
\]
and
\[
- y \circ \{(a \circ n) m z t\} \equiv a \circ n \circ \{y m z t\} + n \circ \{a y m z t\} \equiv \{((a \circ m) \circ n) y z t\}.
\]
Thus
\[
\{((a \circ n) \circ m - (a \circ m) \circ n) y z t\} \equiv 0,
\]
i.e. $(n, a, m)_+ \in N_t$.
\end{proof}

\begin{lemma}
Let $J$ be a Jordan algebra and let $V \subseteq J$ be a subspace such that $V D_{a, b} \subseteq V$ for all $a, b \in J$. Then
\[
\operatorname{id}_J\langle V \rangle = V + VJ + (VJ)J.
\]
\end{lemma}

\begin{proof}
It is enough to show it for $R_a R_b R_c$. We write $A \equiv 0$, $A \in \mathrm{Mult}\,J$ iff for all $v \in V$,
$
v A \in V + VJ + (VJ)J.
$
We have $R_a R_b R_c \equiv -R_c R_b R_a$ (LJI), therefore
\[
\begin{aligned}
2 R_a R_b R_c &= R_a R_b R_c - R_c R_b R_a \\
&= D_{a, b} R_c - D_{c, b} R_a + R_b D_{a, c}.
\end{aligned}
\]
Then $2 V R_a R_b R_c \subseteq V \cdot c + V \cdot a + (V \cdot b) D_{a, c}$. Moreover,
\[
(V \cdot b) D_{a, c} \subseteq (V D_{a, c}) \cdot b + V \cdot (b D_{a, c}) \subseteq V \cdot b + V \cdot (a, b, c)_+ \subseteq VJ.
\]
We conclude $V R_a R_b R_c \subseteq V + VJ + (VJ)J$.
\end{proof}

\begin{theorem}
For all $n_1, n_2, n_3 \in N_t$, let $m = (n_1, n_2, n_3)_+$. Then $\operatorname{id}_{SJ\langle X \rangle}\langle m \rangle \subseteq N_t$.
\end{theorem}

\begin{proof}
Let $V$ be the vector space $\langle (n_1, n_2, n_3)_+ : n_i \in N_t \rangle$. For all $a, b \in J = SJ\langle X \rangle$ we have
\[
\begin{aligned}
(n_1, n_2, n_3)_+ D_{a, b} &= (n_1 D_{a, b}, n_2, n_3)_+ + (n_1, n_2 D_{a, b}, n_3)_+ \\
&\quad + (n_1, n_2, n_3 D_{a, b})_+,
\end{aligned}
\]
which lies in $V$ by the previous lemma. Thus $V D_{a, b} \subseteq V$ for all $a, b \in J$. By the previous lemma it is enough to prove that $\operatorname{id}\langle V \rangle = V + VJ + (VJ)J \subseteq N_t$.

Let $a \in J$. Using the identity above we have
\[
(n_1, n_2, n_3)_+ \circ a = (n_1, n_2 \circ a, n_3)_+ - (n_1, a, n_3)_+ \circ n_2.
\]
By Lemma 2 the two elements on the right are in $N_t$, therefore $(n_1, n_2, n_3)_+ \circ a \in N_t$.

Let $a, b \in J$. We will show that $(m \circ a) \circ b \in N_t$. We have
\[
(m \circ a) \circ b = (m, a, b)_+ + m \circ (a \circ b),
\]
and by the previous item $m \circ (a \circ b) \in N_t$. Moreover,
\[
(m, a, b)_+ + (a, b, m)_+ + (b, m, a)_+ = 0
\]
(this identity is valid in any commutative algebra). By Lemma 1, $(b, m, a)_+ \in N_t$. Thus modulo $N_t$ we have (using flexibility)
\[
(m, a, b)_+ \equiv -(a, b, m)_+ \equiv (m, b, a)_+.
\]
The first and third terms are symmetric with respect to $a$ and $b$. Therefore it is enough to show that $(m, a, a)_+ \in N_t$. Observe that
\[
(m, a, a)_+ = (m \circ a) \circ a - m \circ a^2
\]
with $m \circ a^2 \in N_t$. Then
\[
\{(m, a, a)_+ y z t\} \equiv \{((m \circ a) \circ a) y z t\}.
\]
Using the linearization identity again,
\[
\{((m \circ a) \circ a) y z t\} \equiv (m \circ a) \circ \{a y z t\} + a \circ \{(m \circ a) y z t\} \equiv 0,
\]
since tetrads are antisymmetric. Hence $(m, a, a)_+ \in N_t$.
\end{proof}

\begin{corollary}
For all $x_1, \ldots, x_6 \in X$,
\[
\operatorname{id}_{SJ\langle X \rangle}\left\langle [[x_1, x_2]^2, x_1]^2,\ [[x_3, x_4]^2, x_3]^2,\ [[x_5, x_6]^2, x_5]^2 \right\rangle \subseteq N_t.
\]
\end{corollary}

\begin{remark}
Observe that $(n_1, n_2, n_3)_+$ has degree $30$ and $n_i$ has degree $10$. Elements in $N_t$ of degree $<10$ are not known, and no elements of degree $<30$ are known in the ideal of $N_t$.
\end{remark}

\begin{definition}
A Jordan algebra $J$ is called a PI Jordan algebra if it satisfies a polynomial identity $f=0$ with $0 \neq f \in \operatorname{Jord}[X]$.
\end{definition}

\begin{theorem}
Let $J$ be a special simple Jordan algebra. Then either $J = H(A, *)$ for some associative $*$-simple algebra $(A, *)$, or $J$ satisfies the polynomial identity
\[
f(x, y, z, t, u, v) = \big( [[x, y]^2, x]^2,\ [[z, t]^2, t]^2,\ [[u, v]^2, v]^2 \big)_+ = 0.
\]
In particular, $J \simeq H(A, *)$ or $J$ is a PI Jordan algebra.
\end{theorem}

\begin{proof}
Let
\[
P = \sum \{ I \triangleleft SJ\langle X \rangle : I \subseteq N_t \}.
\]
Define
\[
P(J) = \{ g(a_1, \ldots, a_n) : g \in P,\ a_i \in J \}.
\]
Then $P(J)$ is an ideal of $J$: if $a = g(a_1, \ldots, a_n) \in P(J)$ and $b \in J$, consider
$g_1(x_1, \ldots, x_{n+1}) = g(x_1, \ldots, x_n) \circ x_{n+1} \in P$, hence
$a \circ b = g_1(a_1, \ldots, a_n, b) \in P(J)$.
Since $J$ is simple, $P(J) = 0$ or $P(J) = J$.

If $P(J)=0$, then $J$ satisfies the identity $f=0$ above.

Assume $P(J)=J$. Since $J$ is special, we have $J = SJ\langle X \rangle / I$ with $I = \hat{I} \cap SJ\langle X \rangle$. Then
\[
J \simeq \frac{SJ\langle X \rangle}{\hat{I} \cap SJ\langle X \rangle}
\simeq \frac{SJ\langle X \rangle + \hat{I}}{\hat{I}}
\subseteq \frac{H\langle X \rangle + \hat{I}}{\hat{I}}
\subseteq \frac{F\langle X \rangle}{\hat{I}} \simeq S(J),
\]
where $S(J)$ is the special enveloping algebra.
The involution $*$ on $F\langle X \rangle$ induces an involution on $S(J)$ with $a^* = a$ for all $a \in J$.
Let $\bar{u}: F\langle X \rangle \to F\langle X \rangle/\hat{I} = S(J)$ be the canonical projection. Since
\[
SJ\langle X \rangle \subseteq H\langle X \rangle = H(F\langle X \rangle, *) \subseteq F\langle X \rangle,
\]
we obtain
\[
J \subseteq H(S(J), *) \subseteq S(J).
\]
Since $P(J)=J$, we have $\{JJJJ\} \subseteq J$, and by the Cohn theorem
\[
H(S(J), *) = \operatorname{alg}_{F\langle X \rangle}^{(+)} \langle J \cup \{JJJJ\} \rangle = J.
\]

In general, $S(J)$ is not simple. To find a simple $A$ with $H(A, *) = J$, let
\[
M = \{ I \triangleleft S(J) : I \cap J = 0 \}.
\]
By Zorn's lemma, pick a maximal $K \in M$ and set $A = S(J)/K$. Then $A$ is an algebra with involution and $J \hookrightarrow A$ since $J \cap K = 0$.
Let $\varphi: S(J) \to A$ be the canonical projection. Then $\varphi(H(S(J), *)) = H(A, *)$:
if $b^* = b$ and $b=\varphi(a)$, then $a-a^* \in \ker \varphi = K$ and
\[
b = \tfrac{1}{2}(b+b^*) = \varphi\!\left(\tfrac{1}{2}(a+a^*)\right) \in \varphi(H(S(J), *)).
\]
The reverse inclusion is clear. Hence $J = H(A, *)$.
Finally, $A$ is $*$-simple: if $0 \neq I \triangleleft A$ with $I^* = I$, then $\varphi^{-1}(I)$ is an ideal of $S(J)$.
Since $J$ is simple, $I \cap J = 0$ or $J$. By maximality of $K$, we must have $I \cap J = J$, hence $J \subseteq I$ and $A=I$.
\end{proof}

Radicals in Jordan algebras.

\begin{definition}
Let $\mathcal{V}$ be a class of algebras. A subclass $\mathcal{R} \subseteq \mathcal{V}$ is called a radical class if
\begin{enumerate}
    \item $\mathcal{R}$ is closed under homomorphisms;
    \item for every $A \in \mathcal{V}$ there exists $R(A) \triangleleft A$ such that $R(A) \in \mathcal{R}$ and for every $I \triangleleft A$ with $I \in \mathcal{R}$ we have $I \subseteq R(A)$;
    \item $R(A / R(A))=0$.
\end{enumerate}
The ideal $R(A)$ is called the $\mathcal{R}$-radical. If $R(A)=0$, then $A$ is called $\mathcal{R}$-semisimple.
\end{definition}

Radicals were introduced by Kurosh and Amitsur. Examples:
\begin{enumerate}
    \item For the class of associative algebras (or Jordan algebras of finite dimension) the subclass of nilpotent algebras is radical. In general, nilpotent algebras do not form a radical subclass (the sum of nilpotent ideals need not be nilpotent).
    \item For Jordan algebras of finite dimension the subclass of solvable algebras is radical.
    \item For the class of power-associative algebras the subclass of nil algebras is radical.
    \item For associative algebras of arbitrary dimension there are the following radical classes: Jacobson radical, prime radical, and locally nilpotent radical.
\end{enumerate}

An algebra is called locally nilpotent if any finitely generated subalgebra of $A$ is nilpotent.

\begin{proposition}
Let $\mathcal{R} \subseteq \mathcal{V}$ be a radical class and $A \in \mathcal{V}$. If $I \triangleleft A$ is such that $I \in \mathcal{R}$ and $A/I \in \mathcal{R}$, then $A \in \mathcal{R}$.
\end{proposition}

\begin{proof}
Consider $\bar{A} = A / R(A)$. By (ii) we have $I \subseteq R(A)$, hence $\bar{A}$ is a homomorphic image of $A/I$. Since $\mathcal{R}$ is closed under homomorphisms, $\bar{A} \in \mathcal{R}$, so $R(\bar{A}) = \bar{A}$. By (iii), $R(\bar{A})=0$, hence $\bar{A}=0$ and $A=R(A) \in \mathcal{R}$.
\end{proof}

\begin{definition}
A class of algebras $\mathcal{R}$ is called locally defined if
\[
A \in \mathcal{R} \Longleftrightarrow \text{every finitely generated subalgebra } B \subseteq A \text{ belongs to } \mathcal{R}.
\]
\end{definition}

Examples of locally defined classes: being nil, being algebraic, being locally nilpotent, and locally solvable.

\begin{lemma}
Let $\mathcal{V}$ be a class of algebras and let $\mathcal{R} \subseteq \mathcal{V}$ be a subclass satisfying the condition of the proposition. If $I_1, I_2 \triangleleft A$ and $I_1, I_2 \in \mathcal{R}$, then $I_1 + I_2 \in \mathcal{R}$.
\end{lemma}

\begin{proof}
We have $I_1 \in \mathcal{R}$ and $(I_1+I_2)/I_1 \simeq I_2/(I_1 \cap I_2) \in \mathcal{R}$. By the proposition, $I_1+I_2 \in \mathcal{R}$.
\end{proof}

\begin{corollary}
Let $\mathcal{V}$ be a class of algebras and let $\mathcal{R} \subseteq \mathcal{V}$ satisfy the condition of the proposition. If $I_1, \ldots, I_n \triangleleft A$ and $I_1, \ldots, I_n \in \mathcal{R}$, then $I_1+\cdots+I_n \in \mathcal{R}$.
\end{corollary}

\begin{proposition}
If $\mathcal{V}$ is a class of algebras and $\mathcal{R} \subseteq \mathcal{V}$ is closed under homomorphisms, locally defined, and satisfies the condition of the proposition, then $\mathcal{R}$ is radical.
\end{proposition}

\begin{proof}
Let $A \in \mathcal{V}$ and define
\[
R(A)=\sum\{I \triangleleft A : I \in \mathcal{R}\}.
\]
We show that $R(A) \in \mathcal{R}$ and $R(A / R(A))=0$. Let $a_1, \ldots, a_k \in R(A)$. Then there exist ideals $I_1, \ldots, I_k \triangleleft A$ with $I_j \in \mathcal{R}$ such that $a_j \in I_j$. By the corollary, $I_1+\cdots+I_k \in \mathcal{R}$. Thus the subalgebra $\langle a_1, \ldots, a_k\rangle$ lies in $\mathcal{R}$, and since $\mathcal{R}$ is locally defined, $R(A) \in \mathcal{R}$.

Now let $\bar{I} \triangleleft A/R(A)$ with $\bar{I} \in \mathcal{R}$. Let $I \triangleleft A$ be such that $R(A) \subseteq I$ and $I/R(A)=\bar{I}$. Since both $I$ and $R(A)$ are in $\mathcal{R}$, we have $I \in \mathcal{R}$ and hence $I \subseteq R(A)$. Thus $\bar{I}=0$, so $R(A/R(A))=0$.
\end{proof}

Locally nilpotent radical.

\begin{lemma}
Let $J$ be a Jordan algebra, $J=Fy+I$, with $J^2 \subseteq I$. If $I$ is nilpotent, then $J$ is nilpotent.
\end{lemma}

\begin{proof}
As proved earlier, for every $m \ge 1$ we have
\[
U(J)^{3m} \subseteq U_J(I)^m + U_J(I)^m U(J),
\]
where
\[
U_J(I)=\operatorname{alg}_{U(J)}\langle R_a : a \in I \rangle.
\]
Let $\pi:U(J)\to \operatorname{Mult}(J)$ be the canonical homomorphism. Set
\[
R_J(I)=\pi(U_J(I))=\operatorname{alg}_{\operatorname{Mult}(J)}\langle R_a : a \in I \rangle.
\]
Applying $\pi$ to the previous inclusion, we get
\[
\operatorname{Mult}(J)^{3m} \subseteq R_J(I)^m + R_J(I)^m\operatorname{Mult}(J).
\]
If $I^m=0$, then $R_J(I)^m=0$ (indeed, for $i_1,\ldots,i_m\in I$ and $a\in J$,
$aR_{i_1}\cdots R_{i_m}\in I^m=0$). Hence $\operatorname{Mult}(J)$ is nilpotent, so $J$ is nilpotent.
\end{proof}

\begin{lemma}
Let $J$ be a finitely generated Jordan algebra. If $J^2$ is nilpotent, then $J$ is nilpotent.
\end{lemma}

\begin{proof}
Write $J=\operatorname{alg}\langle a_1,\ldots,a_m\rangle$. Then
\[
J=Fa_1+\cdots+Fa_m+J^2.
\]
Define
\[
I_0=J^2,\quad I_1=J^2+Fa_1,\quad I_2=I_1+Fa_2,\ \ldots,\ I_m=I_{m-1}+Fa_m.
\]
Each $I_k$ is an ideal (it is a subspace containing $J^2$), and
\[
(I_{k+1})^2\subseteq I_k.
\]
Since $I_0=J^2$ is nilpotent, repeated use of the previous lemma gives nilpotency of $I_1,\ldots,I_m$. Therefore $I_m=J$ is nilpotent.
\end{proof}

\begin{lemma}
If $J$ is a finitely generated Jordan algebra, then $J^2$ is finitely generated.
\end{lemma}

\begin{proof}
We may assume $J$ is a free Jordan algebra. Let
\[
I=(J^2)^3\triangleleft J,
\]
and set $\bar{J}=J/I$. Then
\[
(\bar{J}^{\,2})^3=0.
\]
Since $\bar{J}$ is finitely generated, the previous lemma implies that $\bar{J}$ is nilpotent. Hence $\bar{J}^{\,m}=0$ for some $m$, so
\[
J^m\subseteq I\subseteq (J^2)^2.
\]
Let $S$ be the set of all monomials in the free generators of $J$ of degrees $2,\ldots,m$. The set $S$ is finite. We claim
\[
J^2=\operatorname{alg}\langle S\rangle.
\]
The inclusion $\operatorname{alg}\langle S\rangle\subseteq J^2$ is clear. For the converse, let $v\in J^2$ be a monomial. If $\deg(v)\le m$, then $v\in S$. If $\deg(v)>m$, then $v\in J^m\subseteq (J^2)^2$, so
\[
v=\sum_i u_i\circ v_i,
\]
with $u_i,v_i\in J^2$ and $\deg(u_i),\deg(v_i)<\deg(v)$. By induction on degree, $u_i,v_i\in\operatorname{alg}\langle S\rangle$, hence $v\in\operatorname{alg}\langle S\rangle$.
\end{proof}

\begin{corollary}
If $J$ is a finitely generated Jordan algebra, then $J^{(m)}$ is finitely generated for every $m\ge 1$.
\end{corollary}

\begin{proof}
Use induction on $m$: $J^{(1)}=J^2$ is finitely generated by the lemma, and
\[
J^{(m+1)}=(J^{(m)})^2.
\]
\end{proof}

\begin{exercise}
Show that if $J$ is a finitely generated Jordan algebra, then $J^m$ is finitely generated for every $m\ge 1$.
\end{exercise}

\begin{lemma}
Every finitely generated solvable Jordan algebra is nilpotent.
\end{lemma}

\begin{proof}
Induct on the solvability index. Assume $J^{(m)}=0$ and the statement is true for solvability index $<m$. Then $J^2$ is finitely generated, and
\[
(J^2)^{(m-1)}\subseteq J^{(m)}=0.
\]
So, by induction, $J^2$ is nilpotent. The second lemma above then implies that $J$ is nilpotent.
\end{proof}

\begin{remark}
The finite generation hypothesis in the previous lemma is necessary (see the example below).
\end{remark}

\begin{lemma}
Let $J$ be a Jordan algebra and $I\triangleleft J$. If both $I$ and $J/I$ are locally nilpotent, then $J$ is locally nilpotent.
\end{lemma}

\begin{proof}
Let $A\subseteq J$ be a finitely generated subalgebra. Then
\[
\bar{A}:=A/(A\cap I)\simeq (A+I)/I\subseteq J/I.
\]
Since $J/I$ is locally nilpotent, $\bar{A}$ is nilpotent; hence $A^m\subseteq I$ for some $m$. Choose $n$ with $2^n\ge m$. Then
\[
A^{(n)}\subseteq A^{2^n}\subseteq A^m\subseteq I.
\]
By the corollary, $A^{(n)}$ is finitely generated; since $I$ is locally nilpotent, $A^{(n)}$ is nilpotent, therefore solvable. Thus $A^{(n+k)}=0$ for some $k$, so $A$ is solvable. By the previous lemma, $A$ is nilpotent. Therefore every finitely generated subalgebra of $J$ is nilpotent, i.e. $J$ is locally nilpotent.
\end{proof}

\begin{theorem}[Zhevlakov]
The property of local nilpotency is radical in the class of Jordan algebras.
\end{theorem}

Denote by $\operatorname{LocNilp}(J)$ the locally nilpotent radical of $J$. Then
\[
\operatorname{LocNilp}(J)\subseteq \operatorname{Nil}(J).
\]

\begin{exercise}
Let $A$ be a finitely generated associative algebra. Show that $A^{(+)}$ is finitely generated as a Jordan algebra.
\end{exercise}

\begin{exercise}
Let $A$ be an associative algebra such that the Jordan algebra $A^{(+)}$ is nilpotent. Show that $A$ is nilpotent.
\end{exercise}

\begin{remark}
There exists a finitely generated associative algebra $A$ such that $A=\operatorname{Nil}(A)$ and $\operatorname{LocNilp}(A)=0$ (Golod--Shafarevich). Then $J=A^{(+)}$ is a nil Jordan algebra, finitely generated but not nilpotent.
\end{remark}

\begin{remark}
Shirshov proved that if $J$ is special, then
\[
\operatorname{LocNilp}(J)=\operatorname{LocNilp}(S(J))\cap J.
\]
\end{remark}

\begin{example}
There exists a Jordan algebra that is solvable but not nilpotent.

Let $V$ be a vector space over $F$ with basis $\{v_i:i\ge1\}$, and let
\[
G(V)=\operatorname{alg}\langle \bar v_i:i\ge1\mid \bar v_i\bar v_j=-\bar v_j\bar v_i\rangle
\]
be the Grassmann algebra without unit. Set $J=V\oplus G(V)$ (direct sum of vector spaces) and define a commutative product $\circ$ by
\[
V\circ V=0,\qquad G(V)\circ G(V)=0,\qquad v_i\circ w=w\circ v_i=w\bar v_i
\]
for all $i\ge1$ and $w\in G(V)$ (here $w\bar v_i$ is the product in $G(V)$).

Then $J$ is a Jordan algebra. Indeed, for $d=w\in G(V)$ and basis elements $a,b,c\in V$,
\[
((w\circ a)\circ b)\circ c+((w\circ c)\circ b)\circ a
= w(\bar a\bar b\bar c+\bar c\bar b\bar a)=0.
\]
Also,
\[
J^2\subseteq G(V),\qquad G(V)\circ G(V)=0,
\]
so $(J^2)^2=0$ and $J$ is solvable.

On the other hand, $J$ is not nilpotent, since for every $n\ge2$,
\[
((\cdots((\bar v_1\circ v_2)\circ v_3)\circ\cdots)\circ v_n)=\bar v_1\bar v_2\cdots\bar v_n\neq0.
\]
\end{example}

\begin{definition}
An algebra $A$ is called a \textbf{prime algebra} if for any ideals $I, J$ of $A$ we have
\[
IJ = 0 \Rightarrow I = 0 \text{ or } J = 0.
\]
An ideal $P$ of $A$ is called \textbf{prime} if $A/P$ is a prime algebra. Thus $P$ is a prime ideal if for any two ideals $I, J$ of $A$ we have
\[
IJ \subseteq P \Rightarrow I \subseteq P \text{ or } J \subseteq P.
\]
\end{definition}

Let $A$ be an algebra. Denote by $\operatorname{rad}(A)$ the intersection of all prime ideals of $A$.

\begin{definition}
A sequence $\{a_i : a_i \in A, i \geq 1\}$ is an \textbf{$m$-sequence} if for any $i \geq 1$ we have $a_{i+1} \in (\operatorname{id}_A(a_i))^2$. We say that an element $a \in A$ \textbf{annihilates $m$-sequences} if any $m$-sequence which contains $a$ also contains zero.
\end{definition}

\begin{definition}
An algebra $A$ is called a \textbf{semi-prime algebra} if for any ideal $I$ of $A$ with $I^2 = 0$ we have $I = 0$.

An ideal $P$ of $A$ is called a \textbf{semi-prime ideal} if $A/P$ is a semi-prime algebra. Then $P$ is a semi-prime ideal if for any ideal $I$ of $A$ with $I^2 \subseteq P$ we have $I \subseteq P$.
\end{definition}

\begin{lemma}
If $a \in A$ is an element which annihilates $m$-sequences, then $a$ belongs to all semi-prime ideals of $A$.
\end{lemma}

\begin{proof}
Suppose there exists a semi-prime ideal $P$ of $A$ such that $a \notin P$. Consider $(\operatorname{id}\langle a\rangle)^2$. If $(\operatorname{id}\langle a\rangle)^2 \subseteq P$, then $\operatorname{id}\langle a\rangle \subseteq P$ and thus $a \in P$. Hence, there exists $a_1 \in (\operatorname{id}\langle a\rangle)^2$ with $a_1 \notin P$. In a similar manner, there exists $a_2 \in (\operatorname{id}\langle a_1\rangle)^2$ such that $a_2 \notin P$. Then we can construct an infinite $m$-sequence containing $a$: $\{a, a_1, a_2, \dots\}$. Thus $a$ does not annihilate $m$-sequences, contradicting the hypothesis.
\end{proof}

\begin{theorem}
Let $A$ be an algebra. Then $\operatorname{rad}(A) = \{a \in A : a \text{ annihilates } m\text{-sequences}\}$.
\end{theorem}

\begin{proof}
Let $M = \{a \in A : a \text{ annihilates } m\text{-sequences}\}$. To show that $\operatorname{rad}(A) \subseteq M$, it is enough to prove that if $a \notin M$, then there exists a prime ideal $P$ of $A$ such that $a \notin P$. Let $m = \{a = a_0, a_1, \dots, a_{n}, \dots\}$ be an $m$-sequence with $0 \notin m$. Consider
\[
\mathcal{M} = \{I \triangleleft A : m \cap I = \varnothing\}.
\]
Then $\{0\} \in \mathcal{M}$ and $\mathcal{M}$ is inductive: if $\{I_1 \subseteq I_2 \subseteq I_3 \subseteq \dots\}$ is a chain of ideals in $\mathcal{M}$, then $I = \bigcup_{k \geq 1} I_k \in \mathcal{M}$. By Zorn's lemma, there exists $P \in \mathcal{M}$ a maximal element. It is clear that $a \notin P$. We show that $P$ is a prime ideal. Let $I, J$ be ideals of $A$ such that $IJ \subseteq P$, $I \nsubseteq P$ and $J \nsubseteq P$. We may suppose without loss of generality that $P \cap I = P \cap J = \varnothing$ (otherwise replace $I$ by $I + P$ and $J$ by $J + P$). Then there exist $a_i \in I$ and $a_j \in J$, since $P$ is maximal. Suppose that $i \leq j$, then $a_j \in I$ (by induction, since $a_{i+1} \in (\operatorname{id}\langle a_i\rangle)^2 \subseteq I$). Therefore
\[
a_{j+1} \in (\operatorname{id}\langle a_j\rangle)^2 \subseteq IJ \subseteq P,
\]
which implies that $a_{j+1} \in P \cap m$, contradicting the fact that $P \in \mathcal{M}$. Thus $P$ is a prime ideal.

The other inclusion is a consequence of the previous lemma, since any prime ideal is semi-prime.
\end{proof}

\begin{corollary}
The ideal $\operatorname{rad}(A)$ is the minimal semi-prime ideal. In particular, if $A$ is semi-prime, then $\operatorname{rad}(A) = 0$.
\end{corollary}

\begin{proof}
We show that $A/\operatorname{rad}(A)$ is semi-prime. Let $I \triangleleft A$ such that $I^2 \subseteq \operatorname{rad}(A)$, or equivalently for each prime ideal $P$ of $A$, we have $I^2 \subseteq P$. This implies that $I \subseteq P$, since $P$ is prime. Therefore $I \subseteq \operatorname{rad}(A)$ and consequently $A/\operatorname{rad}(A)$ is a semi-prime algebra.
\end{proof}

\begin{corollary}
For any algebra $A$ we have $\operatorname{rad}(A/\operatorname{rad}(A)) = 0$.
\end{corollary}

\begin{proof}
Note that for any prime ideal $\bar{P}$ of $\bar{A} = A/\operatorname{rad}(A)$, there exists a prime ideal $P \triangleleft A$ with $\operatorname{rad}(A) \subseteq P$ and $\bar{P} = P/\operatorname{rad}(A)$. Conversely, for each prime ideal $P \triangleleft A$, the ideal $\bar{P} = P/\operatorname{rad}(A)$ is a prime ideal of $\bar{A}$. Then for any $\bar{a} \in \bigcap \{\bar{P} \triangleleft \bar{A}\}$, we have that if $a \in A$ is such that $\bar{a} = a + \operatorname{rad}(A)$, then $a \in P$ for each prime ideal $P$ of $A$. Thus
\[
a \in \bigcap\{P \triangleleft A\} = \operatorname{rad}(A) \Rightarrow \bar{a} = \bar{0}.
\]
Consequently, $\operatorname{rad}(\bar{A}) = \bigcap \{\bar{P} \triangleleft \bar{A}\} = 0$.
\end{proof}

\begin{definition}
We say that an algebra $A$ is a subdirect product of algebras $\{A_\alpha : \alpha \in I\}$ if $A$ is isomorphic to a subalgebra $B \subseteq \prod_{\alpha \in I} A_\alpha$ with the property that for any $\alpha \in I$ the projection $\pi_\alpha : B \to A_\alpha$ is surjective.
We write $A \simeq \prod_{\alpha \in I} A_\alpha$ in the case when $A$ is a subdirect product of $\{A_\alpha : \alpha \in I\}$.
\end{definition}

We have the following characterization of subdirect product.

\begin{proposition}
An algebra $A$ is isomorphic to a subdirect product of algebras $\{A_\alpha : \alpha \in I\}$ if and only if there exists a family of ideals $\{K_\alpha \triangleleft A : \alpha \in I\}$ such that
\begin{enumerate}
\item[(i)] for any $\alpha \in I$ we have $A / K_\alpha \simeq A_\alpha$,
\item[(ii)] $\bigcap_{\alpha \in I} K_\alpha = 0$.
\end{enumerate}
\end{proposition}

\begin{proof}
If $A \simeq \prod_{\alpha \in I} A_\alpha$ then $A \subseteq \prod_{\alpha \in I} A_\alpha$. Let $K_\alpha = \ker \pi_\alpha \cap A$, with $\pi_\alpha : \prod A_\alpha \to A_\alpha$ the projection. We have $K_\alpha \triangleleft A$. If $a \in \bigcap K_\alpha$ then $\pi_\alpha(a)=0$ for any $\alpha$, hence $a=0$. Further,
\[
A / K_\alpha = A / (A \cap \ker \pi_\alpha) \simeq (A + \ker \pi_\alpha)/\ker \pi_\alpha \simeq \pi_\alpha(A) = A_\alpha.
\]

Conversely, if there exists $\{K_\alpha \triangleleft A : \alpha \in I\}$ with $\bigcap K_\alpha = 0$, we denote by $A_\alpha$ the quotient $A / K_\alpha$ and consider $\varphi : A \to \prod_{\alpha \in I} A_\alpha$ defined by $\varphi(a)_\alpha = a + K_\alpha \in A_\alpha = A / K_\alpha$ for all $\alpha \in I$. We have $\varphi$ is a homomorphism, $\ker \varphi = \bigcap K_\alpha = 0$, so $A$ embeds into $\prod_{\alpha \in I} A_\alpha$, and for any $\alpha \in I$, $\pi_\alpha(\varphi(A)) = A_\alpha$.
\end{proof}

\begin{corollary}
An algebra $A$ is semiprime if and only if $A$ is isomorphic to a subdirect product of prime algebras.
\end{corollary}

\begin{proof}
If $A$ is semiprime, then $0 = \operatorname{rad} A = \bigcap \{P : P \text{ is a prime ideal of } A\}$. By the proposition, $A$ is a subdirect product of $\{A / P : P \text{ prime ideal of } A\}$, and each $A / P$ is a prime algebra.

Conversely, if $A \simeq \prod_{\alpha \in I} A_\alpha$ with $A_\alpha$ prime for all $\alpha \in I$, by the proposition there exists a family $\{P_\alpha \triangleleft A : \alpha \in I\}$ such that $A / P_\alpha = A_\alpha$ for all $\alpha$ and $\bigcap_{\alpha \in I} P_\alpha = 0$. Then $\operatorname{rad} A = \bigcap \{P : P \text{ is prime ideal of } A\} \subseteq \bigcap_{\alpha \in I} P_\alpha = 0$, so $A$ is semiprime.
\end{proof}

\begin{theorem}
Let $\mathcal{V}$ be a homogeneous variety of algebras with the property that locally nilpotent algebras in $\mathcal{V}$ are radical. Then for any algebra $A \in \mathcal{V}$ we have $\operatorname{LocNilp}(A) = \bigcap \{P \triangleleft A : P \text{ prime}, \operatorname{LocNilp}(A/P)=0\}$.
\end{theorem}

\begin{proof}
Let $L(A)=\bigcap\{P\triangleleft A : P \text{ prime}, \operatorname{LocNilp}(A/P)=0\}$.
We have $L(A)\supseteq \operatorname{LocNilp}(A)$ since if $\operatorname{LocNilp}(A/P)=0$, then $P\supseteq \operatorname{LocNilp}(A)$.
For the other inclusion, let $a\in A\setminus \operatorname{LocNilp}(A)$.
We will show that there exists a prime ideal $P$ of $A$ such that $\operatorname{LocNilp}(A/P)=0$ and $a\notin P$.

Since $a\notin \operatorname{LocNilp}(A)$, the ideal $\langle a\rangle$ is not locally nilpotent.
Then there exists a finitely generated subalgebra $S\subseteq \langle a\rangle$ which is not nilpotent.
Consider
\[
M=\{I\triangleleft A : a\notin I \text{ and } S^n\nsubseteq I, \forall n\geq 1\}.
\]
$M$ is not empty, since $0\in M$.
Moreover $M$ is inductive: if $I_1\subseteq I_2\subseteq \cdots \subseteq I_n\subseteq \cdots$ is a chain in $M$, we show that $I=\bigcup I_n\in M$ (and $I$ is the supremum for the given chain).
If $I\notin M$, then there exists $n\geq 1$ such that $S^n\subseteq I$.
Let
\[
S_0=\{v : v \text{ is a product of generators of } S \text{ with } n\leq d(v)\leq 2n\},
\]
where $d(v)$ denotes the number of generators of $S$ in $v$.
Since $S_0\subseteq I$, there exists $m\geq 1$ with $S_0\subseteq I_m$.
We prove that $S^n\subseteq I_m$.

Let $u\in S^n$ be a product of generators with $d(u)\geq n$.
We prove that $u\in I_m$ using induction on $d(u)$.
If $d(u)\leq 2n$, we have $u\in S_0\subseteq I_m$.
If $d(u)>2n$, then $u=\sum_i v_i w_i$ with $d(v_i)+d(w_i)=d(u)>2n$.
We have that $v_i$ (or $w_i$) satisfies $n<d(v_i)$ really.
By induction $v_i\in I_m$, the same holds for every part of $u$, we conclude $u\in I_m$.
Then $S^n\subseteq I_m$, which is impossible since $I_m\in M$.
Thus $M$ is inductive.
By Zorn's lemma, there exists a maximal ideal $P\in M$.
Since $\langle a\rangle\nsubseteq P$, we have $a\notin P$.
We prove that $P$ is a prime ideal.

If there exist ideals $I$ and $J$ of $A$ with $P\subsetneq I$, $P\subsetneq J$ and $IJ\subseteq P$, then there exist $m,n\geq 1$ with $S^m\subseteq I$ and $S^n\subseteq J$.
Then $S^m S^n\subseteq IJ\subseteq P$, so $(S^n)^2\subseteq P$.
Consider $\bar{A}=A/P$, then $\bar{S}=S+P/P\subseteq \bar{A}$ is not nilpotent but $(\bar{S}^n)^2\subseteq P/P=0$.
We have $B=\bar{S}^n\Delta \bar{S}$, $B^2=0$ and $(\bar{S}/B)^n=0$.
Since local nilpotency is a radical, we conclude that $\bar{S}$ is locally nilpotent and therefore nilpotent, a contradiction.
\end{proof}

\begin{theorem}
Let $A$ be an algebra, then
\[
\widetilde{\operatorname{Nil}}(A)=\bigcap\{P\triangleleft A : P \text{ prime}, \widetilde{\operatorname{Nil}}(A/P)=0\}.
\]
\end{theorem}

\begin{proof}
By the previous theorem and what we already know, $\operatorname{rad}(A)\supseteq \operatorname{LocNilp}(A)\supseteq \operatorname{Nil}(A)$.
\end{proof}

\begin{remark}
The first inclusion is also strict: there exist prime locally nilpotent Jordan algebras.
Example comes again from the cone of prime locally nilpotent associative algebras.
\end{remark}

\subsection*{Semiprime Jordan algebras}
We show that homomorphic images of $H\langle X\rangle$ are special Jordan algebras.
First
\[
\operatorname{Sym}\langle X\rangle \subseteq H\langle X\rangle \subseteq F\langle X\rangle.
\]

\begin{lemma}
Let $I\triangleleft H\langle X\rangle$.
Then for every $a\notin I$ we have $\operatorname{id}_{F\langle X\rangle}\langle a^2\rangle \cap H\langle X\rangle \subseteq I$.
\end{lemma}

\begin{proof}
Let $f\in \operatorname{id}_{F\langle X\rangle}\langle a^2\rangle \cap H\langle X\rangle$.
\end{proof}

\begin{proof}
Let $f \in \operatorname{id}_{F\langle x\rangle}\langle a^{2}\rangle \cap H\langle x\rangle$. Then
\[
f = \sum_i \{u_i a^{2} v_i\}, \quad \text{where } u_i \text{ and } v_i \text{ are monomials in } x.
\]
Since $H\langle x\rangle = \operatorname{alg}_{F\langle x\rangle} + \langle x \cup \{x_i x_j x_k x_\ell\}\rangle$, define
\[
\tilde{f}(t) = \sum_i \{u_i t v_i\} \in H\langle x \cup \{t\}\rangle.
\]
Then
\begin{align*}
\tilde{f}(t) &= j_1(x_1,\dots,x_n,\{x_i x_j x_k x_\ell\},\dots, t) \\
&\quad + j_2(x_1,\dots,x_n,\{x_i x_j x_k x_\ell\},\dots,\{x_i x_j x_k t\}).
\end{align*}
Now $f = \tilde{f}(a^{2})$. Substituting $t$ by $a^{2}$ in $j_1$ we have
\[
j_1(x_1,\dots,x_n,\{x_i x_j x_k x_\ell\},\dots, a^{2}) \in I,
\]
since $a^{2} \in I$. And since
\[
\{x_i x_j x_k a^{2}\} = j(x_i, x_j, x_k, a^{2}) + a \cdot \{x_i, x_j, a\},
\]
where $j(x_i, x_j, x_k, a^{2}) \in I$ and $a \cdot \{x_i, x_j, a\} \in I$, substituting $t$ by $a^{2}$ in $j_2$ we get
\[
j_2(x_1,\dots,x_n,\dots, i_1 + i_2) \in I.
\]
Thus $f \in I$.
\end{proof}

\begin{corollary}
$\operatorname{id}_{F\langle x\rangle}\langle a^{2}\rangle \cap H\langle x\rangle \leq I$.
\end{corollary}

\begin{proof}
$I^{2} = \{\sum_i a_i^{2} : a_i \in I\}$.
\end{proof}

\begin{lemma}
Let $J$ be a special Jordan algebra and let $A$ be an associative algebra such that $J \subseteq A^{(+)}$. Suppose that $A = \operatorname{alg}\langle J\rangle$. If $I$ is an ideal of $J$ with $I^{2}=0$, then $I \leq \operatorname{rad} A$.
\end{lemma}

\begin{proof}
Consider $\bar{A} = A/\operatorname{rad} A$ and $\bar{J} = (J + \operatorname{rad} A)/\operatorname{rad} A$. Then we have
\[
\bar{J} \leq \bar{A}^{(+)}, \quad \bar{A} = \operatorname{alg}\langle \bar{J}\rangle.
\]
Without loss of generality we may assume $\operatorname{rad} A = 0$, and prove that $I = 0$.

Let $a, b \in I$, then $a \cdot b = 0$ in $A$. Consider $c = [a, b] \in A$ and $(x, y, z)_{+}$ the associator of $x, y, z$ with the operation $\cdot$. For any $x \in J$ we have
\begin{align*}
[x, c] &= [x, [a, b]] = 4(a, x, b)_{+} = (a \circ x) \cdot b - a \cdot (x \circ b) \in I^{2} = 0.
\end{align*}
Since $A = \operatorname{alg}\langle J\rangle$, we have $[c, A] = 0$, i.e., $c \in Z(A)$. From the identity in associative algebras we have
\[
[a, b] \cdot x = [a \cdot x, b] - a \cdot [x, b].
\]
Then
\begin{align*}
2c^{2} &= 2[a, b]^{2} = [a, b] \cdot [a, b] = [a \cdot [a, b], b] - a \cdot [[a, b], b] \\
&= [[a^{2}, b], b] - a \cdot [[a, b], b] \\
&= -4(a^{2}, b, b)_{+} + 8a \cdot (a, b, b)_{+} = 0.
\end{align*}
Since $c$ is central, we have $(cA)^{2} = c^{2}A^{2} = 0$. Thus $c = 0$ because $\operatorname{rad} A = 0$. Hence for all $a, b \in I$, $[a, b] = 0$. Since $a \circ b = 0$ implies $ab = 0$ for all $a, b \in I$.

Denote $\hat{I} = \operatorname{id}_{A}\langle I\rangle$. We prove that $\hat{I} = I + IA = I + AI$. To show it, it is enough to see that $I + IA$ and $I + AI$ are ideals of $A$. For any $x \in A$, $a \in I$ we have
\[
xa = (xa + ax) - ax \in I + IA.
\]
Thus $J(I + IA) \subseteq I + IA$. Since $J$ generates $A$, we have $A(I + IA) \subseteq I + IA$. Moreover, $I + IA$ is a right ideal, thus $I + IA$ is an ideal. Analogously we show that $I + AI$ is an ideal, hence
\[
(\hat{I})^{2} = (I + AI)(I + IA) \subseteq A I I A = 0.
\]
Therefore $\hat{I} = 0$, so $I = 0$.
\end{proof}

\begin{corollary}
If $J$ is a Jordan algebra, and $J$ is special with $A$ such that $J \subseteq A^{(+)}$ and $A = \operatorname{alg}\langle J\rangle$, then $\operatorname{rad} J = J \cap \operatorname{rad} A$.
\end{corollary}

\begin{proof}
Let $\bar{A}=A/\operatorname{rad}A$, then $\bar{J}\subseteq\bar{A}^{(+)}$ and $\bar{A}=\operatorname{alg}\langle\bar{J}\rangle$. Let $\bar{I}\triangleleft\bar{J}$ with $\bar{I}\cdot\bar{I}=0$. By the previous lemma, $\bar{I}\subseteq\operatorname{rad}\bar{A}=0$, so $\bar{J}$ is semiprime and $\operatorname{rad}\bar{J}=0$, thus $\operatorname{rad}J\leq\operatorname{rad}A$. Therefore $\operatorname{rad}J\leq\operatorname{rad}A$.

For the reverse inclusion, let $a\in\operatorname{rad}A\cap J$. Consider any $m$-sequence $\{a_i:i\geq1\}$ in $J$ containing $a$. Since $a_i\in J$, we have $\operatorname{id}_J\langle a_i\rangle\subseteq\operatorname{id}_A\langle a_i\rangle$, so
\[
a_{i+1}=(\operatorname{id}_J\langle a_i\rangle)^2\subseteq(\operatorname{id}_A\langle a_i\rangle)^2.
\]
Thus $\{a_i:i\geq1\}$ is an $m$-sequence in $A$. Since $a\in\operatorname{rad}A$, this $m$-sequence contains $0$. Hence $a\in\operatorname{rad}J$. This gives the other inclusion.
\end{proof}

\begin{theorem}
Let $J=H\langle X\rangle/P$. If $\operatorname{rad}J=0$, then $J$ is special.
\end{theorem}

\begin{proof}
Let $\widehat{P^2}=\operatorname{id}_{F\langle X\rangle}\langle P^2\rangle$. We have $\widehat{P^2}\cap H\langle X\rangle\leq P$ by the first corollary. Then
\[
\widehat{P^2}\cap P\leq\widehat{P^2}\cap H\langle X\rangle\subseteq\widehat{P^2}\cap P,
\]
so
\[
\widehat{P^2}\cap H\langle X\rangle=\widehat{P^2}\cap P.
\]
We have
\[
\frac{H\langle X\rangle+\widehat{P^2}}{\widehat{P^2}}\Big/\frac{P+\widehat{P^2}}{\widehat{P^2}}\simeq\frac{H\langle X\rangle}{H\langle X\rangle\cap\widehat{P^2}}\Big/\frac{P}{P\cap\widehat{P^2}}\simeq H\langle X\rangle/P=J
\]
as Jordan algebras.

Since $\operatorname{rad}J=0$, we have
\[
\operatorname{rad}\left(\frac{H\langle X\rangle+\widehat{P^2}}{\widehat{P^2}}\right)\subseteq\frac{P+\widehat{P^2}}{\widehat{P^2}}.
\]
Let $B=\operatorname{rad}\left(F\langle X\rangle/\widehat{P^2}\right)$. Then by the second corollary,
\begin{align*}
B\cap\frac{H\langle X\rangle+\widehat{P^2}}{\widehat{P^2}}&=\operatorname{rad}\left(\frac{F\langle X\rangle}{\widehat{P^2}}\right)\cap\frac{H\langle X\rangle+\widehat{P^2}}{\widehat{P^2}}\\
&=\operatorname{rad}\left(\frac{H\langle X\rangle+\widehat{P^2}}{\widehat{P^2}}\right)\leq\frac{P+\widehat{P^2}}{\widehat{P^2}}.
\end{align*}
On the other hand, since $\left(\frac{P+\widehat{P^2}}{\widehat{P^2}}\right)^2=0$, we have
\[
\frac{P+\widehat{P^2}}{\widehat{P^2}}\leq\operatorname{rad}\left(\frac{F\langle X\rangle}{\widehat{P^2}}\right)=B.
\]
Thus $\frac{P+\widehat{P^2}}{\widehat{P^2}}\subseteq B\cap\frac{H\langle X\rangle+\widehat{P^2}}{\widehat{P^2}}$, so equality holds:
\[
B\cap\frac{H\langle X\rangle+\widehat{P^2}}{\widehat{P^2}}=\frac{P+\widehat{P^2}}{\widehat{P^2}}.
\]
Finally,
\[
J\simeq\frac{\frac{H\langle X\rangle+\widehat{P^2}}{\widehat{P^2}}}{\frac{P+\widehat{P^2}}{\widehat{P^2}}}=\frac{H\langle X\rangle+\widehat{P^2}}{\widehat{P^2}}\Big/\left(B\cap\frac{H\langle X\rangle+\widehat{P^2}}{\widehat{P^2}}\right)\hookrightarrow\frac{F\langle X\rangle/\widehat{P^2}}{B},
\]
so $J$ is special.
\end{proof}

\begin{lemma}
Let $J$ be a finite-dimensional simple Jordan algebra over an algebraically closed field $F$. If $J$ satisfies $f$, then $J$ is either $\Delta$ or a Jordan algebra of a bilinear form.
\end{lemma}

\begin{proof}
Assume $\bar{E}=F$. Write $1=e_1+\cdots+e_n$ with $e_1,\dots,e_n$ pairwise orthogonal idempotents. We must show that if $n\geq 3$, then $f$ is not an identity for $F$ (or $J(V,f)$ is a simple algebra with $e_1$, or $e_1Qe_2$ idempotents).

If $n\geq 3$, then $H_n(D)\cong H_n(F)\supseteq H_3(F)$. Thus it suffices to show that $H_3(F)$ does not satisfy $f$. In $H_3(F)$ we may write
\[
[[a,b]^2,b]^2=n(a,b)
\]
(consider $H_3(E)\subseteq A^{(+)}$ with $A$ associative; the operation in $n(a,b)$ is that of $A$).

Denote $u_{ij}=1[i j]$ for $i\leq j\leq 3$. Also write $e_i=e_{ii}$ for $i=1,2,3$, so $u_{ij}=e_i+e_j$, while $2u_{ij}u_{jk}=u_{ik}$.

Let $a=e_2-e_3$ and $b=u_{12}+u_{13}$. Then in any $A^{(+)}$ we have
\begin{align*}
[a,b] &= -2\bigl((a^2,b,b)_+ - 2a\cdot(a,b,b)_+\bigr), \\
(e_2,b,b)_+ &= (e_2,u_{12}+u_{13},u_{12}+u_{13}) = \frac{1}{2}u_{12}(u_{12}+u_{13}) \\
&\quad -e_2\bigl(e_1+e_2+e_1+e_3+u_{23}\bigr) = -\frac{1}{2}\bigl(e_1+e_2+\frac{1}{2}u_{23}\bigr)-e_2-\frac{1}{2}u_{23} \\
&= \frac{1}{2}(e_1-e_2)-\frac{1}{4}u_{23}.
\end{align*}
Similarly, $(e_3,b,b)_+=\frac{1}{2}(e_1-e_3)-\frac{1}{4}u_{23}$.

Then
\begin{align*}
(a^2,b,b)_+ &= (e_2+e_3,b,b)_+ = e_1-\frac{1}{2}(e_2+e_3)-\frac{1}{2}u_{23} = \frac{1}{2}\bigl(-1+3e_1-u_{23}\bigr), \\
(a,b,b)_+ &= (e_2-e_3,b,b)_+ = -\frac{1}{2}(e_2-e_3), \\
2a\cdot(a,b,b)_+ &= -(e_2-e_3)^2 = -e_2-e_3.
\end{align*}
Hence $[a,b]^2 = -2\bigl(\frac{1}{2}(-1+3e_1-u_{23})+e_2+e_3\bigr) = -\bigl(1+e_1-u_{23}\bigr)$ and
\[
[[a,b]^2,b]^2 = [e_1-u_{23},b]^2 = -2\bigl(((e_1-u_{23})^2,b,b)_+ - 2(e_1-u_{23})\cdot(e_1-u_{23},b,b)_+\bigr).
\]
Since $(e_1-u_{23})^2 = e_1+e_2+e_3 = 1$, we have $((e_1-u_{23})^2,b,b)_+ = 0$. Further,
\begin{align*}
(e_1,b,b)_+ &= (e_1,u_{12}+u_{13},u_{12}+u_{13}) = \frac{1}{2}(u_{12}+u_{13})^2 - e_1(u_{12}+u_{13})^2 \\
&= \frac{1}{2}\bigl(e_1+e_2+e_1+e_3+u_{23}\bigr) - \frac{1}{2}(2e_1) = \frac{1}{2}\bigl(e_2+e_3+u_{23}\bigr), \\
(u_{23},b,b)_+ &= (u_{23},u_{12}+u_{13},u_{12}+u_{13}) = \frac{1}{2}(u_{13}+u_{12})^2 - u_{23}(u_{13}+u_{12})^2 \\
&= \frac{1}{2}\bigl(1+e_1+u_{23}\bigr) - u_{23}\bigl(1+e_1+u_{23}\bigr) \\
&= \frac{1}{2}\bigl(1+e_1+u_{23}\bigr) - u_{23} - e_2 - e_3 = \frac{1}{2}\bigl(-1+3e_1-u_{23}\bigr).
\end{align*}
Then
\begin{align*}
2(e_1-u_{23})\cdot(u_{23},b,b)_+ &= -e_1+u_{23}+3e_1+u_{23}^2 = 1+e_1+u_{23}, \\
2(e_1-u_{23})\cdot(e_1,b,b)_+ &= -\frac{1}{2}u_{23} - \frac{1}{2}u_{23} - e_2 - e_3 = e_2 - e_3 - u_{23}.
\end{align*}
Finally,
\begin{align*}
[[a,b]^2,b]^2 &= 2\bigl(-e_2-e_3-u_{23} - 1 - e_1 - u_{23}\bigr) = -4\bigl(1+u_{23}\bigr), \\
&= \bigl[[e_2-e_3, u_{12}+u_{13}], u_{12}+u_{13}\bigr] = -4\bigl(1+u_{23}\bigr), \\
&= \bigl[[e_1-e_3, u_{12}+u_{23}], u_{12}+u_{23}\bigr] = -4\bigl(1+u_{13}\bigr).
\end{align*}
Since
\[
n_3 = \bigl[[e_1-e_2, u_{13}+u_{23}], u_{13}+u_{23}\bigr] = -4\bigl(1+u_{12}\bigr),
\]
and
\begin{align*}
(u_{23},u_{13},u_{12}) &= \frac{1}{2}u_{12}^2 - \frac{1}{2}u_{23}^2 = \frac{1}{2}\bigl(e_1+e_2 - e_2 - e_3\bigr) \\
&= \frac{1}{2}\bigl(e_1-e_3\bigr) \neq 0 \quad \Rightarrow \quad (n_1,n_2,n_3)_+ \neq 0 \text{ in } H_3(F).
\end{align*}
Thus $H_3(F)$ does not satisfy $f$.

Now consider a simple Jordan PI-algebra. Let $X$ be a set of variables. We have proved that for all $a,b\in F\langle X\rangle$, $[[a,b]^2,a]^2\in F\langle X\rangle$ is a tetrad eater. If $x,y\in X$, since $[[x,y]^2,x]^2\in F\langle x,y\rangle$, there exists a polynomial $n(x,y)\in F\{x,y\}$ such that $[[x,y]^2,x]^2 = n(x,y)$. Define
\[
f(x_1,x_2,x_3,x_4,x_5,x_6) = \bigl(n(x_1,x_2), n(x_3,x_4), n(x_5,x_6)\bigr).
\]

\begin{lemma}
Let $J$ be a simple, central, finite-dimensional Jordan algebra over $F$ (with $F$ infinite). If $J$ satisfies the identity $f=0$, then $J\simeq F$ or $J\simeq J(V,f)$.
\end{lemma}
\end{proof}

\begin{proof}
We may assume $F$ algebraically closed.
Consider $\bar{J}=\bar{F}\otimes_F J$ with $\bar{F}$ algebraic closure of $F$; then $\bar{J}$ is simple, central finite dimensional over $\bar{F}$. If $\bar{J}=\bar{F}$ or $\bar{J}=J(V,f)$, the same is valid for $J$.

We know that $1=e_1+\cdots+e_n$, $e_1,\dots,e_n$ primitive pairwise orthogonal idempotents. It is enough to prove that $n\leq 2$.

Suppose $n\geq 3$, then $H_3(F)\leq J$ (since $J=H_n(D)$, $D$ composition algebra, and $H_n(D)\geq H_n(F)\geq H_3(F)$). We will show that $H_3(F)$ does not satisfy the identity $f=0$.

Denote $u_{ij}=1[i j]$ --- elementary matrix, $1\leq i<j\leq 3$, and denote $e_i=e_{ii}$, $i=1,2,3$. Then $u_{ij}^2=e_i+e_j$, $2u_{ij}u_{jk}=u_{ik}$.

Let $a=e_2-e_3$ and $b=u_{12}+u_{13}$. Then we have $[a,b]^2=-2\bigl((a^2,b,b)_+ - 2a\cdot(a,b,b)_+\bigr)$. Observe that
\begin{align*}
(e_2,b,b)_+ &= (e_2,u_{12}+u_{13},u_{12}+u_{13}) \\
&= \frac{1}{2}u_{12}(u_{12}+u_{13}) - e_2\bigl(e_1+e_2+e_1+e_3+u_{23}\bigr) \\
&= \frac{1}{2}\bigl(e_1+e_2+\frac{1}{2}u_{23}\bigr)-e_2-\frac{1}{2}u_{23} = \frac{1}{2}(e_1-e_2)-\frac{1}{4}u_{23}.
\end{align*}
Analogously, $(e_3,b,b)_+=\frac{1}{2}(e_1-e_3)-\frac{1}{4}u_{23}$.
Then
\begin{align*}
(a^2,b,b)_+ &= (e_2+e_3,b,b)_+ = e_1-\frac{1}{2}(e_2+e_3)-\frac{1}{2}u_{23} \\
&= \frac{1}{2}\bigl(-1+3e_1-u_{23}\bigr), \\
(a,b,b)_+ &= (e_2-e_3,b,b)_+ = -\frac{1}{2}(e_2-e_3),
\end{align*}
and
\[
2a\cdot(a,b,b)_+ = -(e_2-e_3)^2 = -e_2-e_3.
\]
Therefore
\begin{align*}
[a,b]^2 &= -2\Bigl(\frac{1}{2}\bigl(-1+3e_1-u_{23}\bigr)+e_2+e_3\Bigr) \\
&= -\bigl(1+e_1-u_{23}\bigr). \\
\Rightarrow\quad \bigl[[a,b]^2,b\bigr]^2 &= \bigl[e_1-u_{23},b\bigr]^2.
\end{align*}
Now we need to calculate
\begin{align*}
\bigl[e_1-u_{23},b\bigr]^2 &= -2\Bigl(\bigl((e_1-u_{23})^2,b,b\bigr)_+ - 2(e_1-u_{23})\bigl(e_1-u_{23},b,b\bigr)_+\Bigr).
\end{align*}
We have $(e_1-u_{23})^2=e_1+e_2+e_3=1$, therefore $\bigl((e_1-u_{23})^2,b,b\bigr)_+=0$.
\begin{align*}
(e_1,b,b)_+ &= (e_1,u_{12}+u_{13},u_{12}+u_{13}) \\
&= \frac{1}{2}(u_{12}+u_{13})^2 - e_1(u_{12}+u_{13})^2 \\
&= \frac{1}{2}\bigl(e_1+e_2+e_1+e_3+u_{23}\bigr) - \frac{1}{2}(2e_1) = \frac{1}{2}\bigl(e_2+e_3+u_{23}\bigr), \\
(u_{23},b,b)_+ &= (u_{23},u_{12}+u_{13},u_{12}+u_{13}) \\
&= \frac{1}{2}(u_{13}+u_{12})^2 - u_{23}(u_{13}+u_{12})^2 \\
&= \frac{1}{2}\bigl(1+e_1+u_{23}\bigr) - u_{23}\bigl(1+e_1+u_{23}\bigr) \\
&= \frac{1}{2}\bigl(1+e_1+u_{23}\bigr) - u_{23} - e_2 - e_3 \\
&= \frac{1}{2}\bigl(-1+3e_1-u_{23}\bigr), \\
2(e_1-u_{23})(u_{23},b,b)_+ &= -e_1+u_{23}+3e_1+u_{23}^2, \\
2(e_1-u_{23})(e_1,b,b)_+ &= -\frac{1}{2}u_{23} - \frac{1}{2}u_{23} - e_2 - e_3 = -e_2 - e_3 - u_{23}.
\end{align*}
Hence
\begin{align*}
\bigl[[a,b]^2,b\bigr]^2 &= 2\bigl(-e_2-e_3-u_{23} -1 - e_1 - u_{23}\bigr) = -4\bigl(1+u_{23}\bigr).
\end{align*}
Set
\begin{align*}
n_1 &= \bigl[[e_2-e_3,u_{12}+u_{13}],u_{12}+u_{13}\bigr] = -4\bigl(1+u_{23}\bigr), \\
n_2 &= \bigl[[e_1-e_3,u_{12}+u_{23}],u_{12}+u_{23}\bigr] = -4\bigl(1+u_{13}\bigr), \\
n_3 &= \bigl[[e_1-e_2,u_{12}+u_{23}],u_{13}+u_{23}\bigr] = -4\bigl(1+u_{12}\bigr).
\end{align*}
Since $(u_{23},u_{13},u_{12}) = \frac{1}{2}u_{12}^2 - \frac{1}{2}u_{23}^2 = \frac{1}{2}\bigl(e_1+e_2-e_2-e_3\bigr) = \frac{1}{2}(e_1-e_3) \neq 0$, it follows that $(n_1,n_2,n_3)_+ \neq 0$.
\end{proof}

\begin{theorem}[Posner]
Let $R$ be an associative prime PI-algebra over $F$ with center $Z$, and let $S=Z\setminus\{0\}$. Then $S^{-1}R$ is an $S^{-1}Z$-algebra which is central, simple and of finite dimension.
\end{theorem}

\begin{remark}
Denote $s^{-1}z$ by $s^{-1}z$.
\end{remark}

\begin{theorem}[Amitsur]
Let $R$ be a ring with involution $*$ and let the set of symmetric elements $H(R,*)$ satisfy a polynomial identity; then $R$ also satisfies a polynomial identity.
\end{theorem}

We want to show an analogue of Posner's theorem for Jordan algebras.

\begin{lemma}
Let $A$ be an associative algebra with involution such that $J = H(A, *)$ is prime and PI. Then $Z = Z(J) \neq 0$ and $Z^{-1}J = \{ z^{-1}a \mid 0 \neq z \in Z, a \in J \}$ is a central simple algebra over a field $Z^{-1}Z$ and the dimension of $Z^{-1}J$ over $Z^{-1}Z$ is finite.
\end{lemma}

\begin{proof}
By Amitsur's theorem, $A$ is a PI associative algebra.

If $A$ is prime, then by Posner's theorem, $Z(A) \neq 0$ and $Z(A)^{-1}A$ is central simple, finite-dimensional over $Z(A)^{-1}Z(A)$. We have $Z(A)^* = Z(A)$. Let $Z_0 = H(Z(A), *)$. If $Z_0 = 0$, then for any $z \in Z(A)$ we have $z^* + z = 0$. Then $z^* = -z$ and $z^2 = -zz^* = 0$. Since the center does not contain non-zero nilpotent elements, we must have $z = 0$. Thus $Z_0 \neq 0$. We have $Z(A)^{-1}A = Z_0^{-1}A$, since
\[
z^{-1}a = (zz^*)^{-1}(z^*a) \in Z_0^{-1}A.
\]
Since $Z_0^{-1}A$ is simple, we have that $H(Z_0^{-1}A, *)$ is simple and $Z^{-1}J = H(Z_0^{-1}A, *)$ has finite dimension over $Z_0^{-1}Z_0 \subseteq Z^{-1}Z$. Then $\dim_{Z^{-1}Z} Z^{-1}J \leq \dim_{Z_0^{-1}Z_0} Z^{-1}J$ is finite.

Suppose that $A$ is not prime. We have
\[
0 = \operatorname{rad} J \supseteq (J \cap \operatorname{rad} A) \Rightarrow J \cap \operatorname{rad} A = 0.
\]
Considering $A/\operatorname{rad} A$ instead of $A$, we may suppose that $A$ is semiprime (since $\operatorname{rad} A$ is invariant under the involution). We show that $A$ is $*$-prime, that is, if $I, K \triangleleft A$ with $I^* = I$, $K^* = K$ and $IK = 0$, then $I = 0$ or $K = 0$. Since $A$ is semiprime,
\[
(I \cap K)^2 = 0 \Rightarrow I \cap K = 0 \Rightarrow KI = 0.
\]
Then in $J$ we have $I \cap J$, $K \cap J \triangleleft J$ and $(I \cap J)(K \cap J) = 0$. Since $J$ is prime, we have $I \cap J = 0$ or $K \cap J = 0$. Suppose $I \cap J = 0$. Then $H(I, *) \subseteq I \cap J = 0$. Then for any $a \in I$ we have
\[
a^* + a, aa^* \in H(I, *) \Rightarrow a^* = -a, a^2 = 0.
\]
Then $A$ is nil of index $\leq 2$. Thus $I^3 = 0 \Rightarrow I = 0$.

Suppose that $0 \neq I, K \triangleleft A$, $IK = 0$ ($I$ and $K$ are not necessarily invariant under $*$). Since $A$ is semiprime, we have $KI = 0$. Consider $K \cap K^*$ and $I \cap I^* \triangleleft (A, *)$ (ideals invariant under $*$). We have
\[
(K \cap K^*)(I \cap I^*) = 0 \Rightarrow K \cap K^* = 0 \quad \text{or} \quad I \cap I^* = 0
\]
since $A$ is $*$-prime. Suppose $K \cap K^* = 0 \Rightarrow KK^* \subseteq K^*K = 0$.

We show next that $K$ is a prime algebra. Let $K_1, K_2$ be ideals of $K$ with $K_1K_2 = 0$. Consider $KK_1K$, $KK_2K \triangleleft A$. We have $(KK_1K)(KK_2K) = 0$. Let
\[
I_1 = \{ a + a^* : a \in KK_1K \}, \quad I_2 = \{ a + a^* : a \in KK_2K \}.
\]
We have that $I_1$ and $I_2$ are ideals of $J$. Indeed, for $a + a^* \in I_1$ we have
\begin{align*}
(a + a^*)h &= (KK_1K + K^*K_1^*K^*)h \\
&= \underbrace{KK_1h}_{\in KK_1K} + \underbrace{K^*K_1^*K^*h}_{=c^*} + \underbrace{hKK_1K}_{\in KK_1K} + \underbrace{hK^*K_1^*K^*}_{=b^*} \in I_1.
\end{align*}
Moreover $I_1I_2 = 0$. Since $J$ is prime, $I_1 = 0$ or $I_2 = 0$. Suppose $I_1 = 0 \Rightarrow \forall a \in KK_1K$ we have $a + a^* = 0$, $aa^* = 0$. As before, we show that $a^2 = 0$ for all $a \in KK_1K$ and consequently
\[
(KK_1K)^3 = 0 \Rightarrow \underbrace{(KK_1K)(KK_1K)}_{\subseteq K_1} \underbrace{(KK_1K)}_{\subseteq K_1} = 0.
\]
Since $KK_1K \triangleleft A$, we conclude $KK_1K = 0$, then
\[
\underbrace{(AK_1A)(AK_1A)}_{\subseteq K} \underbrace{(AK_1A)}_{\subseteq K} = 0 \Rightarrow AK_1A = 0 \Rightarrow K_1 = 0.
\]
Thus $K$ is a prime algebra, $K \cap K^* = 0$. By Posner's theorem $\widetilde{Z} = Z(K) \neq 0$ and $\widetilde{Z}^{-1}K$ is simple, central, finite-dimensional over $\widetilde{Z}^{-1}\widetilde{Z}$. Let $I = \{ a + a^* \mid a \in K \}$, $I \triangleleft J$ and $J_0 = \{ z + z^* \mid z \in \widetilde{Z} \} \neq 0$ (as in the previous case). Then $J_0 \subseteq Z(J)$ and $J_0^{-1}I = J_0^{-1}J$ is a simple central finite-dimensional over $J_0^{-1}J_0$, and the result follows as in the previous case.
\end{proof}

\begin{lemma}
Let $J$ be a prime Jordan algebra, generated by $3$ elements. Suppose that $J$ is i-special and satisfies an identity $f = 0$, where $f$ is $*$-alternating. Then $Z = Z(J) \neq 0$, $J$ is of finite dimension and Clifford type: $Z^{-1}J \simeq Z^{-1}J$ or $Z^{-1}J$ is a Clifford type algebra.
\end{lemma}

\begin{proof}
We have $J \simeq F\langle x, y, z \rangle / P$ for some $P \triangleleft F\langle x, y, z \rangle$. By Cohn's theorem we have $F\langle x, y, z \rangle = H\langle x, y, z \rangle$. Then by a lemma from the previous discussion we have that $J$ is special, and $J = H(S(J), +)$. Then by the previous lemma, the result follows.
\end{proof}

\begin{lemma}
Let $\mathcal{J}=J(V,f)$ with $\dim V \geqslant 2$ and $f$ non-degenerate.
\begin{enumerate}
\item[(i)] If $z \in \mathcal{J}$ is such that for any $x \in \mathcal{J}$, $(z,x,x)$ is nilpotent, then $z \in F1$.
\item[(ii)] If $z \in \mathcal{J}$ is such that for any $x \in \mathcal{J}$, there exists $n$ such that $(zx)^n=0$, then $z=0$.
\end{enumerate}
\end{lemma}

\begin{proof}
First we prove that $J(V,f)$ is quadratic, that is $a^2=\operatorname{tr}(a)a-n(a)1$, with $\operatorname{tr}(a),n(a) \in F$. In fact, we write $x=\alpha 1+x_1$, $\alpha \in F$, $x_1 \in V$. Then
\[
x^2=\alpha^2 1+f(x_1,x_1)1+2\alpha x_1=2\alpha x+\bigl(f(x_1,x_1)-\alpha^2\bigr)1,
\]
thus $\operatorname{tr}(x)=2\alpha$ and $n(x)=\alpha^2-f(x_1,x_1)$. In other words, $\operatorname{tr}(x)=2\alpha$ and $n(x)=\alpha^2-f(x_1,x_1)$. So $1$, $a$, $a^2$ are linearly dependent.

If in $\mathcal{J}$ exists a nilpotent element $a$, then $a^2=0$ (since if $a^n=0$ and $a^{n-1}\neq0$ implies $1,a,\dots,a^{n-1}$ are linearly independent). And if $a^2=0$, then $\operatorname{tr}(a)=0$ and $n(a)=0$.

Let $z=\alpha 1+v$ with $\alpha \in F$, $v \in V$ and take $x \in V$. We have $(z,x,x)=(\alpha 1+v,x,x)=(v,x,x)$. Then if $(z,x,x)$ is nilpotent, we compute
\begin{align*}
0&=(z,x,x)^2=(v,x,x)^2 \\
&=f(v,x)^2 f(x,x)1+f(x,x)^2 f(v,v)1-2f(v,x)^2 f(x,x)1 \\
&=f(x,x)\bigl(f(v,x)^2+f(x,x)f(v,v)-2f(v,x)^2\bigr)1 \\
&=f(x,x)\bigl(f(x,x)f(v,v)-f(v,x)^2\bigr)1.
\end{align*}
Thus if $f(x,x)\neq0$, then $f(v,x)^2=f(x,x)f(v,v)$.

We need to consider two cases:
\begin{enumerate}
\item[(a)] If $f(v,v)=0$, then $f(v,x)=0$ for all $x$, hence $v=0$.
\item[(b)] If $f(v,v)\neq0$, then $V=Fv\oplus(Fv)^\perp$. If $x \in (Fv)^\perp$ is such that $f(x,x)\neq0$, then
\[
0=f(v,x)^2=f(x,x)f(v,v)\neq0,
\]
which is a contradiction. Thus (b) is impossible.
\end{enumerate}
Therefore $v=0$ and consequently $z=\alpha 1 \in F1$. For part (ii), let $z=\alpha 1+v$. Taking $x=z$ we have $z^2=0$ and therefore $\operatorname{tr}(z)=0 \Rightarrow \alpha=0$. Then $z=v$, with $f(v,v)=0$. Taking $x \in V$, we have $xz=f(v,x)1 \in F1$, which implies $f(v,x)=0$, then $v=0$.
\end{proof}

\begin{lemma}
Let $\mathcal{J}$ be a Jordan algebra which satisfies $f=0$. Then
\[
\operatorname{Nil}(\mathcal{J})=\{a \in \mathcal{J}: a, ax \text{ are nilpotent for all } x \in \mathcal{J}\}.
\]
\end{lemma}

\begin{proof}
Let $\widetilde{N}=\{a \in \mathcal{J}: a, ax \text{ are nilpotent for all } x \in \mathcal{J}\}$. We already have that $\operatorname{Nil}(\mathcal{J}) \subseteq \widetilde{N}$. To show the other inclusion, let $a,b \in \widetilde{N}$ and $x,y \in \mathcal{J}$; it is enough to show that $a+b \in \operatorname{Nil}(\operatorname{alg}\langle a,b,x\rangle)$ and $(ax)y \in \operatorname{Nil}(\operatorname{alg}\langle a,ax,y\rangle)$. Then $\widetilde{N}$ will be an ideal.

Since all relations involve only three elements, we may assume that $\mathcal{J}=\operatorname{alg}\langle a,b,x\rangle$. We know that algebras generated by three elements are in $\mathcal{J}$ and satisfy $f=0$.

Consider $\overline{\mathcal{J}}=\mathcal{J}/\operatorname{Nil}(\mathcal{J}) \hookrightarrow \prod_\alpha \mathcal{J}_\alpha$, with $\mathcal{J}_\alpha$ simple and $\operatorname{Nil}(\mathcal{J}_\alpha)=0$ for all $\alpha$. Then $Z(\mathcal{J}_\alpha)\neq0$ and $Z(\mathcal{J}_\alpha)^{-1}\mathcal{J}_\alpha$ is a central simple algebra.

Suppose, for example, there exist $a,b \in \widetilde{N}$ such that $a+b \notin \operatorname{Nil}(\mathcal{J})$. Then there exists $\alpha$ and an epimorphism $\varphi_\alpha:\mathcal{J} \to \mathcal{J}_\alpha$ such that $\varphi_\alpha(a+b)\neq0$. Since $\varphi_\alpha(a)=\varphi_\alpha(b)=0$ in $\mathcal{J}_\alpha^{-1}\mathcal{J}_\alpha$ and we have the identity $f=0$, it follows that $\varphi_\alpha(a+b)=0$, which gives a contradiction. The other case is done analogously.
\end{proof}

\begin{definition}
Let $\mathcal{J}$ be a Jordan algebra. We define
\[
\operatorname{Alt}(\mathcal{J})=\{z \in \mathcal{J}:(z,x,x)=0 \quad \forall x \in \mathcal{J}\}
\]
the alternative center of $\mathcal{J}$.
\end{definition}

In $F\langle x,y\rangle$ we have the following formula
\[
[x,y]=-2(x^2,y,y)+4x\cdot(x,y,y);
\]
in particular $[x,y]$ is a Jordan element.

Let $\mathcal{J}$ be a Jordan algebra, $x,y \in \mathcal{J}$; we denote the element $-2(x^2,y,y)+4x(x,y,y) \in \mathcal{J}$ by $[x,y]$.

\begin{lemma}
Let $\mathcal{J}$ be a Jordan algebra in $\mathcal{S}$ which satisfies the polynomial identity $f=0$ given before, and suppose $\operatorname{Nil}(\mathcal{J})=0$. Then for all $a,b \in \mathcal{J}$ we have $[a,b]^2 \in Z(\mathcal{J})$.
\end{lemma}

\begin{proof}
Let $c=[a,b]^2$. To show that $c \in \mathcal{J}(y)$ we prove:
\begin{enumerate}
    \item $(c,p,p)$ is nilpotent for every $p \in \mathcal{J}$.
    \item For all $p,q \in \mathcal{J}$, $(c,p,p)q$ is nilpotent.
    \item $(c,p,p) \in \operatorname{Vie}(y)$ for all $p \in \mathcal{J}$ (thus $c \in \operatorname{Jalt}(y)$).
    \item If $z \in \operatorname{alt}(y)$ is nilpotent then $z=0$.
    \item For all $x,y \in \mathcal{J}$, $(x,c,y)$ is nilpotent.
    \item $c \in Z(y)$.
\end{enumerate}
For (1) consider $\mathcal{J}_0 = \mathcal{J}_y\langle a,b,p\rangle$. Then $\mathcal{J}_0 \in \operatorname{Jord}$ and $f=0$ in $\mathcal{J}_0$. It is enough to show that $(c,p,p) \in \operatorname{Nil}(\mathcal{J}_0)$. We have
\[
\mathcal{J}_0/\operatorname{Nil}(\mathcal{J}_0) \longleftrightarrow \prod_\alpha \mathcal{J}_\alpha,
\]
with $\mathcal{J}_\alpha^{-1}\mathcal{J}_\alpha \in \{ \}_\alpha^{-1}\}_\alpha$, $\forall \alpha (V,f)$. Let $\varphi_\alpha: \mathcal{J} \to \mathcal{J}_\alpha$ be the epimorphism associated to the subdirect product. If $(c,p,p) \notin \operatorname{Nil}(\mathcal{J})$, then there exists $\alpha$ such that $\varphi_\alpha(c,p,p) \neq 0$ in $\mathcal{J}_\alpha$. In $\mathcal{J}_\alpha$ we have $\varphi_\alpha(c)=[\varphi_\alpha(a),\varphi_\alpha(b)]^2$. If we show that for every $x,y \in \mathcal{J}_\alpha$ we have $[x,y]^2 \in \mathcal{J}(\mathcal{J}_\alpha)$, then $\varphi_\alpha((c,p,p))=0$, contradicting the assumption that $(c,p,p) \notin \operatorname{Nil}(\mathcal{J}_0)$. To show that $([x,y],\mathcal{J}_\alpha,\mathcal{J}_\alpha)=0$, we prove $([x,y],\mathcal{J}_\alpha,\mathcal{J}_\alpha)=0$. If $\mathcal{J}_\alpha^{1/7}$ is a field it is valid. Suppose that $\mathcal{J}_\alpha^{-1}\mathcal{J}_\alpha = \mathcal{J}(v,f)$. Since $[x,y]=-2((x^2,y,y)-2x(x,y,y))$, we have
\[
(x^2,y,y)-2x(x,y,y)=0.
\]
\end{proof}

\begin{lemma}
Let $\mathcal{J}$ be a unital Jordan algebra. Every Jordan algebra satisfies the identity
\[
n_1(x,y)\cdot x^2 - n_2(x,y)\cdot x + n_3(x,y)=0,
\]
where
\begin{align*}
& n_1(x,y)=[x,y]^2, \quad n_2(x,y)=[x^2,y]\circ[x,y], \\
& n_3(x,y)=\frac{1}{2}\left([x^2,y^2]^2 - 2[y\circ x^2,x]\circ[y,x]\right).
\end{align*}
\end{lemma}
\begin{proof}
Exercise (but it is enough to prove for special Jordan algebras).
\end{proof}

\begin{lemma}
Let $\mathcal{J}$ be a unital Jordan algebra over an infinite field $F$, such that for every $x,y \in \mathcal{J}$, $[x,y] \in F$. If there exist $a,b \in \mathcal{J}$ with $[a,b]^2 \neq 0$, then $\mathcal{J}$ is quadratic over $F$.
\end{lemma}
\begin{proof}
By the previous lemma, if $x \in \mathcal{J}$ is such that there exists $y \in \mathcal{J}$ with $[x,y]^2 \neq 0$, then $x$ is quadratic in $y$. Observe that for all $x,y \in \mathcal{J}$, $n_i(x,y) \in F$, since $n_2$ and $n_3$ are linearizations of $n_1$. In particular, elements $a$ and $b$ are quadratic. For $x \in \mathcal{J}$, consider for $\beta \in F$:
\[
[a,x+\beta b]^2 = [a,x]^2 + \beta[a,x]\cdot[a,b] + \beta^2[a,b]^2.
\]
Since $F$ is infinite, there exist $\beta_1 \neq \beta_2 \in F$ such that $[a,x+\beta_1 b]^2 \neq 0$ and $[a,x+\beta_2 b]^2 \neq 0$. Then $x+\beta_i b$ is quadratic for $i=1,2$. Thus there exist $\alpha_i \in F$ such that
\[
(x+\beta_i b)^2 - \alpha_i(x+\beta_i b) \in F.
\]
Since $b$ is quadratic we have
\[
x^2 + 2\beta_i(x \cdot b) - \alpha_i x \in F + F b, \quad i=1,2.
\]
Multiplying the first equation by $\beta_2$ and the second by $\beta_1$ and subtracting one from the other we obtain
\[
x^2 - \alpha x = \gamma + \delta b, \quad \alpha,\gamma,\delta \in F.
\]
We do the same for $a$ to get
\[
x^2 - \alpha' x = \gamma' + \delta' b.
\]
If $\delta=0$ or $\delta'=0$, then $x$ is quadratic. If $\delta\delta' \neq 0$, then
\[
\delta^2(\delta')^2[a,b]^2 = [x^2-\alpha x-\gamma, x^2-\alpha' x-\gamma']^2 = 0,
\]
which implies $[a,b]^2=0$. Then $\delta=0$ or $\delta'=0$, and thus $x$ is quadratic.
\end{proof}

\begin{corollary}
If $\mathcal{J}$ satisfies the conditions of the previous lemma, then $\mathcal{J} \simeq \mathcal{J}(v,f)$.
\end{corollary}
\begin{proof}
It is enough to observe that any quadratic Jordan algebra is a Jordan algebra of a bilinear form. The space $V$ consists of trace zero elements.
\end{proof}

\begin{theorem}
Let $\mathcal{J} \in \mathcal{S}$ be a prime Jordan algebra which satisfies $f=0$. Suppose that $[x,y]^2 \neq 0$ in $\mathcal{J}$. Then $\mathcal{J}^{-1}\mathcal{J} \simeq \mathcal{J}(v,f)$.
\end{theorem}

\begin{proof}
By Lemma 1, for all $x, y \in \mathcal{J}$ we have $[x, y]^2 \in \delta\{\}$. Since there exist $a, b \in J$ with $0 \neq [a, b]^2 \in \xi(y)$, then $j = z(y)$ is a domain and we may consider $3^{-1} \mathcal{J}_2$, which is a Jordan algebra over $\left.3_1^{-1}\right\}$, and for all $x, y \in 3^{-1} \mathcal{J}_1$, $[x, y]^2 \in 3_1^1 \mathcal{J}_2$. Then we are in the conditions of Lemma 2. Then by Corollary,
\[
j^{-1} j = j(v, f).
\]
Since $[x, y] = -2(x^2, y, y) - 2x(x, y, y)$, we have
\begin{align*}
& (x^2, y, y) - 2x(x, y, y) = (\operatorname{tr}(x)x - n(x), y, y) \\
& - 2x(x, y, y) = \operatorname{tr}(x)(x, y, y) - 2x(x, y, y) \\
& = (\operatorname{tr}(x) - 2x)(x, y, y).
\end{align*}
Next, observe that $\operatorname{tr}(x) - 2x \in V$ (if $x = \alpha + x_1$, with $\alpha \in F$ and $x_1 \in V$, then $\operatorname{tr}(x) = 2\alpha$ and $\operatorname{tr}(x) - 2x = -x_1 \in V$) and $(x, y, y) \in V$ (if $y = \beta + y_1$, with $\beta \in F$ and $y_1 \in V$, then $(x, y, y) = (x, y_1, y_1) \in V$). Then
\[
[x, y]^2 = -2(\operatorname{tr}(x) - 2x)(x, y, y) \in V \cdot V \leq F \cdot \mathcal{J}(v+1).
\]
Thus we conclude that $(c, p, p)$ is nilpotent for all $p \in \mathcal{J}$.

To show (ii), let $\mathcal{J}_0 = \operatorname{alg}_y\langle c, p, q\rangle$. We prove that $(c, p, p) \in \operatorname{Nil}(\mathcal{J}_0)$. Then we consider the embedding
\[
\mathcal{J}_0 / \operatorname{Nil}(\mathcal{J}_0) \hookrightarrow \prod_\alpha \mathcal{J}_\alpha.
\]
We have, for any $\alpha$, $3^{-1} \mathcal{J}_\alpha$ is of type $\mathcal{J}(v, f)$ or is a field. For any $\bar{p} \in \bar{\mathcal{J}}_\alpha$, $(\bar{c}, \bar{p}, \bar{p})$ is nilpotent. Thus $\bar{c} \in \mathcal{Z}(\mathcal{J}_\alpha)$ (since $\mathcal{J}_\alpha = \mathcal{J}(v, f)$). Then $(\bar{c}, \bar{p}, \bar{p}) = 0$, equivalently $(c, p, p) \in \operatorname{Nil}(\mathcal{J}_0)$. Consequently, $(c, p, p)q$ is nilpotent.

The proof of (iii) follows from Lemma 2 and (i) and (ii) of this lemma. We conclude that $(c, p, p) \in \operatorname{Nil}(\mathcal{J}) = \{0\}$. Thus $c \in \mathcal{Z}(\mathcal{J})$.

For (iv), let $\mathcal{J}_0 = \operatorname{alg}_y\langle z, x\rangle$, $x \in \mathcal{J}$. We will show that $z \in \operatorname{Nil}(\mathcal{J}_0)$ and then by Lemma 2 we have $z \in \operatorname{Nil}(\mathcal{J})$.
\end{proof}

\begin{remark}[Quadratic ideals]
We consider Jordan algebras satisfying the identity $[x, y]^2 = 0$.
\end{remark}

\begin{theorem}
Let $\mathcal{J}$ be a nondegenerate Jordan algebra which satisfies the identity $[x,y]^2=0$. Then $\mathcal{J}$ is associative and therefore is a domain.
\end{theorem}

\begin{proof}
We show that $(a,b,c)$ is an absolute zero divisor for any $a,b,c \in \mathcal{J}$. Since $\mathcal{J}$ is nondegenerate this will imply that $(a,b,c)=0$ for all $a,b,c \in \mathcal{J}$. Since $[x,y]^2=0$ we have $(x^2,y,y)=2x(x,y,y)$. Since the associator in the second entry is a derivation we have $2(x,y,y)y=(x,y^2,y)(x)$ and since $[R_y,R_{y^2}]=0 \Rightarrow (x,y^2,y)=(x,y,y^2)$.

Linearizing $(*)$ in $y$ we obtain
\[
(x^2,y,z)+(x^2,z,y)=2x(x,y,z)+2x(x,y,y).
\]
In particular
\[
(***)\quad (x^2,y,y^2)+(x^2,y^2,y)=2x\left((x,y,y^2)+(x,y,y^2)\right).
\]
Again, since $[R_y,R_{y^2}]=0$,
\[
(****)\quad 2(x^2,y,y^2)=4x(x,y,y^2).
\]
Observe that $4(x,(x,y,y),y)=4(x(x,y,y))y$
\[
\begin{aligned}
&-4x((x,y,y)y)=2(x^2,y,y)y-2x(x,y,y^2)= \\
&\stackrel{(**)}{=}(x^2,y,y^2)-(x^2,y,y^2)=0.
\end{aligned}
\]
Thus
\[
(x,(x,y,y),y)=0 \quad (*v) \text{ for all } x,y \in \mathcal{J}.
\]
Linearizing $(*)$ in $x$ we get
\[
(xz,y,y)=x(z,y,y)+z(x,y,y). \quad (*6)
\]
Then
\[
(xy,y,y)=x(y,y,y)+y(x,y,y)=\frac{1}{2}(x,y,y^2). \quad (*7)
\]
Linearizing $(****)$ in $x$,
\[
(xy,y^2,y)=\frac{1}{2}(x,y^2,y^2). \quad (*8)
\]
Then $(x,y,y^2)y \stackrel{(*7)}{=} 2(xy,y,y)y \stackrel{(**)}{=} (xy,y,y^2)$.
\[
0=[R_y,R_{y^2}](xy,y^2,y) \stackrel{(*8)}{=} \frac{1}{2}(x,y^2,y^2).
\]
Thus $(x,y^2,y^2)=0$. \quad (*9)

Now,
\[
\begin{gathered}
(x,y,y)(x,y,y)=(x,y(x,y,y),y)-y(x,(x,y,y),y) \\
\stackrel{(*v)}{=} \frac{1}{2}\left(x,\left(x,y,y^2\right),y\right). \quad (*10)
\end{gathered}
\]
From one side we have
\[
\left(x,\left(x,y,y^2\right)\right)y \underset{(****)}{=} \frac{1}{2}(x^2,y,y^2)y \underset{(*9)}{=} \frac{1}{4}(x^2,y^2,y^2),
\]
from the other side $x\left(\left(x,y,y^2\right)y\right) \stackrel{(*9)}{=} \frac{1}{2}x\left(x,y^2,y^2\right) \stackrel{(*)}{=} \frac{1}{4}\left(x^2,y^2,y^2\right)$.
Therefore from $(*10)$ follows $(x,y,y)^2=0$. \quad (*11)

Next we show that $(x,y,y)$ is an absolute zero divisor or equivalently that $U_{(x,y,y)}=0$. For any $z \in \mathcal{J}$, by definition
\[
\begin{aligned}
& zU_{(x,y,y)}=zR_{(x,y,y)}^2-zR_{(x,y,y)^2}= \\
=& zR_{(x,y,y)}^2=(z(x,y,y))(x,y,y).
\end{aligned}
\]
The linearization of $(*11)$ in $x$ results in
\[
(x,y,y)(z,y,y)=0. \quad (*12)
\]
Then
\[
\begin{aligned}
& U_{(x,y,y)}z=(z(x,y,y))(x,y,y) \stackrel{(*6)}{=} (x,z,y,y)(x,y,y) \\
& \quad - (x(z,y,y))(x,y,y) \stackrel{(*12)}{=} - (x(z,y,y))(x,y,y). \quad (*13)
\end{aligned}
\]
Now, observe
\[
\begin{aligned}
& ((z,y,y)x)(x,y,y)=(x,y,y)((z,y,y)x) \\
& = -((x,y,y),(z,y,y),x)+\underbrace{((x,y,y)(z,y,y))x}_{=0 \text{ by } (*12)}.
\end{aligned}
\]
Also
\[
((z,y,y)x)(x,y,y)=((z,y,y),x,(x,y,y))-(z,y,y)(x(x,y,y)).
\]
Since $(x(x,y,y))=\frac{1}{2}(x^2,y,y)$, we have
\[
\begin{aligned}
((z,y,y)x)(x,y,y)&=((z,y,y),x,(x,y,y))-\frac{1}{2}\underbrace{(z,y,y)(x^2,y,y)}_{=0 \text{ by } (*12)} \\
&=((z,y,y),x,(x,y,y)).
\end{aligned}
\]
Therefore $((x,y,y),(z,y,y),x)=((z,y,y),x,(x,y,y))$.

Now, by the Jacobi identity for the associator,
\[
0=((x,y,y),x,(z,y,y))+(x,(z,y,y),(x,y,y))+((z,y,y),(x,y,y),x).
\]
Using the previous equality, this becomes
\[
0=((x,y,y),x,(z,y,y))+((x,y,y),(z,y,y),x)+((z,y,y),(x,y,y),x).
\]
But $((x,y,y),(z,y,y),x)=((z,y,y),x,(x,y,y))$, so
\[
0=((x,y,y),x,(z,y,y))+((z,y,y),x,(x,y,y))+((z,y,y),(x,y,y),x).
\]
Since $((z,y,y),(x,y,y),x)=((z,y,y)(x,y,y))x=0$ by $(*12)$, we obtain
\[
0=((x,y,y),x,(z,y,y))+((z,y,y),x,(x,y,y)).
\]
Thus $((x,y,y),x,(z,y,y))=-((z,y,y),x,(x,y,y))$.

Now, from $(*13)$,
\[
U_{(x,y,y)}z = - (x(z,y,y))(x,y,y) = -((x,y,y),x,(z,y,y)).
\]
But we have shown $((x,y,y),x,(z,y,y)) = -((z,y,y),x,(x,y,y))$, and by symmetry $((z,y,y),x,(x,y,y)) = -((x,y,y),x,(z,y,y))$, so $((x,y,y),x,(z,y,y))=0$. Hence $U_{(x,y,y)}z=0$ for all $z$, so $(x,y,y)$ is an absolute zero divisor. Since $\mathcal{J}$ is nondegenerate and contains no nonzero absolute zero divisors, we have $(x,y,y)=0$ in $\mathcal{J}$. Then $(y,y,x)=0$ as well.

Therefore $\mathcal{J}$ is alternative. In particular the associator is a derivation in any entry. Then
\[
\begin{aligned}
4(x,y,(x,y,z)) &= -4(x,(x,y,z),y) \\
&= -4(x(x,y,z))y+4x((x,y,z)y)-2(x,y,z)y+2x(x,y,z) \\
&= -\left(x^2,y^2,z\right)+\left(x^2,y^2,z\right)=0.
\end{aligned}
\]
Thus $(x,y,(x,y,z))=0$. Now,
\[
(x,y,z)^2=(x,y,z)(x,y,z)=(x,y,z(x,y,z))-\underbrace{(x,y,(x,y,z))}_{=0}z=\left(x,y,\left(x,y,z^2\right)\right)=0.
\]
Finally, for $t \in \mathcal{J}$,
\[
tR_{(x,y,z)}^2 - tR_{(x,y,z)^2} = tR_{(x,y,z)}^2 = (t,(x,y,z))(x,y,z)=t(x,y,z)^2=0.
\]
Thus $(x,y,z)$ is an absolute zero divisor and hence $(x,y,z)=0$ in $\mathcal{J}$.

Therefore $\mathcal{J}$ is associative.
\end{proof}

\begin{definition}
Let $\mathcal{J}$ be a Jordan algebra. A subspace $K \subseteq \mathcal{J}$ is called a quadratic ideal if for any $a \in K$, $\mathcal{J} U_a \subseteq K$, or equivalently, for every $x \in \mathcal{J}$ the element $\{a x a\}$ is in $K$.

Recall: $\{a b c\}=(a b) c+(c b) a-(a c) b$.
\end{definition}

Denote $(K : \mathcal{J}) = \{ a \in K : a \mathcal{J} \subseteq K \}$.

\begin{remark}
If $J$ is a Jordan algebra with unit element, then every quadratic ideal of $J$ is a subalgebra of $J$. Indeed, let $K$ be a quadratic ideal of $J$. Since $U_{a,b}=R_aR_b+R_bR_a-R_{ab}$, we already know that $K$ is a subspace of $J$. For any $a,b\in K$ we have $2ab=1U_{a,b}1\in K$. Therefore, $K$ is a subalgebra of $J$.
\end{remark}

\begin{lemma}
Let $J$ be a Jordan algebra with unit $1$, and let $K$ be a quadratic ideal of $J$. Then
\begin{enumerate}
\item[(i)] $(K:J)K\subseteq K$,
\item[(ii)] $K/(K:J)$ is special,
\item[(iii)] $(K:J)$ is a quadratic ideal of $J$,
\item[(iv)] $(K:J)^2\subseteq((K:J):J)$.
\end{enumerate}
\end{lemma}

\begin{proof}
To show (i), let $a\in(K:J)$, $k\in K$ and $x\in J$. We have $xU_{a+k}\in K$ since $a+k\in K$. Next,
\[
xU_{a,k}=\frac{1}{2}\left(xU_{a+k}-xU_a-xU_k\right)\in K.
\]
Since
\[
xU_{a,k}=(xa)k+(xk)a-x(ak)
\]
and observing that $(xa)k,(xk)a\in K$ (since $a\in(K:J)$), it follows that $x(ak)\in K$. Consequently, $ak\in(K:J)$, which means that $(K:J)$ is an ideal in $K$.

To show (ii), consider the quotient space $J/K$. Observe that for all $k\in K$ we have $KR_k\subseteq K$ since $K$ is a subalgebra of $J$. Then we take $\varphi:K\to\operatorname{End}(J/K)$ defined as $\varphi(a)=\overline{R_a}$, $a\in K$, where $\overline{R_a}(x+K)=ax+K$ for all $x\in J$. Then
\[
\operatorname{Ker}\varphi=\{k\in K:Jk\subseteq K\}=(K:J).
\]
By the homomorphism theorem,
\[
\overline{\varphi}:K/(K:J)\to\operatorname{End}(J/K)
\]
is linear. We will show that $\overline{\varphi}$ is a monomorphism of Jordan algebras into $\operatorname{End}(J/K)^{(+)}$. For $a,b\in K$ and $x\in J$ we have
\begin{align*}
\overline{\varphi}(\bar{a}\bar{b})&=\overline{R_{ab}}=\overline{R_aR_b+R_bR_a-U_{a,b}}\\
&=\overline{R_a}\,\overline{R_b}+\overline{R_b}\,\overline{R_a}=\overline{R_a}\cdot\overline{R_b}+\overline{R_b}\cdot\overline{R_a}\\
&=\overline{\varphi}(\bar{a})\cdot\overline{\varphi}(\bar{b}).
\end{align*}

To prove (iii), let $a\in(K:J)$ and $x\in J$. Since $xa\in K$, we have $xU_a=2(xa)a-xa^2=ka-xa^2$ for some $k\in K$. Thus $ka\in(K:J)$. We need to show that $(xU_a)y\in K$ for all $y\in J$. Then we need to show that $(xa^2)y\in K$. We have
\begin{align*}
(xa^2)y&=a^2(xy)+(a^2,x,y)\\
&=k'+2(a,x,ya)\\
&=k'+2(a,x,k'')\in K+K\subseteq K.
\end{align*}

Finally, to prove (iv),
\[
((K:J):J)=\{k\in(K:J):Jk\subseteq(K:J)\}=\{k\in(K:J):(kx)y\in K\text{ for all }x,y\in J\}.
\]
Let $a\in(K:J)$ and $x,y\in J$. Then
\begin{align*}
(a^2x)y&=a^2(xy)+(a^2,x,y)\\
&=k+2(a,x,ay)\\
&=k+2((ax)(ay)-a(x(ay)))\in K.
\end{align*}
\end{proof}

\begin{lemma}
The following identity is satisfied by every Jordan algebra:
\[
\left(z^2U_{x_1}\cdots U_{x_r}\right)^n=\left(x_r^2U_{x_{r-1}}\cdots U_{x_1}U_z\right)^{n-1}U_zU_{x_1}\cdots U_{x_r}.
\]
\end{lemma}

\begin{proof}
Prove it for $x_1,\dots,x_r$ generators of a free Jordan algebra.
\end{proof}

\begin{theorem}
Let $J$ be a Jordan algebra with unit element and let $K$ be a quadratic ideal of $J$ and $I=\operatorname{id}_J\langle (K:J)^2\rangle$. Then
\begin{enumerate}
\item[(i)] For every $a\in I$, there exists $b\in J$ such that $bU_{1-a}\in J+(K:J)$.
\item[(ii)] If $F$ has more than $\dim J$ elements, then for all $a\in I$, there exists $n=n(a)$ such that $a^n\in(K:J)$.
\end{enumerate}
\end{theorem}

Primitive algebras: We say that a Jordan algebra with unity is primitive if there exists $K\neq J$ maximal quadratic ideal of $J$, such that for all non-zero ideals $I$ of $J$ we have $K+I=J$. Equivalently, $J$ is primitive if there exists $K\neq J$ maximal quadratic ideal of $J$, which does not contain non-zero ideals of $J$.

\begin{lemma}
Let $J$ be a Jordan algebra with unit element. Let $K$ be a quadratic ideal of $J$ and $I\triangleleft J$. Then $K+I$ is a quadratic ideal of $J$.
\end{lemma}

\begin{lemma}
Let $J$ be a Jordan algebra and suppose that $J$ is primitive and contains unity. Then $J$ is prime and $\operatorname{Nil}(J)=0$.
\end{lemma}

\begin{proof}
Suppose there exists $I\neq0$, $I\triangleleft J$, with $I+K=J$, then there exist $k\in K$ and $a\in I$ with $1=k+a$. Then $1-k=a$ is nilpotent $\Rightarrow\exists n>0$ such that $(1-k)^n=0\Rightarrow1\in\operatorname{id}_J\langle k\rangle\subseteq K$, this we get a contradiction. Hence $\operatorname{Nil}(J)=0$.

Suppose now that there exist non-zero ideals $I$ and $L$ of $J$ such that $IL=0$. We have $K+I=J=K+L$. Thus there exist $k,k'\in K$, $a\in I$ and $b\in L$ such that $k+a=1=k'+b$. Since $(1-k)(1-k')=ab\in IL=0$, we have $1=k+k'-kk'\in K$ which is a contradiction. Consequently, $J$ is prime.
\end{proof}

\begin{lemma}
Let $J$ be a Jordan algebra, let $J$ be primitive and $K$ a maximal quadratic ideal. Then $K_1^2=0$, where $K_1=(K:J)$.
\end{lemma}

\begin{proof}
Let $I=\operatorname{id}_J\langle K_1^2\rangle$. Suppose $K_1^2\neq0$. Then $I\neq0$ and $K+I=J$. Then there exist $k\in K$ and $a\in I$, such that $1=k+a$, that is $1-a=k$. By Theorem 1(i) there exists $b\in J$ such that $bU_{1-a}\in J+K_1$. But $bU_{1-a}=bU_k\in K$. Thus $1\in K$, which is a contradiction. Then $I=0$. In this condition $K_1^2=0$.
\end{proof}

Let $\Omega$ be the $T$-ideal of all $S$-identities, that is $\Omega=\ker\pi$, where $\pi:J[x]\to J\langle x\rangle$ is a surjective homomorphism of Jordan algebras such that $\pi|_x$ is the identity on $x$. Then $J\in\overline{S\text{-ford}}\Longleftrightarrow S(J)=0$. ($S(J)$ is defined as the set $\{f(a_1,\dots,a_n):f\in S,\ a_i\in J\}$.)

\begin{definition}
Let $J$ be a Jordan algebra. We define $S_0=\bigcap\{0\neq I\triangleleft J\}$ the heart of Jordan algebra $J$.
\end{definition}

\begin{lemma}
Let $J$ be a primitive Jordan algebra with unit element such that $F\notin S\text{-ford}$. Then $S_0=S(J)^3\neq0$ is the unique minimal bilateral ideal of $J$.
\end{lemma}

\begin{proof}
Since $\operatorname{Nil}(J)=0$ and $S(J)\neq0$, we have $S(J)^3\neq0$. It is enough to show that $S(J)^3$ is contained in every non-zero ideal of $J$.

Let $I\triangleleft J$, $I\neq0$. Then $J=I+K$, where $K$ is a maximal quadratic ideal of $J$. Since $K/K_1$ is special, we have $S(K/K_1)=0$. Then $S(K)\subseteq K_1$. Since $I$ is an ideal, we have
\[
S(J)\subseteq I+S(K)\subseteq I+K_1\Rightarrow S(J)^2\subseteq I+K_1^2=I,
\]
since $K_1^2=0$. Thus $S(J)^3\subseteq I$.
\end{proof}

Simple $i$-exceptional algebra: Let $J$ be a primitive Jordan algebra with unity. We have seen that $0\neq S_0=S(J)^3$ is minimal ideal of $J$.

Fix $f=f(x_1,\dots,x_m)\in S^3$ such that $f|_J\neq0$. Let $n=\deg f$. We define for $s\in S_0$,
\[
\operatorname{Spec}(s)=\{\alpha\in F:1-\alpha^{-1}s\text{ not invertible}\},
\]
and for $K\triangleleft J$, we denote $\check{K}$ the maximal bilateral ideal of $J$ contained in $K$.

\begin{lemma}
Let $s\in S_0$ and $\alpha\in\operatorname{Spec}(s)$, $\alpha\neq0$. If $P=JU_{s-\alpha1}$, then $f|_P=0$.
\end{lemma}

\begin{proof}
$P\neq J$ since otherwise we would have $aU_{s-\alpha1}=1$ for some $a\in J$, and thus $s-\alpha1$ is invertible, which contradicts the hypothesis.

By a former lemma, there exists a maximal ideal which contains $P$. If $K^*\neq0$ then $S_0\subseteq K^*$ and then $s(s-\alpha^{-2}1)=1U_{s-\alpha1}\in P\subseteq K$, which implies that $1\in K$. Thus we have $K=0$.

Since $f\in S^3$ we will write
\[
f=\sum_i(f_i,f_i)f_i\quad\text{where }f_{ij}\in S.
\]
Moreover $K/K_1$ is special. Then
\[
f_{ij}(K/K_1)=0\Rightarrow f_{ij}(K)\subseteq K_1,\ \forall i,j.
\]
Then
\[
f(K)\subseteq\sum(f_i(K)f_i(K))f_i(K)\subseteq K_1^3\subseteq K_1^2=0.
\]
Since $P\subseteq K$, we have $f|_P=0$.
\end{proof}

\begin{lemma}
For any $s \in S_0$ the set $\operatorname{spec}(s) \setminus \{0\}$ has at most $2n$ elements, where $n = \deg f$.
\end{lemma}

\begin{proof}
Suppose there exist $\alpha_1,\dots,\alpha_{2n+1} \in \operatorname{spec}(s)$ non-zero and $\alpha_i \neq \alpha_j$ if $i \neq j$.
Let $P_i = J U_{s-\alpha_i 1}$. Then $P_i$ is an ideal, $P_i$ is quadratic and by a previous lemma $f/P_i = 0$ for all $i=1,\dots,2n+1$.
Consider $g(x) = \prod_{i=1}^{2n+1} (x-\alpha_i)$, $g_i(x) = \frac{g(x)}{x-\alpha_i}$.
The polynomials $g_i(x)$ are pairwise prime. Then there exist $h_1,\dots,h_{2n+1} \in F[x]$ such that $\sum_{i=1}^{2n+1} h_i(x) g_i(x) = 1$.
In particular, $\sum_{i=1}^{2n+1} h_i(s) g_i(s) = 1$.
Then $1 = U_1 = U_{\sum_{i=1}^{2n+1} h_i(s) g_i(s)} = \sum_{i,j} U_{h_i(s) g_i(s), h_j(s) g_j(s)}$.
Let $x^{(1)},\dots,x^{(m)} \in \mathcal{J}$ be arbitrary elements. Then
\[
x^{(k)} = \sum_{i,j} x_{ij}^{(k)} \quad \text{with } x_{ij}^{(k)} \in J U_{h_i(s) g_i(s), h_j(s) g_j(s)}.
\]
Since the field $F$ is infinite, we may assume $f$ is homogeneous. Let $\sum_k f_k$ be the linearization of $f$, then
\begin{align*}
f\left(x^{(1)},\dots,x^{(m)}\right) &= f\left( \sum_{i,j}^{2n+1} x_{ij}^{(1)}, \dots, \sum_{i,j}^{2n+1} x_{ij}^{(m)} \right) \\
&= \sum_k f_k\left( x_{i_1 j_1}^{(1)}, \dots, x_{i_m j_m}^{(m)} \right).
\end{align*}
Since $m \leq n$, there exists $r \leq 2n+1$ such that $r \notin \{i_1,j_1,\dots,i_m,j_m\}$. We show that $x_{i_1 j_1}^{(1)},\dots,x_{i_m j_m}^{(m)} \in P_r$.
We have $P_r = J U_{s-\alpha_r 1}$ and $(x-\alpha_r 1)$ divides $g_i(x)$ for all $i \neq r$. Now,
\[
U_{ab,ac} = U_{b,c} U_a = U_a U_{b,c} \quad \forall a,b,c \in F[x].
\]
In fact, since the algebra generated by $x$ is special, we have
\begin{align*}
y U_{ab,ac} &= \frac{1}{2}\big((ab)y(ac) + (ac)y(ab)\big) \\
&= \frac{1}{2}\big(a(byc + cyb)a\big) = \frac{1}{2}(byc + cyb) U_a \\
&= y U_{b,c} U_a.
\end{align*}
Thus $U_{ab,ac} = U_{b,c} U_a$. Moreover, since the operators commute we have $U_{b,c} U_a = U_a U_{b,c}$.
Writing $h_i(x) g_i(x) = (s-\alpha_r 1) f_i(x)$, we have
\begin{align*}
x_{ij}^{(k)} &\in J U_{h_i(s) g_i(s), h_j(s) g_j(s)} = J U_{(s-\alpha_r 1) f_i(s), (s-\alpha_r 1) f_j(s)} \\
&= U_{f_i(s), f_j(s)} U_{s-\alpha_r 1} \subseteq P_r.
\end{align*}
Since $f/P_r = 0$, it implies that $\sum_k f_k\left( x_{i_1 j_1}^{(1)}, \dots, x_{i_m j_m}^{(m)} \right)$ is zero. Since $x^{(1)},\dots,x^{(m)}$ are arbitrary, we obtain $f/J = 0$, a contradiction. Hence $\operatorname{spec}(s) \setminus \{0\}$ has at most $2n$ elements.
\end{proof}

\begin{lemma}
The set $S_0$ is algebraic over $F$.
\end{lemma}

\begin{proof}
Let $s \in S_0$. By the previous lemma $\operatorname{spec}(s)$ is finite. Since $F$ is infinite, $F \setminus \operatorname{spec}(s)$ is infinite and for every $\alpha \notin \operatorname{spec}(s)$ we have $s-\alpha 1$ is invertible.
Consider the set $\{(s-\alpha 1)^{-1} : \alpha \in F \setminus \operatorname{spec}(s)\}$.
This set is linearly dependent since $\#F > \dim J$. Thus there exist $\alpha_1,\dots,\alpha_m \in F \setminus \operatorname{spec}(s)$ and $\beta_i \in F$, $\beta_i \neq 0$ for $i=1,\dots,m$ such that $\sum_{i=1}^m \beta_i (s-\alpha_i 1)^{-1} = 0$.
Multiplying by $(s-\alpha_1 1)\cdots(s-\alpha_m 1)$ we have
\begin{align*}
0 &= (s-\alpha_1 1)\cdots(s-\alpha_m 1) \sum_{i=1}^m \beta_i (s-\alpha_i 1)^{-1} \\
&= \sum_{i=1}^m \beta_i f_i(s), \quad \text{where } f_i(s) = \frac{(s-\alpha_1 1)\cdots(s-\alpha_m 1)}{s-\alpha_i 1}, \quad i=1,\dots,m.
\end{align*}
We have that $g(x) = \sum_{i=1}^m \beta_i f_i(x) \neq 0$ since $x-\alpha_i$ divides all summands except one. On the other hand $g(s)=0$, so $s$ is algebraic over $F$.
\end{proof}

\begin{lemma}
The algebra $S_0$ is simple.
\end{lemma}

\begin{proof}
Let $I$ be a non-zero ideal of $S_0$. Since $\operatorname{Nil}(J)=0$ we have $\operatorname{Nil}(S_0)=0$. Then $I$ is not nil, is algebraic and thus contains an idempotent $e \neq 0$ (since any non-zero element in an algebraic algebra generates a finite-dimensional associative algebra).
Let $I_1 = I$, $I_{k+1} = (I_k)^3$ for all $k \geq 1$. Then for any $k \geq 1$ we have $I_k$ an ideal of $S_0$. Set $\widetilde{I} = \bigcup_k I_k \subseteq S_0$.
We prove that $\widetilde{I}$ is an ideal of $J$.
It is enough to show that $I_{k+1} J \subseteq I_k$.
By induction, $k=1$: if $a,b,c \in I_k$ and $x \in J$ then
\[
((ab)c)x = (ab)(cx) + (ab,c,x),
\]
with $(ab)(cx) \in (I_k)^2 I_{k-1} \subseteq (I_k)^3 = I_{k+1}$ and
\[
(ab,c,x) = -(ax,c,b) - (bx,c,a) \in (I_{k-1}, I_k, I_k) \subseteq (I_{k-1})^3 = I_k.
\]
Then $\widetilde{I} \square J$ is an ideal containing a non-zero idempotent. Thus $S_0 \subseteq \widetilde{I}$. Then
\[
S_0 \subseteq \widetilde{I} \subseteq I \square S_0 \Rightarrow I = S_0.
\]
\end{proof}

\begin{lemma}
Let $F=\bar{F}$. Then $J$ is an Albert algebra.
\end{lemma}

\begin{proof}
We have $J$ an algebraic algebra. Let $e_1,\dots,e_m$ be pairwise orthogonal idempotents in $J$. We will show that $m \leq 2n$ where $n=\dim J$.

Let $s=\alpha_1 e_1+\dots+\alpha_m e_m \in J$ with $\alpha_i \neq 0$ and $\alpha_i \neq \alpha_j$ if $i \neq j$. Observe that $\alpha_i \in \operatorname{Spec}(s)$ (in fact $s-\alpha_i 1$ is not invertible, since $U_{s-\alpha_i 1}$ is not invertible: $e_i U_{s-\alpha_i 1}=(s-\alpha_i 1) \circ e_i (1-\alpha_i 1)=0$). Then $m \leq 2n$.

Since $F$ is algebraically closed, any idempotent can be decomposed into a sum of absolutely primitive idempotents, pairwise orthogonal. Let $m$ be maximal with $e=e_1+\dots+e_m$, $e_i$ pairwise orthogonal, absolutely primitive. Then
\[
J = J_1(e) \oplus J_{1/2}(e) \oplus J_0(e) \text{ with }
\]
$J_0(e)$ nil (otherwise there exists another idempotent). Thus $J_0(e)=0$. Analogously we prove that $J_{1/2}(e)=0$. Then $J=J_1(e) \Rightarrow e=1$.

So $1=e_1+\dots+e_m \in J$. Since $J$ is simple, $e_1,\dots,e_m$ are connected $\Rightarrow \forall i,j$ $\exists a_{ij} \in J_{ij}$ such that $a_{ij}^2=e_i+e_j$. By the coordinatization theorem $J \cong H_3(D,*)$ with $(D,*)$ alternative with involution.

If $D$ is associative, then $J$ is special, which does not happen since $J$ is not special. Hence $n=3$, and $(D,*)$ is simple alternative non-associative $\Rightarrow$ $J$ is an Albert algebra.
\end{proof}

\begin{theorem}
Let $J$ be a primitive Jordan algebra with unit over $F=\bar{F}$, and $\#F > \dim J$. If $J$ is not special, then $J$ is an Albert algebra.
\end{theorem}

\begin{proof}
We will show that $J_{\text{spe}} = J$. We know $1 \in J_{\text{spe}}$. We prove next that $e \in Z(J)$ and equivalently $e \in N(J)$, since $J$ is commutative.

Let $a, b \in J$. We have
\begin{align*}
(e, a, b) &= (e^2, a, b) = 2(e, a, b e) = 2(e^2, a, b e) \\
&= 4(e, a(b e), e) = 4(e, a, b e).
\end{align*}
From $2(e, a, b e)=4(e, a, b e)$ it follows that $(e, a, b e)=0 \Rightarrow (e, a, b)=0$. By the flexibility law $(b, a, e)=0$ as well and since $(a, b, c)+(b, c, a)+(c, a, b)=0$, we get $(a, e, b)=0 \Rightarrow e \in Z(J)$.

For any $a \in J$ we have $a = a e + (a - a e)$; $a e \in J e$ (since $e \in Z(J)$) and $a - a e \in I$, with $I:=\{a - a e : a \in J\}$. We show that $I$ is an ideal of $J$: for $a, b \in J$ we have
\[
(a - a e) b = a b - (a e) b \stackrel{e \in Z(J)}{=} a b - a(e b) = a b - a b e \in I.
\]
Moreover $J e \cap I = 0$ because
\begin{align*}
x = a e = b - e b &\Rightarrow x e = a e = e b - e(e b) = 0 \\
&\Rightarrow a e = 0 \Rightarrow x = 0.
\end{align*}
Then $J = J e \oplus I$. If $I \neq 0$, $I$ would be a prime ideal. Thus $I=0$ and $J = J e = J_{\text{spe}}$.
\end{proof}

\begin{corollary}
If $J$ is a Jordan algebra with unity, simple central over $F$ and $J$ is not special, then $J$ is an Albert type algebra.
\end{corollary}

\begin{definition}
We say that a Jordan algebra $J$ over $F$ is of \emph{Albert type} if there exists $K \supset F$ such that $J_K = K \otimes_F J$ is an Albert algebra over $K$.
\end{definition}

\begin{proof}[Proof of corollary]
If $\#F \leq \dim J$, we consider an extension of $F$ whose cardinality is greater than $\dim J$. Without loss of generality we may assume $K=\bar{K}$. Consider $J_K = K \otimes_F J$. Then $J_K$ is simple and central over $K$ and $J_K$ is not special. By the last theorem $J_K$ is an Albert algebra.
\end{proof}

\subsection*{Simple non-special algebras}
Suppose $J$ is not special. Let $\rho$ be the $T$-ideal generated by $\{t(n_1, n_2, n_3)^+ : n_i \in N_t(J\langle x \rangle)\}$. By previous results we know that $\rho \subseteq N_t$. In particular, $\forall a, b, c, d \in J\langle x \rangle$, we have $\{a b c\} \in \rho$, $\{a b c d\} \cdot \rho \subseteq \rho J\langle x \rangle$ (since modulo $\rho J\langle x \rangle$ we have $\{a b c d\} \cdot \rho = \{(a \cdot \rho) b c d\} + \{\rho b c d\} \cdot a \equiv 0$).

Recall
\[
\{a, b, c\} = (a b) c + (c b) a - b(a c) = b U_{a, c}.
\]
We denote $\{a, b, c\} = a V_{b, c}$. Then we may rewrite
\begin{gather*}
V_{b, c} = [R_b, R_c] + R_{b c} = D_{b, c} + R_{b c}, \\
D_{b, c} := [R_b, R_c].
\end{gather*}

Let $I \triangleleft J$. Define
\[
\operatorname{Ann}(I) = \{b \in J : \{b a x\} = 0 \text{ for all } a \in I \} = \{b \in J : b I = \{b I J\} = 0\}.
\]

\begin{lemma}
If $I \triangleleft J$ then $\operatorname{Ann}(I) \triangleleft J$.
\end{lemma}

\begin{proof}
Observe that
\[
b \in \operatorname{Ann}(I) \Leftrightarrow bI = (bJ)I = 0.
\]
In fact, it is enough to show that if $bI = 0$, then
\[
\{bIj^*\}=0 \Leftrightarrow (bJ)I = 0.
\]
Then for $a \in I$ and $x \in J^{\#}$ we have
\[
\{b,a,x\} = \underbrace{(ba)x}_{=0} + \underbrace{(xa)b}_{=0} - (bx)a,
\]
so $\{b,a,x\}=0 \Leftrightarrow (bx)a = 0$.

Thus to prove that $I$ is an ideal of $J$, it is enough to show that $(bx)I = ((bx)y)I = 0$ for $b \in \operatorname{Ann}(I)$ and $x \in J$. The first condition was proved already. For the second, $\forall y \in J$, $a \in I$:
\begin{align*}
((bx)y)a &= (b(xy))a + (b,x,y)a \\
&= (b,a x,y) - (b,a,y)x \in (b,I,y) + (b,I,y)J = 0,
\end{align*}
since $(b,I,y)=0$.
\end{proof}

For $A, B \subset \mathcal{J}$ denote $V(A,B) = \{ V_{a,b} \mid a \in A, b \in B \}$.

\begin{lemma}
For every $a \in H = H(F\langle X \rangle,*)$ there exists $m \in \mathbb{N}$ such that
\[
a V(P, \operatorname{Sym}\langle X \rangle^{\#})^{m} \subseteq \operatorname{Sym}\langle X \rangle.
\]
\end{lemma}

\begin{proof}
We first prove that for all $x,y,z,t,x_1,\dots,x_s, a \in X$ and $p \in P$, we have
\begin{enumerate}
\item[(i)] $\{xyzt\} R_{x_1} \dots R_{x_s} V_{p,a} \in \operatorname{Sym}\langle X \rangle$;
\item[(ii)] $\{xyzt\} R_{x_1} \dots R_{x_s} R_p \in \operatorname{Sym}\langle X \rangle$.
\end{enumerate}
To verify them we use induction on $s$.

For $s=0$, (ii) is proved already and for (i) note that modulo $\operatorname{Sym}\langle X \rangle$,
\[
\{xyzt\} V_{p,a} = \{xyzt\}(D_{p,a} + R_{pa}) \equiv \{xyzt\} D_{p,a} \quad \text{since } pa \in P.
\]
Now
\begin{align*}
\{xyzt\} D_{p,a} &= (p,\{xyzt\},a)_+ = \frac{1}{4}[\{xyzt\},[p,a]] \\
&= \{x D_{p,a} y z t\} + \{x y D_{p,a} z t\} + \{x y z D_{p,a} t\} + \{x y z t D_{p,a}\} \equiv 0,
\end{align*}
since $x D_{p,a}, t D_{p,a} \in D_{p,a}$. This proves the case $s=0$.

Suppose (i) and (ii) hold for some $s \geq 0$. Then
\begin{align*}
&\{xyzt\} R_{x_1} \dots R_{x_{s+1}} R_p \\
&= \{xyzt\} R_{x_1} \dots R_{x_s} R_p R_{x_{s+1}} - \{xyzt\} R_{x_1} \dots R_{x_s} [R_p, R_{x_{s+1}}] \\
&\equiv -\{xyzt\} R_{x_1} \dots R_{x_s} (V_{p,x_{s+1}} + R_{p x_{s+1}}) \equiv 0,
\end{align*}
by the induction hypothesis. Thus (ii) is valid for $s+1$.

Next,
\begin{align*}
\{xyzt\} R_{x_1} \dots R_{x_{s+1}} V_{p,a} &= \{xyzt\} R_{x_1} \dots R_{x_{s+1}} (D_{p,a} + R_{pa}).
\end{align*}
Observe that
\begin{align*}
&\{xyzt\} R_{x_1} \dots R_{x_{s+1}} D_{p,a} \\
&\equiv \{ \{xyzt\} R_{x_1} \dots R_{x_s} D_{p,a} \} R_{x_{s+1}} + \{xyzt\} R_{x_1} \dots R_{x_s} [R_{x_{s+1}}, D_{p,a}] \\
&\equiv 0 \quad \text{(by induction)} \\
&\equiv \{xyzt\} R_{x_1} \dots R_{x_s} R_{x_{s+1} D_{p,a}}, \quad \text{since } [R_a, D] = R_{aD}
\end{align*}
for any derivation $D$.

Then the first summand belongs to $\operatorname{Sym}\langle X \rangle$ by the induction hypothesis. The second summand is in $\operatorname{Sym}\langle X \rangle$ as well, by (ii). Therefore the right-hand side is in $\operatorname{Sym}\langle X \rangle$, proving (i).

Let $a \in H$. Then by Cohn's theorem we have $a = j(t_1,\dots,t_r, x_1,\dots,x_s)$ with $t_i$ tetrads and $x_j \in X$. Let $d(a)$ denote the degree of $a$ in tetrads, that is, how many tetrads appear in $a$. We prove the main result of the lemma by induction on $d(a)$.

If $d(a)=0$, then $a$ is already a Jordan element. Suppose $d(a) \geq 1$. Then $a$ is a sum of terms $t R_{a_1} \dots R_{a_k}$, with $t$ a tetrad and $a_i$ monomials in $t_j, x_j$ for all $i=1,\dots,k$. Without loss of generality we may assume that $a$ is a single monomial, that is, $a = t R_{a_1} \dots R_{a_k}$. We have $\sum_{i=1}^k d(a_i) = d(a)-1$. Then for any $p \in P$ and $x \in J$, we have
\[
a V_{p,x} = t R_{a_1} \dots R_{a_k} V_{p,x} = \sum W_i,
\]
with $W_i = j_i(a_1,\dots,a_k,x)$ (using that $\{xyzt\} R_{x_1} \dots R_{x_k} V_{p,a} = j(x,y,z,t,x_1,\dots,x_k,a)$ and an evaluation homomorphism $x_i \to a_i$). We have $d(W_i) < d(a)$ for each $i$.

By the induction hypothesis, there exists $m$ such that
\[
W_i V(P, \operatorname{Sym}\langle X \rangle^{\#})^{m} \subseteq \operatorname{Sym}\langle X \rangle.
\]
Then $a V(P, \operatorname{Sym}\langle X \rangle^{\#})^{m+1} \subseteq \operatorname{Sym}\langle X \rangle$.
\end{proof}

Let $J = \operatorname{Sym}\langle X \rangle / I$ with $P(J)=0$ and $I$ is \dots


\begin{proof}
Let $\dot{I}=\operatorname{id}_{F\langle x\rangle}\langle I\rangle$ and let $f \in \delta y\langle x\rangle \cap \hat{I}$.
By the Cohn criterion we have to show that $f \in I$. Write
\[
f=\sum_{k=1}^{l}\left\{i_{n}a_{i}^{(k)}\ldots a_{n_{k}}^{(k)}\right\} \quad \text{with } i_{k}^{(k)} \in X.
\]
Consider $V=\{y_{1},\ldots,y_{s}\}$, the algebra $F\langle Y \cup X\rangle$ and
\[
\tilde{f}=\sum_{i=1}^{l}\left\{y_{k}a_{i}^{(k)}\ldots a_{i_{k}}^{(k)}\right\} \in H\langle V \cup X\rangle.
\]
By the previous lemma there exists $m \geq 1$ such that
\[
\tilde{f}V\left(p, S y\langle Y \cup X\rangle^{\#}\right)^{m} \subseteq S y\langle Y \cup X\rangle.
\]
Let $\varphi: S y\langle X \cup Y\rangle \rightarrow S y\langle X\rangle$ be the homomorphism such that $\varphi|_{X}=\operatorname{id}_{X}$ and $\varphi(y_{k})=i_{k}$ for all $k=1,\ldots,s$. Applying $\varphi$ we obtain
\[
\varphi(\tilde{f})=f \quad \text{and} \quad f V\left(\left(P, S y\langle Y \cup X\rangle^{\#}\right)^{m}\right) \subseteq \operatorname{id}_{S y\langle X\rangle}\langle i_{1},\ldots,i_{s}\rangle \subseteq I.
\]
Then in $y$ we have
\[
j V\left(p(y), j^{\#}\right)^{m-1} V\left(p(y), j^{\#}\right)=\bar{f} V\left(p(y) j^{m}\right)=0.
\]
Thus $\bar{f} V\left(p(y), y^{*}\right)^{m-1} \in \operatorname{ann} p(z)$. Since $z$ is prime and $p(y) \neq 0$, we have $\bar{f} V\left(p(y), y^{*}\right)^{m-1}=0$. We repeat the process until we get $\bar{f} \equiv 0$ in $J$. Then $f \in I$.
\end{proof}

\begin{corollary}
Let $J \in \overline{S 70}$ be a simple Jordan algebra. Then we have the following possibilities:
\begin{enumerate}
\item $J \simeq H(A,*)$, where $(A,*)$ is $*$-simple associative.
\item $p(y)=0$ and $[x,y]^{2} \in y$.
\begin{enumerate}
\item If $[x,y]^{2}/y \not\equiv 0$, then $J \simeq J(V,f)$.
\item If $[x,y]^{2}/y \equiv 0$ and $J$ is nondegenerate, then $J$ is associative.
\end{enumerate}
\end{enumerate}
\end{corollary}

\begin{theorem}
Let $J$ be a special Jordan algebra. Then $J \simeq H(A,*)$ for some $(A,*)$, or $J$ is a PI Jordan algebra.
\end{theorem}

\begin{theorem}
Let $J \in \widehat{\operatorname{SJord}}$ be a prime algebra satisfying $f=0$. Suppose $[x,y]^{2}=0$ in $J$. Then $\partial^{-1}J \simeq J(V,f)$.
\end{theorem}

\begin{remark}
The difference between prime, nondegenerate and primitive algebras: the class of prime and nondegenerate Jordan algebras is much larger than the class of primitive algebras. Primitivity imposes additional conditions related to the existence of specific types of inner ideals. Over a field of characteristic $\neq 2$, a primitive Jordan algebra satisfying a PI must be simple and unital (e.g., finite dimensional over its center, or Jordan algebras of a nondegenerate quadratic form).
\end{remark}

\begin{remark}
Inner/quadratic ideals serve as the counterpart to one-sided ideals in associative algebras. Unlike standard ideals, which are closed under outer multiplication by any element of the algebra, inner ideals are closed under the sandwich operation: for $b \in J$,
\[
U_{b}(J)=\{U_{b}(x) \mid x \in J\}.
\]
We introduce the McCrimmon radical and show that it is locally nilpotent. A simple degenerate algebra is locally nilpotent.
\end{remark}

\end{document}
